\documentclass[11pt,a4paper,oneside,openany]{book}

\usepackage[xetex,top=24mm,bottom=25mm,left=27mm,right=25mm,headheight=15pt]{geometry}
\usepackage{fontspec,xeCJK}
\setmainfont{Times New Roman}
\setCJKmainfont[BoldFont=SimHei,ItalicFont=KaiTi]{SimSun}
\setCJKsansfont{Microsoft YaHei}
\setCJKmonofont{SimSun}
\usepackage{amsmath,amssymb,amsthm,mathtools}
\usepackage{booktabs,longtable,array,graphicx}
\usepackage{fancyhdr}
\usepackage{xcolor}
\usepackage[unicode,bookmarks=false,colorlinks=true,linkcolor=blue!45!black,citecolor=blue!45!black,urlcolor=blue!45!black]{hyperref}
\newcommand{\Z}{\mathbb Z}
\newcommand{\Q}{\mathbb Q}
\newcommand{\F}{\mathbb F}
\newcommand{\R}{\mathbb R}
\newcommand{\SL}{\operatorname{SL}}
\newcommand{\GL}{\operatorname{GL}}
\newcommand{\diag}{\operatorname{diag}}
\newcommand{\rad}{\operatorname{rad}}
\newcommand{\cont}{\operatorname{cont}}
\newcommand{\Um}{\operatorname{Um}}
\newcommand{\ord}{\operatorname{ord}}
\newtheorem{theorem}{定理}[chapter]
\newtheorem{lemma}[theorem]{引理}
\newtheorem{proposition}[theorem]{命题}
\newtheorem{corollary}[theorem]{推论}
\theoremstyle{definition}
\newtheorem{definition}[theorem]{定义}
\newtheorem{example}[theorem]{例}
\newtheorem{remark}[theorem]{说明}
\renewcommand{\proofname}{证明}
\renewcommand{\contentsname}{目录}
\renewcommand{\bibname}{参考文献}
\renewcommand{\tablename}{表}
\renewcommand{\figurename}{图}
\renewcommand{\chaptername}{讲}
\renewcommand{\appendixname}{附录}
\numberwithin{equation}{chapter}
\setcounter{tocdepth}{1}
\setcounter{secnumdepth}{2}
\setlength{\parindent}{2em}
\setlength{\parskip}{4pt plus 1pt}
\linespread{1.16}
\emergencystretch=3em
\allowdisplaybreaks[2]
\raggedbottom
\pagestyle{fancy}
\fancyhf{}
\fancyhead[L]{\small 素数平均：从例子到完整证明}
\fancyhead[R]{\small\thepage}
\renewcommand{\headrulewidth}{0.3pt}
\fancypagestyle{plain}{\fancyhf{}\fancyfoot[C]{\thepage}\renewcommand{\headrulewidth}{0pt}}
\makeatletter
\def\@makechapterhead#1{\vspace*{8mm}{\parindent\z@\raggedright\normalfont
 \ifnum\c@secnumdepth>\m@ne\if@mainmatter
 {\Large\sffamily\color{blue!40!black}\@chapapp\ \thechapter\par}\vskip12pt\fi\fi
 \interlinepenalty\@M\LARGE\bfseries #1\par\nobreak\vskip24pt}}
\def\@makeschapterhead#1{\vspace*{8mm}{\parindent\z@\raggedright\normalfont
 \interlinepenalty\@M\LARGE\bfseries #1\par\nobreak\vskip24pt}}
\makeatother
\hypersetup{pdftitle={素数平均：从例子到完整证明},pdfsubject={中文数学讲义}}

\begin{document}
\frontmatter
\begin{titlepage}
\centering
\vspace*{25mm}
{\sffamily\Large 数学讲义\par}
\vspace{15mm}
{\Huge\bfseries 素数平均\par}
\vspace{7mm}
{\LARGE 从例子到完整证明\par}
\vspace{16mm}
{\large 面向数学本科生的逐步讲解\par}
\vfill
\begin{minipage}{.79\textwidth}
\large
对奇素数 $p$，当 $n\ge 2p+1$ 时，\par
有限次 $p$ 元平均可达全等\par
\hfill 当且仅当\quad $G$ 是 $p$ 的幂。
\end{minipage}
\vfill
{\large 通俗展开版\quad 2026年9月\par}
\end{titlepage}
\chapter*{先把问题讲明白，再进入证明}
\addcontentsline{toc}{chapter}{先把问题讲明白，再进入证明}
这份讲义从一个很具体的问题出发：一排数，每次选几个，把它们都换成平均值，能不能经过有限次操作使它们完全相等？你只需要先知道平均数、最大公因数、模运算和二阶矩阵乘法。后面会用到群和环，但我们会先说明这些对象解决什么问题，再给定义。

读数学证明时，最容易迷路的地方不是某个计算太长，而是不知道为什么突然要算它。例如，为什么研究平均数的人开始写 $2\times2$ 矩阵？为什么需要一个局部化环？为什么把两个数变成模 $q$ 相同，竟然就离精确相同不远了？本版把这些转折作为讲解的重点。

\section*{先记住两个动作}
第一个动作是\textbf{下降}。如果每一步都让一个非负整数变小，就不能一直做下去。平方和正是我们的第一个下降量。但为了使平方和一直为整数，选择平均组时要满足整除条件，还要留下某些同余差异。

第二个动作是\textbf{输送}。位置较少时，强求每一步都是整数平均会卡住。我们改为把整排数整理成三个等值块，只用两个参数描述它们。几次平均合在一起，就成为两个参数上的一个矩阵。取得足够的矩阵后，直接把参数送到一个能收尾的地方。

这两种方法都要先回答同一个问题：哪些输入从一开始就绝无可能成功？答案来自一个最大公因数 $G$。前两讲会算清楚它，不要求你先接受任何算术群结论。

\section*{第一次和第二次读，任务可以不同}
第一次沿着每讲开头的例子和问题走，目标是能用自己的话解释每个工具为什么出现。第5讲的五平均十五元计数表、第8讲的具体输送矩阵、第10讲的五平均十一元分组，尤其值得亲手算一遍。

第二次再逐行核对各个正式命题的证明。所有原讲义的正式命题、证明段和有限计数表都保留在相应讲次中，讲解插在它们前后。新讲解会把长公式拆成可跟踪的任务，但不会把有限分类替换成随机实验，也不会省略原来的参数分支。

第10讲以后是最费计数的部分。此时先问三件事：这张表从各块取了几个位置？一共取够了没有、取多了没有？取完以后想得到哪一个矩阵？理解这三件事后，再看配比公式，公式就不再是一串没有来由的字母。

\section*{几个会反复出现的词}
\begin{longtable}{p{.20\textwidth}p{.71\textwidth}}
\toprule 讲义用语 & 先用普通语言理解\\\midrule
本原 & 将整数坐标共同的最大公因数约掉。\\
危险素数 & $n$ 的素因子中不等于平均元数 $p$ 的那些。\\
见证 & 两个余数不同的位置，它们证明“现在还没有全同余”。\\
重值 & 某个数值至少出现 $p$ 次。数值的大小不限。\\
块 & 一组当前数值相等的位置。位置并未被合并成一个。\\
返回、宏 & 一段真实平均操作，做完后恢复同一种块重数。\\
射影 & 暂时只看两个参数的比例；实际配置仍带着共同尺度。\\
单位 & 在所讨论的环里有乘法逆元的数。\\
剪切、根 & 给一个坐标加上另一个坐标的某个倍数。\\
同余层 $q$ & 要求矩阵模 $q$ 为单位矩阵；$q$ 越大通常要求越严。\\
同余纤维 & 所有模 $q$ 给出同一个列向量的本原参数。\\
内容 & 整数列两个坐标的最大公因数。\\\bottomrule
\end{longtable}
“安全”在本讲义中是数学定义：操作后仍保留所要求的同余差异。“深核”在原稿中指较深的主同余子群；正文会尽量直接写“已经能实现的 $\Gamma(q,R)$”，让你看见实际用了哪个模数。

这是一份对既有证明的教学改写。Morris 和 Serre 的结果作为外部定理引用，其作用和前提在第8讲说明。小参数中的精确有限证书完整保留。新增例子和教学叙述的检查，不等于对整项研究作了一次独立同行审稿。

\tableofcontents
\mainmatter
\chapter{先玩一个例子：为什么接近相等还不够}\label{sec:main}
\section{每一步到底允许做什么}
看这七个数：
\[(-3,0,0,0,1,1,1).
\]
每次选三个位置，把这三个数同时改成平均值。先选 $(-3,0,0)$，它们都变成 $-1$。现在有三个 $-1$、三个 $1$、一个 $0$。再选 $(-1,0,1)$，它们都变成零；这样重复三次，整排数就全为零。总共只做了四次平均。

这里的三个零是三个不同位置，不是一个数只写了一遍。以后我们用 $0^3$ 表示“三份零”，用 $a^m$ 表示“$m$ 个位置都取值 $a$”。在这种配置记号里，上标表示份数。

现在把“三个”换成一个素数 $p$，把总位置数记为 $n$。两个字母管不同的事：$p$ 是一次操作选多少个，$n$ 是总共能使用多少个位置。所有操作都在原来的 $n$ 个位置上进行。

\section{终点其实已经确定了}
一次操作把 $p$ 个数换成 $p$ 份平均值，总和没有变化。因此，假如最终全相等，公共值必然是
\[A=\frac{x_1+\cdots+x_n}{n}.
\]
我们不必猜终点，只需判断能不能到达这个确定的 $A$。

“越来越接近 $A$”和“某一步恰好全等于 $A$”之间有很大差别。数列 $1,1/2,1/4,\ldots$ 越来越小，却没有一项是零。平均确实会让差异变小，但仅凭这一点还不能回答有限步的问题。

\section{先把题目化成整数归零}
把所有坐标减去 $A$，终点就变成零。再同时乘一个数清除分母，最后除去坐标共同的最大公因数。这不改变成功的位置选择序列，因为取平均与整体平移、整体缩放交换。

例如输入 $(1,2,4)$ 的均值为 $7/3$。减去均值得 $(-4/3,-1/3,5/3)$，再乘 $3$ 得
\[X=(-4,-1,5).
\]
三个坐标的最大公因数是 $1$，不用再约分。我们把这样的整数代表叫作\textbf{本原中心化代表}：“中心化”表示总和为零，“本原”表示已经约掉共同整数因子。

一般写成
\[X=c(x-A\mathbf1)\in\Z^n,\qquad \sum_iX_i=0,\qquad\gcd_iX_i=1.
\]
对非恒定输入，这个代表只差一个整体正负号。现在定义
\begin{equation}\label{eq:G}
G(X)=\gcd_{i<j}|X_i-X_j|.
\end{equation}
这是\emph{差的}最大公因数，与坐标自己的最大公因数不同。上例坐标本原，但三个差为 $3,9,6$，所以 $G=3$。下一讲会解释，为什么恰好该看这些差。

\section{我们要证明的答案}
\begin{theorem}[素数平均的充要条件]\label{thm:prime}
设 $p$ 为奇素数，$n\ge2p+1$。非恒定配置 $x\in\Q^n$ 可经有限次 $p$ 平均达到全等，当且仅当 $G(x)=p^e$，其中 $e\ge0$。对每个奇素数 $p$，在 $n=2p$ 存在 $G=1$ 的不可达配置。
\end{theorem}
这里 $p^0=1$ 也允许。这个定理不是说位置足够多以后所有配置都能平均；它说此时障碍恰好由 $G$ 是否含非 $p$ 因子决定。
\begin{example}[如何算判据]
取五平均、十一位置的配置
\[
X=(1^{10},-10).
\]
它已经零和而且本原。任意两个不同坐标的差都是 $11$，所以 $G=11$。
虽然位置数已达到 $2p+1=11$，它仍然不能归零：模 $11$ 的每个数都等于 $1$，
五平均也保持这个余数。
相比之下，$(-5,0^5,1^5)$ 同样有十一位置，且 $G=1$；
先平均 $(-5,0^4)$，便化成相反值和一个零，从而可解。
\end{example}
\begin{corollary}[最优最终阈值]\label{cor:threshold}
记 $N(p)$ 为使上述充要条件对所有 $n\ge N(p)$ 成立的最小整数。则奇素数满足 $N(p)=2p+1$；二平均满足 $N(2)=4$。
\end{corollary}
“最终阈值”是从某个位置数开始，\emph{每一个更大的位置数}都满足判据。例如 $n=p$ 时全体平均一次就完成，但 $n=2p$ 反而存在 $G=1$ 的失败例子。因此不能把“第一个成功的维数”和“以后永远成立的界”混为一谈。

\section{后面的证明怎样走}
先证明有非 $p$ 障碍必然失败。这部分很短。真正长的是反方向：没有这些障碍，为什么一定能选出一条成功路径？

位置较多时，我们让平方和作为非负整数不断下降；这需要精心挑选平均组。位置紧张时，我们把整排数整理成三个等值块，用两个参数记录；再用几次平均组成一个矩阵动作，把参数送到容易收尾的位置。二者最后都会遇到一个共同的好消息：一旦选出和为零的 $p$ 个数，接下来已有独立的收尾方法。

你现在只需记住这三件事：目标可以化成零；$G$ 负责判断最初的算术障碍；充分性要同时解决“下一步怎么选”和“为什么有限步结束”。接下来逐件说明。

\chapter{为什么一个最大公因数能阻止所有平均}\label{sec:obstruction}
\section{七个数模七都等于一，会发生什么}
考虑三平均的配置 $(1^6,-6)$。总和是零，但每个数模 $7$ 都等于 $1$。任取三个，其和模 $7$ 等于 $3$；除以 $3$ 后，余数仍是 $1$。所以所有位置在每一步都模 $7$ 等于 $1$，不可能全变成零。

可能你会担心：平均后出现分数，还能谈模 $7$ 吗？可以，因为这些分母只是 $3$ 的幂，模 $7$ 都有逆元。例如 $1/3\equiv5\pmod7$，因为 $3\cdot5\equiv1$。所以“除以三”在模七下仍然说得清。

这个解释也说明为什么模平均元数 $p$ 不一样：$p$ 没有模 $p$ 的逆元。因此同样的阻碍不能拿来禁止 $G$ 中的 $p$ 因子。

\section{把例子写成一般证明}
先确认上一讲的归一化。
\begin{lemma}[归一化与操作序列相容]\label{lem:affine}
任一平均词 $W$ 满足 $W(a\mathbf1+bx)=a\mathbf1+bWx$。因此配置达到均值，等价于其本原中心化代表达到零。
\end{lemma}
\begin{proof}
每个 $T_S$ 线性且固定 $\mathbf1$，对词长归纳即可。
\end{proof}
如果 $\ell\mid G$，意思正是所有坐标模 $\ell$ 相等。为什么共同余数不能恰好是零？因为那会使所有整数坐标都有公因子 $\ell$，违背本原性。这一句解释了归一化为何必须约掉共同因子。
\begin{lemma}[非平均素数处的共同余数障碍]\label{lem:obstruction}
非零本原整数零和配置满足 $G\mid n$。若素数 $\ell\ne p$ 整除 $G$，则配置不可归零。
\end{lemma}
\begin{proof}
模 $G$ 的全部坐标同余于某个 $a$。本原性给出 $\gcd(a,G)=1$，零和给出 $na\equiv0\pmod G$，所以 $G\mid n$。若 $\ell\mid G$，共同余数模 $\ell$ 非零。所有中间值属于 $\Z[1/p]$，该环可约化至 $\F_\ell$；每次平均保持这个非零共同余数，故不能归零。
\end{proof}
其中 $G\mid n$ 很有用：将来要防范的素数不是所有素数，只是 $n$ 的非 $p$ 素因子。

称当前整数零和配置\textbf{合法}，是指它在每个 $\ell\mid n,\ell\ne p$ 处都不全同余。称保持这项性质的整数平均为\textbf{安全操作}。记这些需要防范的不同素数总数为
\[d_p(n)=\#\{\ell:\ell\mid n,\ell\ne p\text{ 为素数}\}.
\]
初始本原配置合法，当且仅当 $G$ 是 $p$ 的幂。操作后的配置未必还本原，所以后面会区分“当前实际整数值的余数”与“约去公因子后的判据”。

\section{合法输入也会被错误的一步带坏}
三平均七元输入 $(3^5,1,-16)$ 的总和为零，$G=\gcd(2,19)=1$。它是合法的。但如果先平均 $(3,1,-16)$，均值为 $-4$，输出就成了
\[(3^4,(-4)^3).
\]
现在所有数模 $7$ 都等于 $3$，再也不能归零。这说明 $G$ 并不是每一步自动保持的量。充分性不能任意平均，必须安排操作以保留必要的余数差异。

\section{为什么恰好少一个位置就会失败}
要证明 $2p+1$ 这个界不能再低，只需在 $2p$ 个位置找一个 $G=1$ 的失败例子。候选是
\[(-p,0^{p-1},1^p).
\]
它含 $0,1$，所以 $G=1$。如果只允许整数平均，唯一非恒等选择是 $(-p,0^{p-1})$，结果为 $((-1)^p,1^p)$。由于 $p$ 奇，这个配置所有数模 $2$ 都等于 $1$，失败。

但这里有个必须补上的问题：为什么不能先做一次产生分数的操作，绕过去？下面用“分母里有几个 $p$”排除它。

对非零有理数写 $x=p^h u/v$，$p\nmid uv$，定义 $v_p(x)=h$。例如 $v_3(18)=2$，$v_3(2/9)=-2$，并规定 $v_p(0)=+\infty$。只需要知道：两项赋值不同时，和的赋值等于较小的那个，因为提取最小 $p$ 幂后，一个单位加一个被 $p$ 整除的数不会变成零模 $p$。

首次产生分数时，恰好有新选中的 $p$ 个位置具有最低负赋值。下一次若只选它们中的 $k$ 个，$1\le k<p$，因 $p\nmid k$，最低项无法被其他较高赋值的项抵消；除以 $p$ 后，负赋值反而再降低一。这就是一直跟踪的东西。
\begin{lemma}[小维数中首次分数化不可挽回]\label{lem:fraction}
设 $p$ 为奇素数，$p\le n\le2p$。整数零和配置的一条成功平均路径不可能首次产生非整数。
\end{lemma}
\begin{proof}
首次非整数平均后，最低 $p$ 进赋值 $h<0$ 恰由一个等值的 $p$ 块承担。若下一选择与此块交于 $k$ 个位置，$1\le k<p$，则其和的赋值为 $h$，新均值的赋值为 $h-1$，并形成新的最低 $p$ 块。选整个最低块时操作恒等。若选择与之不交，则只能 $n=2p$ 且选中补集；零和迫使结果成为 $(a^p,(-a)^p)$。按引理\ref{lem:affine}，它等价于 $(1^p,(-1)^p)$，被模 $2$ 的共同非零余数阻碍。故最低负赋值既不能消失，也不能通过不交操作进入可解配置。
\end{proof}
因此整数分支和分数分支都已排除，反例才真正完整。
\begin{proposition}[位于2p的反例]\label{prop:lower}
对奇素数 $p$，配置 $(-p,0^{p-1},1^p)$ 本原、零和、$G=1$，但不可归零。
\end{proposition}
\begin{proof}
由引理\ref{lem:fraction}，成功路径必须始终采用整数平均。若一次选择包含 $j$ 份 $1$ 和 $e$ 份 $-p$，整性要求 $j-ep\equiv0\pmod p$，故 $j=0$ 或 $p$。唯一非恒等可能是选择 $(-p,0^{p-1})$，得到 $((-1)^p,1^p)$；引理\ref{lem:obstruction}排除该配置。
\end{proof}

\section{二平均为什么单列}
二平均的最终阈值是 $4$，而不是把 $p=2$ 代入奇素数公式所得的 $5$。文献\cite{mixing}的条件MC说：若所有整数输入在一个奇模数下相等，则目标均值也必须在那个余数类里。对本原零和输入，这正好等价于 $G$ 没有奇素因子。

三元反例为 $(-3,1,2)$。任意首次平均后都成 $(a,a,-2a)$、$a\ne0$；以后非恒等平均仍是同形状，只把 $a$ 变为 $-a/2$。它越变越小但不会是零。四元以上的充分性引用上述论文；其两个输入、两个输出的混合器按时间顺序执行，就是本题两个原位置同时取平均。

\chapter{先学会安全地走一步}\label{sec:descent}
\section{平方和下降，到底还缺什么}
令 $E(x)=\sum_i x_i^2$。平均一组 $S$，其均值为 $a$，则
\[
E(x)-E(T_Sx)=\sum_{i\in S}x_i^2-pa^2
=\sum_{i\in S}(x_i-a)^2.
\]
右边非负，且除非选中的值全相等，否则严格为正。这是“平均让数更接近”的精确版本。

要从这里得到有限终止，还缺两个条件。第一，中间数最好都是整数，这样下降量不可能无穷趋近于零。第二，别走到上一讲的共同余数陷阱里。我们将把这两个条件和一种位置储备一起保持：\textbf{有两个不同数值，各出现至少 $p$ 次。}

这种数值叫“重值”。比如五平均中 $a$ 出现六次，它是重值；出现四次则不是。重值不是说 $a$ 的绝对值大。

\section{留在外面的两个数，为什么能保护这一步}
假如未选位置里还留着一个奇数和一个偶数，那么不管选中的位置怎样变化，输出都不可能全同奇偶。模任意素数也是同一个道理。因此，与其直接检查新的均值，不如事先在外面保留一对不同余位置。

什么时候这会紧张？只有绝大多数位置都属于同一个余数类、少数例外很容易被一次全拿走的时候。
\begin{definition}[危险素数与危险集]
称 $\ell\mid n$ 且 $\ell\ne p$ 为危险素数。
若模 $\ell$ 有一个余数类包含至少 $n-p$ 个位置，称它为主类，
将其余位置构成的集合记为 $F_\ell$。
如果一次平均会使全部位置同余，那么未选的 $n-p$ 个位置已经必须同余。
因此只有这种主类存在时，才有一步制造障碍的风险。
\end{definition}
例如五平均二十元中，若十八个位置是偶数、两处是奇数，那么不要把两处奇数同时拿走；保留一处，就保留了奇偶差异。对每个危险素数留一个例外位置，总共最多留 $d_p(n)$ 个。
\begin{lemma}[由未选位置保证安全]\label{lem:safety}
设 $n>2p$。若模危险素数 $\ell$ 存在大小至少 $n-p$ 的余类，其补集 $F_\ell$ 唯一，且 $2\le|F_\ell|\le p$。不包含任何完整 $F_\ell$ 的整数平均安全。以下整数选择也安全：
\begin{enumerate}
\item $(a^{p-1},x)$，其中至少一份 $a$ 未选；
\item 由若干份 $a,b$ 及一份第三值 $x$ 组成，且 $a,b$ 均有未选副本；
\item 每种被选数值均有未选副本。
\end{enumerate}
\end{lemma}
\begin{proof}
$n-p>n/2$ 给出主类唯一性。例外数不能为 $1$，否则零和模 $\ell$ 迫使唯一例外仍属主类。避开完整 $F_\ell$ 时，未选位置保留主类及例外。对后三种选择，若输出全同余，未选副本及选中和为 $p$ 倍新均值依次迫使全部被选值也属于共同余类，矛盾。
\end{proof}
这里还给了几个更省位置的选择。最常用的是 $(a^{p-1},x)$：留下的一份 $a$ 会迫使任何假想的共同余数等于 $a$，而平均方程再迫使 $x$ 也等于这个余数，矛盾。正是这一份重复值，省掉了另找见证的麻烦。

\section{为什么要保留两个重值}
一次平均会产生 $p$ 份相等的新值。如果另外还保留着一个旧的 $p$ 重值，而且新旧两值不同，就又回到“两个重值”的状态，可以继续下一轮。

注意，新均值若恰好等于保留的旧值，就只剩一个重值。这种情况叫“碰撞”。它不是说原配置不可解，只是当前的储备没有保持住。下一讲会通过碰撞等式重新选一组。

初始化时，在分析中先考虑 $pX$。所有初值都被 $p$ 整除，所以第一次任取 $p$ 项的均值是整数。选出位置序列后，在原 $X$ 上执行同一序列即可。后面的整数尺度固定下来，不在下降途中随意改变。

一个小技巧反复有用：对 $p+1$ 个候选数，删掉一个后剩 $p$ 个。如果所有删法都得到同一个和，比较 $S-x_i$ 与 $S-x_j$ 就知 $x_i=x_j$。因此，多留一个候选位置便有机会避开指定均值。
\begin{lemma}[建立两个不同重值]\label{lem:initialize}
若 $n\ge3p+d$，从 $pX$ 出发至多两次整数平均可建立两个不同重值。
\end{lemma}
\begin{proof}
在任一候选池中，为仍完整包含在池内的每个 $F_\ell$ 留出一个位置，至多留出 $d$ 个。若池内所有值同余模 $p$、每种实际值不足 $p$ 份，余下 $p$ 个位置给安全非恒等整数平均；余下 $p+1$ 个位置时还可避开一个指定均值，因为所有删一子集同和会迫使这 $p+1$ 个值全相等。
已有两个重值时无须操作。仅一个重值 $a$ 且重数至少 $2p-1$ 时，选择 $(a^{p-1},b)$，留下 $p$ 份 $a$ 并生成另一重值。若重数至多 $2p-2$，在 $a$ 外的池中避开均值 $a$。若初始没有重值，先制造一个；它的重数至多 $2p-1$，其外至少 $d+p+1$ 个未改变的初值仍被 $p$ 整除，故第二次可制造不同重值。
\end{proof}

\section{整除要求交给一个零和定理}
安全性安排好以后，还要使所选和被 $p$ 整除。只看模 $p$ 的余数，这就是找 $p$ 项零和。Erdős--Ginzburg--Ziv 定理在此处的结论很直观：$2p-1$ 个候选永远够；少一个时，失败的形状也能描述。
\begin{example}[五平均的边界池]
模 $5$ 的池 $(0^4,1^4)$ 没有五项零和。
五项中必须包含 $1,2,3$ 或 $4$ 个 $1$，和都不被 $5$ 整除。
添加任意一个第九位置以后，零和引理保证出现五项零和。
逆命题说，长度为八的失败池在仿射改名后只能有这种两类各四份的形状。
\end{example}
\begin{lemma}[素数零和选择及逆边界]\label{lem:egz}
任意 $2p-1$ 个 $\F_p$ 元素包含 $p$ 项零和。长度 $2p-2$ 的序列没有 $p$ 项零和，当且仅当它为 $u^{p-1}v^{p-1}$、$u\ne v$。
\end{lemma}
\begin{proof}
先说明有限域计数：若 $m$ 元多项式 $F_i$ 的次数和小于 $m$，其共同零点数模 $p$ 等于
\[\sum_{z\in\F_p^m}\prod_i(1-F_i(z)^{p-1})=0.
\]
这里把 $1-F_i(z)^{p-1}$ 看成一个指示函数：$F_i(z)=0$ 时它为 $1$，否则为 $0$。
对一个变量，有 $\sum_{t\in\F_p}t^k=0$，除非 $k$ 是正的 $p-1$ 倍数；当 $k=0$ 时求和也是 $p=0$。
展开后每个单项式的总次数小于 $m(p-1)$，至少一个变量的指数不是正的 $p-1$ 倍数，对它求和为零。应用于
\[\sum_{i=1}^{2p-1}z_i^{p-1}=0,\qquad
\sum_{i=1}^{2p-1}a_i z_i^{p-1}=0.
\]
由于 $z_i^{p-1}$ 在 $z_i\ne0$ 时为 $1$，第一个方程要求非零坐标个数被 $p$ 整除。
一共只有 $2p-1$ 个变量，所以非零解的支持恰有 $p$ 项；第二个方程要求该支持对应的 $a_i$ 之和为零。
每个合格支持允许每一非零坐标任取 $p-1$ 个值，零点数为 $1+(p-1)^p M$，故 $M\equiv1\pmod p$。

证明逆命题。出现 $p$ 次同一余数时已有零和；仅两类时只能为所示形式。若至少三类且各重数至多 $p-1$，可选两对不交位置，使两对的差非零且不互为正负。取重数最大的类 $a$，它至少有两份；若两个其他类对 $a$ 的差不互为正负即可。否则仅有 $a,b,c$ 三类，且 $b,c$ 关于 $a$ 对称；$b,c$ 至少一类有两份，改用 $(a,b),(b,c)$，两差相差倍率 $2$。
余下 $2p-6$ 项最大重数至多 $p-2$，可再配 $p-4$ 对不同余数，留下两个单点。所用匹配事实为：多重数序列的异类最大匹配数是 $\min(\lfloor M/2\rfloor,M-f_{\max})$。若某类超过一半，将其与其余类配对；否则将各类连续排成一列，前后半列配对即达到该界。
前两对各取一项的和集有四个元素。以后每加一对，和集至少增大一，除非已为 $\F_p$：一个真子集不能在非零平移下不变。故 $p-2$ 对的选择和遍历 $\F_p$，加上两个单点得到 $p$ 项零和。
\end{proof}

第一次读这个证明，可以把它拆成两幅图景。前半段用 $z_i^{p-1}$ 表示“第 $i$ 项选不选”：它在零处为零，其余地方为一。第一方程迫使选项数是 $p$，第二方程迫使所选余数之和是零。有限域求和证明这样的选择确实存在。

后半段处理少一个候选的边界。把一些位置两两配对，每对可以选左边或右边。这提供一串二选一的自由度。每增加一对，可能的余数和至少增加一种；只要一开始比普通情况多出一种，就能覆盖全部 $p$ 个余数。唯一躲开这个过程的，是两类各 $p-1$ 份的形状。

到此，我们有了三件工具：留下余数见证保证安全；重复值保证下一轮还有材料；零和定理保证平均仍为整数。下一讲只负责把它们安排成一个覆盖所有状态的选择方法。

\chapter{让一次失败告诉我们下一次怎么选}\label{ch:descent-finish}
\section{分类之前，先看我们在怕什么}
现在有两个不同重值 $a,b$。如果完整留住 $p$ 份 $b$，另找一组做整数平均，唯一担心是新均值又等于 $b$。如果它不等于 $b$，我们就有新旧两个重值，任务完成。

这看似是麻烦，其实是额外信息。设 $a=b+\delta$，尝试
\[(a^i,b^{p-1-i},x).
\]
总和为 $pb+i\delta+x-b$。当 $p\nmid\delta$ 时，模 $p$ 可以唯一选择 $i\in\{0,\ldots,p-1\}$ 使平均为整数。如果均值恰为 $b$，就不只是一条同余，而是严格等式
\[x=b-i\delta.
\]
所以“所有尝试都失败”会强迫所有第三值落在一段有限等差列里。这时零和条件往往就不容许它们继续失败。

\section{副本多时：先拿背景值补足人数}
如果 $b$ 有至少 $2p-1$ 份，我们最多拿走 $p-1$ 份，仍留下一个完整 $p$ 块。因此它是很好的“补足人数”的背景值。

下面证明里出现的长选择式有一个简单来源。负方向的值是 $b-j_t\delta$，系数总和第一次达到 $T\ge p$。我们希望新均值为 $b-\delta$，即所选的相对系数和为 $-p$。已经拿了系数和 $-T$，就补 $T-p$ 份系数为 $1$ 的 $a$；余下位置用系数为零的 $b$ 填满。于是得到
\[(b-j_1\delta,\ldots,b-j_s\delta,a^{T-p},b^{2p-s-T}).
\]
不是凭空猜到这个式子，而是“先指定总和，再补足项数”。证明中的 $T+s\le2p$ 正好保证最后的补足数不为负。
\begin{lemma}[有充足背景副本时的续步]\label{lem:background}
有两个不同重值且其中一个重数至少为 $2p-1$ 时，存在保持不变量的安全非恒等整数平均，或配置可由相反值消去完成。
\end{lemma}
\begin{proof}
取背景 $b$、另一重值 $a=b+\delta$。若 $p\mid\delta$，选择 $(b^{p-1},a)$。否则对每个第三值 $x$，唯一选 $0\le i<p$ 使 $(a^i,b^{p-1-i},x)$ 整数平均。若均值不是 $b$，即告成功；否则 $x=b-i\delta$。因而失败时所有值属于 $b+\delta\{1,0,-1,\ldots,-(p-1)\}$，合法性给 $\gcd(\delta,n)=1$。
若另一个值也有 $2p-1$ 份，用两个这样的值调节一份第三值，两个旧重值均保留。因此设非 $b$ 重数均至多 $2p-2$，记正系数 $1$ 的重数为 $A$、负系数 $-j$ 的重数为 $c_j$，$S=\sum j c_j$。
若 $S<p$，零和给 $n\mid A-S$，而 $0<A-S<n$，矛盾。\par\noindent\textbf{跨越构造。}
若有 $j\ge2$，将负系数降序排列，取首次跨过 $p$ 的前缀 $j_1,\ldots,j_s$，总和 $T$。有 $2\le s\le p-1$，且 $T+s\le2p$：由 $(s-1)j_s\le p-1$ 和 $(s-2)(p-s)\ge0$ 得 $j_s+s\le p+1$，再用 $T\le p-1+j_s$。选择
\[(b-j_1\delta,\ldots,b-j_s\delta,a^{T-p},b^{2p-s-T})
\]
这个选择的项数为
\[
s+(T-p)+(2p-s-T)=p,
\]
而相对于 $b$ 的系数总和为 $-T+(T-p)=-p$，所以均值正好是 $b-\delta$。
所有计数都非负：$T\ge p$ 与 $T+s\le2p$ 分别保证后两项。
它留下至少 $p$ 份 $b$ 以及一份 $a$；二者在所有危险处不同余，故安全。
剩下仅有负系数 $-1$。若其重数为 $C$，则 $n\mid A-C$，且 $|A-C|<n$，故 $A=C,b=0$。反复选择 $(\delta,-\delta,0^{p-2})$ 完成。
\end{proof}

\section{重值种类多时：数的大小顺序就够用了}
如果有四个不同重值，先把最大的一个留在外面，只用另外三个。无论怎样平均，新均值都严格小于那个最大值，碰撞自动被排除。

只剩三个重值时，位置稍紧，需要先看有没有轻值、极端值有没有多出一份。最后的困难形状是两端各恰有 $p$ 份，中间有 $p+r$ 份。先拿一份中间值；如果碰撞，再拿两份。两次碰撞方程不能同时成立，除非已经是相反数加零的终端。
\begin{lemma}[三个及更多重值的续步]\label{lem:heavy}
当 $n\ge3p+2$ 时，至少三个不同重值的配置具有保持不变量的安全整数平均或相反值终端。两个重值同余模 $p$ 的情形在 $n\ge3p+d$ 时同样成立。
\end{lemma}
\begin{proof}
至少四个重值时，取 $a_1<a_2<a_3<b$。前三者同余则选 $(a_1^{p-1},a_2)$；否则以两个不同余值调节第三个，选 $i$ 使 $(a^i,c^{p-1-i},x)$ 整数平均。新均值小于未动的 $b$，并由引理\ref{lem:safety}安全。
恰三重值 $a<b<c$ 时，两者同余可用两个相反配比，均值之差为 $(p-2)(a-b)/p\ne0$，至少一个避开第三重值。三者不同余且有轻值 $x$ 时，以位于 $x$ 同侧的两重值调节，保留另一极端；若某极端有 $p+1$ 份，则只取它一份，用另两者调节。
\par\noindent\textbf{三个值的最后形状。}
只剩 $(a^p,b^{p+r},c^p)$，$r\ge2$。取 $i$ 使 $(a^i,c^{p-1-i},b)$ 整数平均。若均值为 $b$，则
\[i(a-b)+(p-1-i)(c-b)=0,\qquad1\le i\le p-2.
\]
令 $j\equiv2i\pmod p$。$j\le p-2$ 时改取 $(a^j,c^{p-2-j},b^2)$；再次碰撞会给 $(p-1)j=(p-2)i$，与 $j=2i$ 或 $2i-p$ 矛盾。若 $j=p-1$，则 $i=(p-1)/2$，从而 $a+c=2b$，零和给 $b=0,a=-c$。
\par\noindent\textbf{两个重值模 $p$ 相同。}
最后设仅有同余重值 $a,b$。有富余副本时用 $(a^{p-1},b)$。否则二者各 $p$ 份，保护后的轻值池至少有 $p$ 个位置；对其减去 $a$ 的部分和取碰撞，得到非空 $T$、$s=|T|\le p$、$\sum_{x\in T}(x-a)\equiv0\pmod p$。分别用 $a,b$ 补足 $p$ 项，所得均值之差为 $(p-s)(a-b)/p$，不可能同时碰撞到另一个重值，否则 $s=2p$。至少一个选择保持不变量。
\end{proof}
这个引理后半还处理“两个重值模 $p$ 相等”的情况。它用的部分和碰撞是抽屉原理：$p+1$ 个部分和只有 $p$ 种余数，总有两个相同；它们之间那一段就是一个非空零和片段。

\section{两边各多一点，也可以合起来用}
不必总要求某一个重值单独有 $2p-1$ 份。设它们分别有 $p+\alpha,p+\beta$ 份，我们把 $\alpha,\beta$ 叫作富余份数。如果一次总共只从两边拿 $p-2$ 份，而 $\alpha+\beta\ge p-2$，就不可能同时把两边都拿到不足 $p$ 份。

这个简单计数允许两个重值共同承担背景值的工作。下面的等差结构依然由“假设所有单轻值选择都碰撞”推出来，再由零和排除。
\begin{lemma}[两个重值共享富余副本]\label{lem:two-heavy}
若恰有两个不同余重值，其重数为 $p+\alpha,p+\beta$，$0\le\alpha,\beta\le p-2$、$\alpha+\beta\ge p-2$，则存在安全整数续步。
\end{lemma}
\begin{proof}
令 $\delta=a-b$。所有单轻值调节失败时，其他值被迫为
\[
b-i\delta\quad(\alpha+1\le i<p),\qquad
a+j\delta\quad(\beta+1\le j<p).
\]
故 $\gcd(\delta,n)=1$。设上下侧系数和为 $H,L$，上侧位置数为 $C$。零和给 $n\mid f(a)+C+H-L$；若 $H,L<p$，则
\[
f(a)+C+H-L\ge f(a)-(p-1)>0,
\]
另一方面，$f(a)+C\le n-f(b)$ 给出
\[
f(a)+C+H-L\le n-f(b)+p-1<n.
\]
这个严格介于 $0,n$ 的整数不可能被 $n$ 整除，矛盾。设下侧总和至少 $p$，取引理\ref{lem:background}的跨越前缀，平均
\[(b-i_1\delta,\ldots,b-i_s\delta,a^{T-p},b^{2p-s-T}).
\]
系数全为 $1$ 且总和达到 $p$ 会产生第三重值，故 $s\le p-1$。两锚点总消耗为 $p-s\le p-2\le\alpha+\beta$，至少一者保留 $p$ 份；新均值为 $b-\delta$，不同于二者。各锚点均留副本，保证安全。上侧情形对称。
\end{proof}

\section{富余也不够时：借出一份，再借出两份}
最后只剩两个不同余重值，并且富余总数最多 $p-3$。此时先各留 $p$ 份，再保护最多 $d_p(n)$ 个位置。如果
\[n\ge4p-2+d_p(n),
\]
还能拿出 $2p-2$ 个候选。上一讲说，若找不到 $p$ 项零和，这个池只能分成两类，各 $p-1$ 份。

现在借一份锚点 $a$ 进来，用两类余数补足。若所有实际选择都碰撞到 $b$，在同一余数类内交换“拿了谁、没拿谁”，碰撞等式会强迫同类位置的实际数值都相等。接着借两份 $a$，便能打破碰撞。若唯一的对称选择不允许借两份，就改从 $b$ 端做；两端不可能同时处于那个坏位置。
\begin{lemma}[整数下降覆盖的维数范围]\label{lem:integer-range}
若 $p\ge5$ 且 $n\ge4p-2+d_p(n)$，则每个合法本原输入可归零。特别所有 $n\ge4p+\lceil\log_2p\rceil$ 均满足判据。
\end{lemma}
\begin{proof}
引理\ref{lem:initialize}建立不变量。由引理\ref{lem:background}、\ref{lem:heavy}和\ref{lem:two-heavy}，只剩两个不同余重值、$\alpha+\beta\le p-3$ 的情形。若轻值 $x\equiv a\pmod p$，选择 $(a^{p-1},x)$；碰撞只可能 $x=pb-(p-1)a$。若它有两份，改用 $(a^{p-2},x^2)$，均值为 $2b-a$，仍安全。故每个锚点余类至多有一个轻值位置。
\par\noindent\textbf{建立边界池。}
预留 $p$ 份 $a,b$ 及至多 $d$ 个危险见证，在余池取 $2p-2$ 项并包含全部锚点富余副本。若有 $p$ 项零和，则其平均完成续步。否则引理\ref{lem:egz}给两个各有 $p-1$ 项的余类。$\alpha>0$ 会强迫 $\alpha\ge p-2$，矛盾；同理 $\beta=0$。故池的两类都不同于锚点余类。
\par\noindent\textbf{借用锚点并排除碰撞。}
取一份 $a$，由两池余类补足为整数平均；两类选取数均在 $1,p-2$ 内。如果所有此类选择均碰撞到 $b$，交换同类的选中与未选位置，说明两类各为实际常值 $u,v$，且 $a+iu+(p-1-i)v=pb$。若 $i\ne(p-1)/2$，用两份 $a$ 和 $j\equiv2i\pmod p$ 份 $u$，余下取 $v$，均值为 $2b-u$ 或 $2b-v$，不再碰撞。否则从 $b$ 端重复构造；两端的 $i$ 不可能相同，因为 $a\not\equiv b\pmod p$。
\par\noindent\textbf{终止。}
以上情况已经穷尽当前状态。只要还没到相反值终端，总有下一次安全操作。
每次非终端操作都是非恒等整数平均，且
\[\sum_i x_i^2-\sum_i(T_Sx)_i^2=\sum_{i\in S}(x_i-\bar x_S)^2>0.
\]
因此整数能量有限下降，最终进入已给出的相反值终端。对 $pX$ 执行所得词，再由引理\ref{lem:affine}回到 $X$。
令 $b=\lceil\log_2p\rceil$。因 $4p+b<2^{b+3}$，有 $4p+b-\lfloor\log_2(4p+b)\rfloor\ge4p-2$。函数 $n-\lfloor\log_2n\rfloor$ 不减，而 $d_p(n)\le\lfloor\log_2n\rfloor$，得到最后结论。
\end{proof}

这里终于可以使用能量终止：已经证明每个非终端都有下一步，而且每步都使非负整数 $E$ 严格减小。因此过程不会无路可走，也不能无限走下去。这两件事都要有，单独一句“能量下降”还不够。

我们得到的是一块宽维数范围。它已经很大，但尚未达到 $2p+1$。后面会看到，位置紧张时该换一种办法：允许分数，用精确输送代替每步整数下降。

\chapter{为什么造出一组零，就已经快结束了}\label{sec:blocks}
\section{先解决十五元五平均，看看矩阵从哪里来}
设十五个位置已是 $(a^5,b^5,c^5)$，且 $a+b+c=0$。按照下面的表重新分成三组：
\[
K_+=\begin{pmatrix}0&3&2\\1&2&2\\4&0&1\end{pmatrix}.
\]
第一行说“从 $a$ 块拿零份，从 $b$ 块拿三份，从 $c$ 块拿两份”；每行和为五，所以每行可以做一次五平均。每列和也为五，所以恰好用完每个旧块，没有多拿、漏拿或复制位置。

三个新值就是 $K_+(a,b,c)^T/5$。令 $x=a,y=a-b$，则 $b=x-y,c=y-2x$。代入第一行，得到
\[a'=-\frac{x+y}{5};
\]
第二行给 $b'=-x/5$，所以 $y'=a'-b'=-y/5$。这三次真实平均，在两个参数上就是
\[(x,y)\longmapsto-\frac15(x+y,y).
\]
暂时忽略整体比例，它把第一坐标加上第二坐标，这叫\textbf{上剪切}。

另一张表
\[K_-=\begin{pmatrix}2&1&2\\3&0&2\\0&4&1\end{pmatrix}
\]
实现 $-(x-y,y)/5$。两个词接起来，得到 $(x,y)/25$。参数的比例回来了，但真实数值整体缩小了；因此“方向上的逆”完全可以由正向平均实现。

\section{一般的三个等大块}
一般 $p\ge5$ 时写 $p=3m+\epsilon$、$\epsilon=\pm1$。在零和空间上，全一矩阵 $J$ 的作用为零。于是我们可以先设计想要的作用 $I\pm D$，再加 $mJ$，把所有计数变成非负。加上去的部分对零和参数没有影响，却让抽象矩阵成为真的分组。

下面会出现 $\SL_2(\Z)$，它只是所有行列式为一的二阶整数矩阵。$U(k)$ 把 $(x,y)$ 变为 $(x+ky,y)$；$L(k)$ 把它变为 $(x,y+kx)$。欧几里得算法正是在反复做这种整数加减。

$\Gamma_0(3,\Z)$ 是这些矩阵中左下项被三整除的那些。证明中的四陪集表，你可以看成一张四个状态的路线表：读入一个生成元后，记录落在哪个状态，并把差额记成已有的 $U,L$。读完整个矩阵词，就得到可以实际执行的剪切词。
\begin{proposition}[3p个位置的独立解法]\label{prop:threep}
若 $p\ge5$ 为素数，则 $3p$ 个位置满足定理\ref{thm:prime}的判据。
\end{proposition}
\begin{proof}
先将位置分为三个 $p$ 组。若其和模 $3$ 全相等，交换不同组中一对不同余值，组和变为 $s+\delta,s-\delta,s$，其中 $\delta\ne0$。这样的跨组数对必存在，否则原配置模 $3$ 全相等。分别平均后得到 $(a^p,b^p,c^p)$，$a+b+c=0$，且模 $3$ 不全相等。

写 $p=3m+\epsilon$，$\epsilon=\pm1$、$m\ge2$。令 $J$ 为三阶全一矩阵，
\[
D=\begin{pmatrix}1&-1&0\\1&-1&0\\-2&2&0\end{pmatrix},\qquad
K_\pm=mJ+\epsilon(I\pm D).
\]
两表非负且行列和为 $p$。每个输出组按一行从三个旧块取数，三次平均实现 $K_\pm/p$。在零和空间上，它们等于 $\epsilon(I\pm D)/p$。由于 $D^2=0$，有
\[
(I+D)(I-D)=I-D^2=I.
\]
实际三步平均仍会带上 $1/p$ 的共同尺度，但两作用在共同尺度意义下互为逆。
因此两个方向的整数次剪切都有正向平均序列。
改用 $(x,y)=(a,a-b)$，它们分别为 $U(1),U(-1)$。块角色交换给
\[
\Sigma=\begin{pmatrix}1&-1\\0&-1\end{pmatrix},\qquad
P=\begin{pmatrix}-2&1\\-3&2\end{pmatrix},\qquad
L(-3)=U(-1)P\Sigma.
\]
因此 $U(\Z),L(3\Z)$ 可实现。


这里 $\Gamma_0(3,\Z)$ 表示行列式为 $1$、左下项被 $3$ 整除的整数矩阵群。
为什么恰好出现 $3$？因为我们正在控制三个等大块，
其模 $3$ 的差分就是剩余的共同余数障碍。
这些矩阵射影生成 $\Gamma_0(3,\Z)$。为验证，取 $S=\left(\begin{smallmatrix}0&-1\\1&0\end{smallmatrix}\right)$、$T=U(1)$，右陪集代表为 $I,ST^{-1},S,ST$，由模 $3$ 底行的四个射影类区分。
陪集表的含义是：从代表 $r$ 出发，右乘一个生成元 $g$，
再乘新陪集代表 $r'$ 的逆，便回到子群。
任意群词可以逐字这样改写为这些回路的乘积，
因此只要表中的回路都由我们已有的矩阵生成，整个子群就由它们生成。
对生成元 $S,T,T^{-1}$，Schreier 回路 $rg(r')^{-1}$ 如下：
\[
\begin{array}{c|ccc}
r&S&T&T^{-1}\\\hline
I&I&T&T^{-1}\\
ST^{-1}&-T&I&L(3)\\
S&-I&I&I\\
ST&T^{-1}&L(-3)&I
\end{array}
\]
欧几里得消元证明 $S,T$ 生成 $\SL_2(\Z)$，上述表便给所需子群；中央符号不影响式\eqref{eq:return}。

本原三块的 $G=\gcd(y,3)$。合法性给 $3\nmid y$，可在 $\alpha x+\beta y=1$ 的裴蜀解中令 $3\mid\alpha$。矩阵
\[
\begin{pmatrix}y&-x\\\alpha&\beta\end{pmatrix}\in\Gamma_0(3,\Z)
\]
将 $(x,y)$ 送到 $(0,1)$，故得到 $(0^p,s^p,(-s)^p)$。
到了相反值配置，令 $h=(p-1)/2$。第一次选 $h$ 份 $s$、$h$ 份 $-s$ 和一个零，
得到 $p$ 个新零；第二次作同样的选择以后，每个符号只剩一份。
最后一组用这两份相反值和 $p-2$ 个零即可。
具体地，连续两次平均 $(s^h,(-s)^h,0)$，最后平均 $(s,-s,0^{p-2})$。
\end{proof}

\section{一个零给我们留了什么余地}
在 $2p+1$ 个零和位置中，如果有一个零，把它留着，其他 $2p$ 个位置正好分两组平均。两个组和相反，所以得到 $(a^p,(-a)^p,0)$。

例如 $p=5$，先拿两份 $a$、两份 $-a$、一个零平均，产生五个零；再做一次，只剩一对 $a,-a$；最后拿这对和三个零。只需三次就消去两个五重块。
\begin{lemma}[含零的2p+1元配置]\label{lem:zero-base}
对奇数 $p$，$2p+1$ 元零和配置若含一个零，则至多五次 $p$ 平均归零。
\end{lemma}
\begin{proof}
保留零，将其余位置分成两个 $p$ 组并平均，得到 $(a^p,(-a)^p,0)$。用命题\ref{prop:threep}末尾所列的三次相反值操作完成。
\end{proof}

\section{为什么可以在较小子系统里收尾}
一个零和 $p$ 元组平均以后，至少出现 $p$ 个零，所以非零位置最多有 $n-p$ 个。如果我们选一个子系统，\emph{包含全部非零位置，再多带一个零}，它的总和仍然是零，目标与整个配置一致。

含零还带来一个算术好处：所有坐标差的最大公因数等于坐标本身的最大公因数。因为差里面就包含 $x_i-0=x_i$。本原化后，这个子系统的 $G$ 必为一。

因此，只要这个子系统的大小属于已经证明的范围，就能在里面完成，外面的零保持不动。下面按 $n$ 的范围选择 $2p+1,3p,4p$ 等已解决的大小。
\begin{lemma}[从一个零和组到全部归零]\label{lem:zero-tail}
若 $p\ge5$，$n\ge2p+1$ 的零和配置含一个零和 $p$ 元组，则可归零。对任意奇数 $p$，此结论在 $2p+1\le n\le3p$ 同样成立。
\end{lemma}
\begin{proof}
先平均该组，产生 $p$ 个零。若某个已解决的维数 $m\le n$ 满足非零支持数小于 $m$，就选一个 $m$ 元子系统包含全部非零位置及一个零。它自身零和，本原化后的差分公因子为 $1$，可在该子系统归零。
对 $n\le3p$，使用引理\ref{lem:zero-base}中的 $m=2p+1$。对 $3p<n<4p$，使用命题\ref{prop:threep}。对 $4p\le n<5p$，使用引理\ref{lem:integer-range}在 $m=4p$ 的结论，因为 $d_p(4p)=1$。对 $n\ge5p$，同一引理的最终上界已适用。
\end{proof}

后面每次说“送到终端”，具体目标都是先让某个 $p$ 元组的和恰好为零。其余操作由本讲接手。这样我们不必要求一个可逆矩阵把非零列直接变成零；它也确实做不到。

\chapter{把很多数整理成三个块，同时留下余数差异}\label{sec:compression}
\section{位置不够了，换一个工作目标}
如果 $n=2p+1$，两个 $p$ 块已经用掉 $2p$ 个位置，只剩一个单点。前面那种“把两块都留住，再到外面找一组”的方法没法继续。

现在我们允许平均产生分数。目标是把所有位置整理成
\[(A^p,B^p,C^r),\qquad n=2p+r.
\]
这还没有解决问题，但复杂度下降了：三种数值受一个零和方程约束，只剩两个自由参数。我们可以集中研究这两个参数怎样变化。

难点是不能随意分组。三平均的 $(3^5,1,-16)$ 已经告诉我们，选错一组会毁掉唯一的模七见证。因此，下面先把余数差异放在容易保护的位置上，再合并其他位置。

\section{同一个整数，要同时躲开几个禁区}
两块互补混合：第一组选 $p-k$ 份 $a$、$k$ 份 $b$，另一组选剩余的互补份数。结果是
\begin{equation}\label{eq:mix-pair}
a_k=a+\frac{k}{p}(b-a),\qquad b_k=b+\frac{k}{p}(a-b).
\end{equation}
固定参照值 $w$。模某个危险素数 $\ell$，$a_k-w$ 是 $k$ 的一次式。只要 $a-w,b-w$ 不同时为零，这个一次式要么一直非零，要么只禁止一个 $k$ 余数。

这时中国剩余定理只能告诉我们某个整数存在，却不保证它在 $1,\ldots,p-1$ 里。实际的 $k$ 是取用份数，不能无限大。因此还要用 $n<4p$ 提供的素因子数量限制，证明这个短区间中就有一个同时可用的 $k$。
\begin{lemma}[短区间中的同时避让]\label{lem:short-choice}
若 $p\ge5$、$2p<n<4p$，对 $n$ 的任意一组不同素因子，各禁止一个剩余类，则 $1,\ldots,p-1$ 中存在同时避开的整数。
\end{lemma}
\begin{proof}
设素因子数为 $h$，乘积为 $M$，$\rho=\varphi(M)/M$。对 $h\le4$，长度 $2^h$ 足够：含 $2$ 时先留一种奇偶类，其余奇素数至多删去 $0,1,3,7$ 项，对应剩余 $1,2,4,8$ 项；不含 $2$ 时至多删去 $1,3,7,15$ 项。$h=3,4$ 时 $M\ge30,210$，从 $M\le n<4p$ 得 $2^h\le p-1$，较小 $h$ 直接成立。
若 $h\ge5$，CRT 将禁止类写成某个整数 $b$ 的各模余数。容斥给避开数目大于 $(p-1)\rho-(2^h-1)$。五个最小素数给 $\varphi(M)\ge480>2^{h+2}$，每增加素因子仍保持此界。因此
\[(p-1)\rho>(n/4-1)\rho\ge\varphi(M)/4-\rho>2^h-1.
\]
\end{proof}
证明中的 $\rho=\varphi(M)/M$ 可以先理解成“模 $M$ 的整数里，与 $M$ 互素的比例”。容斥给出主项 $L\rho$，再减去一个可控制的边界误差。如果主项比误差大，至少剩一个可用整数。

\section{把分散的差异集中到一对数上}
例如某块 $a$ 与单点 $w$ 已在模七不同，但模十一还相同；另一块 $b$ 在模十一与 $w$ 不同。用互补混合，选一个同时避开两个坏配比的 $k$，就能使同一个 $a_k-w$ 在两个素数处都非零。

如果差异在一个单点 $x$ 上，先平均 $(b^{p-1},x)$。留下的一份旧 $b$ 接替单点角色，新值接替 $b$ 块。对应二元变换为
\[(b,x)\mapsto\left(\frac{(p-1)b+x}{p},b\right).
\]
它模危险素数可逆，又固定常值配置，所以不会把原有的整体差异全部抹掉。
\begin{lemma}[把同余见证集中到一个差分]\label{lem:concentrate}
已有两个 $p$ 块及单点时，可以固定一个单点 $w$，使某个块值 $a$ 在全部危险素数处不同于 $w$。
\end{lemma}
\begin{proof}
若某个缺失素数处第二块 $b\ne w$，用式\eqref{eq:mix-pair}合并支持。否则 $a=b=w$，另有单点 $x\ne w$。平均 $(b^{p-1},x)$，留下旧 $b$ 一份作为新单点，再混合两块。这一阶段至多三次平均，严格增加支持。块--单点变换
\begin{equation}\label{eq:exchange}
(b,x)\longmapsto(((p-1)b+x)/p,b)
\end{equation}
在每个危险素数处可逆且保持常值方向，故保持整体见证。至多 $d_p(n)$ 个阶段完成。
\end{proof}

\section{最紧的情况，怎样先造出两个块}
还不能默认我们一开始已有两个块。先保护各素数的不同余位置，造出一个块。剩下是多个单点。对任一单点，刚才的块--单点交换都可以执行。

接下来的交换子公式看起来突然，其目的很简单：让一个目标单点的差值增加或减少另一个单点差值的倍数，并让其余多数差值保持不动。它就是高维坐标中的一次“加倍数”操作。

此时只需要在模 $\rad(n)$ 下得到这个作用。有限可逆矩阵的逆等于它的某个非负次幂，所以公式里的逆可以通过多做正向交换得到相同的余数作用。这不把实际有理数值倒回去；入口阶段只需要余数正确。
\begin{lemma}[保留见证的双块入口]\label{lem:two-block-entry}
设 $1\le r<p$、$n=2p+r$、$\gcd(p,n)=1$。任意合法输入可进入具有完整见证的两个 $p$ 块及 $r$ 个单点。
\end{lemma}
\begin{proof}
固定一个位置并为每个危险素数保留一个不同余位置。保护数至多 $d+1\le p+r$：否则 $n\ge2^d\ge2^{p+r}>n$。保护集外平均一个 $p$ 组，写剩余单点为 $c_i$、$m=p+r$、$\delta_i=c_i-a$。令 $\lambda=-1/p$，式\eqref{eq:exchange}在差分中的矩阵为
\[
A_i=I+v_i e_i^T,\qquad v_i=\lambda\mathbf1-e_i.
\]

为了说明交换子的由来，取互异 $i,j,k$，置
\[
w=e_k-e_j,\qquad h=e_j^T+\lambda e_k^T,
\qquad T=A_k^{-1}A_jA_k.
\]
秩一矩阵公式给出 $A_k^{-1}v_j=w$、$e_j^TA_k=h$，故 $T=I+wh$。
另外 $e_i^Tw=0$、$hv_i=\lambda(1+\lambda)$。
将这两个等式用于共轭，再乘 $A_i^{-1}$，得到
\begin{equation}\label{eq:entry-commutator}
[A_k^{-1}A_jA_k,A_i]
=I+(1+\lambda)(e_k-e_j)e_i^T.
\end{equation}
取 $L$ 为各模 $\ell$ 的 $2\times2$ 可逆矩阵群阶的公倍数。
这里 $\rad(n)$ 是 $n$ 的全部不同素因子的乘积。
按实际时间顺序，用下面八段正向操作实现交换子的模作用：
\[
 A_i^{L-1},\ A_k,\ A_j^{L-1},\ A_k^{L-1},\
 A_i,\ A_k,\ A_j,\ A_k^{L-1}.
\]
这些段一共含 $4L$ 次块--单点交换。
交換矩阵模 $\rad(n)$ 的逆可由 $L-1$ 次正向交换实现，故式\eqref{eq:entry-commutator}的任意非负幂可用真实交换实现其模作用。
系数 $(p-1)/p$ 失效的素数整除 $\gcd(p-1,n)=\gcd(p-1,r+2)$，其数量至多 $r$。
确实，若有 $r+1$ 个这样的素数，则其乘积至少为 $2^{r+1}$，又整除 $r+2$，
与 $2^{r+1}>r+2$ 矛盾。这个小计数解释了保留 $r$ 个单点为什么足够：
不能用交换子移动的素数，恰好可以逐个留一个原有见证。为这些素数分别选一个非零差分单点，补足为 $r$ 个保留角色，固定其中一个目标 $j$。对每个其余素数 $\ell$，若目标差为零，选 $\delta_i\ne0$、$k\notin\{i,j\}$，用 CRT 选指数在模 $\ell$ 为 $1$、其他非失效素数模下为 $0$。式\eqref{eq:entry-commutator}使目标差新增非零值，并保留已处理支持和全部失效处数据。最终保留这 $r$ 个单点，将另外 $p$ 个单点平均。旧块和保留单点已提供全部见证。
\end{proof}
最关键的容量配合是：交换子失效的素数都整除 $r+2$，其个数不超过剩余单点数 $r$。失效处用单点保存，其他处用交换子搬运。因此即使 $n$ 有很多素因子，也不需要为每个素因子永久占据一个额外位置。

\section{为什么核心要写成这个样子}
设三块零和：$pA+pB+rC=0$，且 $\gcd(p,r)=1$。在整数代表中 $p\mid C$，写 $C=-pz$，令 $A=a$，便得到
\[B=rz-a.
\]
所以
\[K_{p,r}(a,z)=(a^p,(rz-a)^p,(-pz)^r)
\]
只是零和方程的一个方便参数化。令 $y=a+pz=A-C$，另外的块差可写成 $nz-2y$。本原性给 $\gcd(y,z)=1$，于是
\[G=\gcd(y,n).
\]
我们将来要控制的“第二坐标是模 $n$ 单位”，正是最初的 $G=1$ 条件。
\begin{proposition}[一般输入的三块归约]\label{prop:entry}
若 $p\ge5$、$2p<n<4p$、$n\ne3p$，每个合法输入可进入
\begin{equation}\label{eq:core}
K_{p,r}(a,z)=(a^p,(rz-a)^p,(-pz)^r),\qquad r=n-2p.
\end{equation}
本原参数满足 $\gcd(a,z)=1$，且合法性为 $\gcd(a+pz,n)=1$。当 $p=3,n=7$ 时该结论仍成立。
\end{proposition}
\begin{proof}
\par\noindent\textbf{第一步：下半带两轮修复。} 先设 $r<p$。由引理\ref{lem:two-block-entry}及\ref{lem:concentrate}取得 $a^p,b^p$ 和单点 $w$，其中 $a-w$ 在全部危险处非零。混合两块后将 $b_k$ 与 $w$ 交换，得到
\[
a_1=a_k,\quad b_1=((p-1)b_k+w)/p,\quad w_1=b_k.
\]
新块差乘 $p$ 为 $D(k)=pa-(p-1)b-w+(2p-1)k(b-a)/p$。在 $\ell\nmid2p-1$ 处选 $k$ 使它非零；在 $\ell\mid2p-1$ 处选 $k$ 使下一轮常数
\[E(k)=pa_1-(p-1)b_1-w_1=w+(a-3b)/2+4k(b-a)
\]
非零。两种条件均至多禁止一个类。第二轮的差为 $E(k)+(2p-1)h(b_1-a_1)/p$，再由引理\ref{lem:short-choice}同时避零。现在两块差在全部危险处非零。平均全部 $r$ 单点及第一块的 $p-r$ 份，留下 $r$ 份旧第一块和完整第二块，即得合法形状\eqref{eq:core}。

\par\noindent\textbf{第二步：上半带的三块折叠。}
若 $n=3p+s$、$1\le s<p$，先保护单点 $w$ 及至多 $d$ 个见证，在其外制造两个 $p$ 块。$d<p$ 由 $2^p>4p>n$ 保证容量。用引理\ref{lem:concentrate}取得 $a-w$ 的全支持，再避开 $w$ 制造第三块。混合第二、第三块以保留已好素数的差，再将第二块与 $w$ 交换，可使三个块自身提供全部见证。将单点与第一块折叠，得到 $(a^p,b^p,c^p,d_0^s)$。
零和模危险素数给 $(a-d_0)+(b-d_0)+(c-d_0)=0$。先混合使 $a-d_0$ 全支持。写 $p=3m+\epsilon$，平均两组同组成 $(m+\epsilon,m,m)$ 的位置，再平均余下的 $p$ 个，得到 $(A^{2p},B^p,C^s)$，其中
\[
A=[m(a+b+c)+\epsilon a]/p,\quad B=a+b+c-2A,\quad C=d_0.
\]
模危险素数有 $A-C=\epsilon(a-d_0)/p\ne0$。写 $A=a,B=-2a+sz,C=-pz$，再应用引理\ref{lem:shape-change}的两次平均，得到 $(p,p,p+s)$ 核心。
\par\noindent\textbf{第三步：把块值写成两个整数参数。}
对整数本原核心，由 $p(A+B)+rC=0$ 及 $\gcd(p,r)=1$ 得 $C=-pz$，进而 $B=rz-a$。坐标公因子为 $\gcd(a,z)$，而差分为 $a+pz$ 和 $nz-2(a+pz)$，所以 $G=\gcd(a+pz,n)$。
当 $p=3,n=7$ 时，引理\ref{lem:two-block-entry}已经给两个三块与一个单点。保留单点与旧块的全见证便给所需合法核心，不需后续两轮修复。
\end{proof}

\section{上半带多出来的一个块，怎样换成大载体}
当 $n=3p+s$ 时，最容易先得到的是 $(2p,p,s)$ 重数：三个完整 $p$ 块先折叠出一个 $2p$ 块。后面的共同矩阵工具则更喜欢两个等大的 $p$ 块，因此要把它换成 $(p,p,p+s)$。

这里不用任何抽象“拆块定理”。先平均一组，从尚未用掉的旧 $A$ 中留下 $p+s$ 份，再平均其余 $p$ 份，一共两次。唯一需要选好的，是第一组从三种值各取多少。
\begin{lemma}[用两次平均改换块重数]\label{lem:shape-change}
设 $p\ge5$、$1\le s<p$、$n=3p+s$。合法核心 $(a^{2p},(-2a+sz)^p,(-pz)^s)$ 可用两次平均进入合法 $K_{p,p+s}$。
\end{lemma}
\begin{proof}
记三个块值为 $A,B,C$。选 $0\le u\le s$、$s\le t+u\le p$，令 $q=3t+u$ 与 $n$ 互素。先平均 $(A^{p-t-u},B^t,C^u)$，从剩余旧 $A$ 保留 $p+s$ 份，再平均其余的 $p$ 个位置。返回为
\[
\frac1p\begin{pmatrix}p-q&st-pu\\-1&0\end{pmatrix}.
\]
可取 $(t,u)=(s,0)$（$3\nmid s$）；若 $3\mid s$，令 $\rho=p\bmod3$，取 $u=s-3+\rho,t=(p-u)/3$，此时 $q=p$。容量 $t+u\ge s$ 等价于 $p-s\ge6-2\rho$；唯一例外为 $\rho=1,s=p-1$，改取 $t=(p+1)/2,u=(p-3)/2$，此时 $q=2p$。三种选择均合法且 $\gcd(q,n)=1$。
若 $y=a+pz$，新见证满足 $y'\equiv-qy/p\pmod n$，行列式为 $st-pu\equiv-pq\pmod n$。故在每个危险素数处保持本原性及见证。
\end{proof}
到这里，任意合法输入已经被送入一个覆盖全部位置的二维参数族。后面用矩阵工作时，别把两个参数误看成只剩两个位置；原来的 $n$ 个位置一个也没少。

\chapter{一段平均操作，为什么可以是一张可逆矩阵}\label{sec:arithmetic}
\section{把参数变化算一次，就知道矩阵从哪里来}
在核心 $K_{p,r}(a,z)$ 中，固定第三块，让前两块作互补混合。第一组拿 $p-k$ 份 $a$ 和 $k$ 份 $rz-a$，所以
\[a'=\frac{p-2k}{p}a+\frac{kr}{p}z,qquad z'=z.
\]
这就是矩阵
\[\binom{a'}{z'}=
\begin{pmatrix}(p-2k)/p&kr/p\\0&1\end{pmatrix}\binom az.
\]
没有增加任何新操作。矩阵只是把两条线性公式写在一起。只要一次或几次平均后重新出现同样的块重数，就能重新读出两个参数，继续用这种记账法。

我们把这样一段有限平均叫作\textbf{返回}，也叫宏。名字里的“返回”是块形状回来，不是数值回到原来。

\section{真正输出与方便的矩阵代表}
回想第5讲的五平均例子，真实输出是 $-(x+y,y)/5$。写成 $U(1)(x,y)$ 更容易看出它是剪切，但实际配置还带着 $-1/5$。

为何可以研究 $U(1)$？因为我们最终要到零。某个后续位置序列能消去 $v$，它也能消去 $cv$，只要 $c\ne0$。因此，在完整配置上，共同尺度不影响可达性。
\begin{definition}\label{def:return}
设 $E_\pi:\Q^2\to\Q^n$ 按分组 $\pi$ 将两个参数展开为等值块配置。若平均词 $W$ 满足
\begin{equation}\label{eq:return}
WE_\pi=\lambda E_{\pi'}M,\qquad \lambda\in\Q^\times,
\end{equation}
且 $\pi'$ 与 $\pi$ 具有相同块重数，称 $M$ 是一个可实现返回。构造须对每个同形状分组适用。可逆矩阵的可实现性按非零共同标量识别。
\end{definition}
把 $E_\pi$ 想成“按照分组名单展开参数”：它将 $(a,z)$ 写到全部 $n$ 个位置。$\pi'$ 是做完后新的分组名单。等式左边先展开再做平均，右边先变参数再按新名单展开，两个结果相同。

\begin{example}[两次正向返回提供逆]
若 $J=\left(\begin{smallmatrix}0&d\\-1&0\end{smallmatrix}\right)$，$d\ne0$，则 $J^2=-dI$。
假如 $J$ 有真实平均词，重复这个词就使全部参数共同乘 $-d$（另计实际词自身的共同尺度）。
所以 $J^{-1}=(-d)^{-1}J$ 在射影意义上由同一条正向词实现。
这不表示一次平均在整个 $\Q^n$ 上可逆；可逆性只发生在带相同块形状的参数空间上。
\end{example}
这不是说某个一般平均投影忽然可逆了。一个任意三元输入平均后丢掉两个差分，当然无法恢复；但我们限制的块族原来就只有两个参数，而且只要求恢复整个配置的比例。额外位置使我们有机会把这两个参数的信息转存到别的块中。

\section{剪切为什么比一般矩阵好用}
记
\[U(t)=\begin{pmatrix}1&t\\0&1\end{pmatrix},qquad
L(t)=\begin{pmatrix}1&0\\t&1\end{pmatrix}.
\]
它们分别给第一坐标加 $t$ 倍第二坐标，或给第二坐标加 $t$ 倍第一坐标。最方便的性质是
\[U(s)U(t)=U(s+t),qquad U(t)^{-1}=U(-t).
\]
如果已经有 $U(2)$ 和 $U(3)$ 以及它们的逆，裴蜀关系 $3-2=1$ 就给 $U(1)$。构造里经常计算若干根参数的最大公因数，原因就在这里。“根子群”在本题中就是让同一类剪切的参数遍历某个加法群。

\section{局部化环不是另一套数，是允许的分母清单}
\begin{definition}[整数的有限局部化]
对整数 $D>1$，
\[
\Z[1/D]=\{a/D^k:a\in\Z,\ k\ge0\}.
\]
等价地，约分后的分母只含 $D$ 的素因子。
其单位是 $\pm\prod_{\ell\mid D}\ell^{e_\ell}$，其中指数 $e_\ell$ 可正可负。
\end{definition}
$\Z[1/6]$ 里有 $1/2,1/3,1/12$，因为它们都能写成 $a/6^k$；没有 $1/5$。数 $5$ 在这个环里，但它的逆不在，所以 $5$ 不是单位。数 $2$ 和 $3$ 则是单位。

在原数值层面，整数输入经 $p$ 平均只能出现 $p$ 的分母。但核心矩阵允许除去一个\emph{整个配置的}共同尺度后表示。它的参数环可能含别的分母。这两句话并不冲突：矩阵中的 $1/2$ 不等于实际坐标被单独免费除以二。

\section{用共轭把一个步长变成另一个步长}
设已取得可逆三角返回 $T$。组合 $TU(t)T^{-1}$ 的意义是：先用 $T^{-1}$ 改参数，再剪切，最后用 $T$ 改回来。计算告诉我们，这相当于换了剪切的步长。
\begin{lemma}[可逆倍率扩展剪切参数]\label{lem:localize}
假设 $U(hR)$ 全部可实现，且上三角矩阵 $T=\left(\begin{smallmatrix}1&b\\0&\kappa\end{smallmatrix}\right)$ 及其逆可实现。则 $U(hR[\kappa,\kappa^{-1}])$ 全部可实现。若 $\kappa=A/B$ 约分，则新环为 $R[1/(AB)]$。
\end{lemma}
\begin{proof}
使用 $TU(t)T^{-1}=U(t/\kappa)$、$T^{-1}U(t)T=U(\kappa t)$ 和 $U(s)U(t)=U(s+t)$。
新环中每个元素是有限个 $r_i\kappa^i$ 的和，其中 $i$ 允许为负数。
先用共轭实现每一项的剪切，再相乘把参数相加，得到所需返回。
全部求和都是有限的，因此最终仍是一条有限平均词。
至于新环中的分母为什么正好来自 $AB$，一方面 $A/B$、$B/A$ 显然都属于 $R[1/(AB)]$；
另一方面，裴蜀等式 $uA+vB=1$ 给出 $u\kappa+v=1/B$；对逆倍率同理得到 $1/A$。
\end{proof}
例如倍率为 $3/2$。它和逆倍率 $2/3$ 允许我们把参数乘这些数。更具体地，$3/2-1=1/2$，$1-2/3=1/3$，所以加法和这些倍率已经生成 $\Z[1/6]$。一般情形中，裴蜀等式承担同一任务。

到此若说“已能实现 $U(hR)$”，意思是对每个属于 $hR$ 的参数，都有某条有限平均词。并非一条词同时实现所有参数，也不只是模几个整数检查到了它们。下一讲将利用这整族剪切，把无界参数之间的精确输送交给群论。

\chapter{同余子群：怎样从有限余数回到精确数值}\label{ch:congruence}
\section{先看一个可以算到最后的输送}
给定两个本原整数列
\[v=\binom23,qquad w=\binom78.
\]
它们模五相同，但实际相差很大。矩阵
\[\gamma=\begin{pmatrix}-4&5\\-5&6\end{pmatrix}
\]
满足 $\det\gamma=1$、$\gamma\equiv I\pmod5$，并且 $\gamma v=w$。

这正是我们想要的工具：不改变指定余数，却能把实际向量精确送到另一个。定义
\[\Gamma(q,R)=\{M\in\SL_2(R):M\equiv I\pmod{qR}\}.
\]
$\SL_2(R)$ 是行列式为一的二阶矩阵群；“主同余子群”就是其中模 $q$ 等于单位矩阵的那些。这个名称虽然长，条件只是四个矩阵项的同余。

若 $R=\Z[1/D]$ 且 $\gcd(q,D)=1$，环里的分数怎样取模？把 $a/D^k$ 送到 $a(D^k)^{-1}\pmod q$。所以 $R/qR$ 与普通的 $\Z/q\Z$ 是同一个有限环。若 $q$ 含已经在 $R$ 中可逆的因子，先把它们去掉即可。例如在 $\Z[1/2]$ 中，模六的条件等价于模三，因为二已经是单位。

\section{为什么必须先取得这个群的真实实现}
模 $q$ 相同的两列有无穷多个。如果它们都能由 $\Gamma(q,R)$ 互相送到，而这个群的每个矩阵又都有真实平均词，那么无限问题就只剩有限问题：先送对余数，剩下由群精确完成。

不过，光在计算机里看到模 $q$ 的矩阵很多，并不能推出这个群已经可实现。模运算会把一大批不同有理矩阵压到一起，丢掉的部分正是我们需要控制的。这是两个外部定理真正进入证明的地方。

“有限指数”表示一个子群把整个群分成有限多个陪集，可以先把它理解成只剩有限种余类。Morris 的定理从足够多的初等矩阵得出有限指数；Serre 的定理再说，对我们的环，这样的子群一定包含某个主同余子群。
\begin{theorem}[Morris；Serre]\label{thm:arithmetic-external}
设 $D>1$、$R=R_D$。对于 $\SL_2(R)$ 的任一有限指数子群 $H$：
\begin{enumerate}
\item $H$ 中的上下初等矩阵生成 $H$ 的一个有限指数子群；
\item $H$ 包含某个非零层的主同余子群 $\Gamma(q,R)$。
\end{enumerate}
\end{theorem}
第一项使用 Morris\cite[Theorem 6.1(2)]{morris} 的二维无限单位情形。第二项使用 Serre\cite{serre} 的有理 $S$ 整数结论，准确陈述见\cite[pp.308--309]{prasad}。这里 $S$ 含实处和 $D$ 的素因子处，$D>1$ 保证至少有一个有限处。若 $D=1$，变成 $\SL_2(\Z)$，第二项一般就不成立。

这两项是本讲引用的深结果。后面从“不知道多大的模数”改进到明确的平方模数，则是我们自己能逐步算出的消元。

\section{为什么最后出现的是平方模数}
我们已经有上、下两族步长在 $hR$ 内的剪切。外部定理先保证一个未知 $\Gamma(q_0,R)$ 已包含其中。接下来想证明：任何 $g\equiv I\pmod{h^2}$ 也属于已生成的群。

取更深的模数 $q=h^2q_0$。在有限环 $R/qR$ 里，先给第一行加第二行的一个合适倍数，使左上角变成单位。然后像高斯消元一样消掉非对角项。剩下的对角项也能写成四个剪切，因为它与一的差被 $h^2$ 整除。

消元只保证与 $g$ 模 $q$ 相同。但剩余误差属于一开始已经取得的 $\Gamma(q_0,R)$，所以能够精确补掉。这就是“先有一个真群，再用有限运算改善它”的顺序。
\begin{lemma}[两条根生成一个确定的同余层]\label{lem:deep}
若 $D>1$、$h\ne0$，则
\[
\Gamma(h^2,R_D)\subseteq\langle U(hR_D),L(hR_D)\rangle.
\]
\end{lemma}
\begin{proof}

\par\noindent\textbf{先得到一个未知深层。}
记右侧为 $H$。将第一个外部定理用于有限指数子群 $\Gamma(h,R)$：
它包含的上下初等矩阵恰是参数在 $hR$ 中的剪切，故 $H$ 有限指数。
定理\ref{thm:arithmetic-external}先给出 $H$ 有限指数，再给出 $\Gamma(q_0,R_D)\subseteq H$。令 $q=h^2q_0$、$B=R_D/qR_D$。

\par\noindent\textbf{在有限环中找单位主元。}
对 $g=\left(\begin{smallmatrix}a&b\\c&d\end{smallmatrix}\right)\in\Gamma(h^2,R_D)$ 的模 $q$ 像，用 CRT 选 $k\in B$，使 $u=a+hkc$ 为单位：剩余特征整除 $h$ 时 $a=1$；其余分量若 $c\ne0$ 只排除一个 $k$ 类，若 $c=0$ 则 $a$ 已为单位。置 $v=b+hkd$，有
\[
g=U(-hk)L(c/u)\diag(u,u^{-1})U(v/u).
\]

从 $g\equiv I\pmod{h^2}$ 可知 $c\in h^2B$、$b\in h^2B$，
而 $v=b+hkd\in hB$。于是上式两侧的非对角消元都已经使用 $hB$ 中的参数。
还需拆开对角矩阵。由于 $u-1\in h^2B$，取 $t\in hB$ 使 $ht=u-1$。恒等式
\[
\diag(u,u^{-1})=U(h)L(t)U(-h/u)L(-tu)
\]
把中间项也分解为参数属于 $hB$ 的根。
\par\noindent\textbf{回到真实参数。}
约化 $hR_D\to hB$ 满射，因而可把有限个剪切参数逐个提升。
提升各参数至 $hR_D$，得到 $e\in H$ 与 $g$ 模 $q$ 相同。于是 $e^{-1}g\in\Gamma(q,R_D)\subseteq H$。
\end{proof}

\section{如果起初只有上剪切怎么办}
上剪切始终固定第一坐标轴。如果一个可逆返回 $J$ 把这条轴送到另一条直线，那么 $JU(t)J^{-1}$ 就固定那条新的直线。把两条直线当成坐标轴，就看到了上、下两种方向。

只是这个换基矩阵未必有整数逆，所以要把分母损失写进同余模数。下面的 $g^3$ 就是为这件事付出的明确代价。它不是多出三个物理位置，也不是三次平均的次数。
\begin{lemma}[横向返回与换基的分母损失]\label{lem:transverse}
假设 $U(hR)$ 以及 $J,J^{-1}$ 可实现，其中 $J=\left(\begin{smallmatrix}a&b\\c&d\end{smallmatrix}\right)\in M_2(R)$、$c\det J\ne0$。写 $c=g\epsilon$、$\epsilon\in R^\times$、$g\in\Z\setminus\{0\}$。则 $\Gamma(g^3h^2,R)$ 可实现。若有多个这样的 $g_i$，可将其层合并为 $\gcd(g_i)^3h^2$。
\end{lemma}
\begin{proof}
令 $P=(e_1,Je_1)=\left(\begin{smallmatrix}1&a\\0&c\end{smallmatrix}\right)$。相乘得
\[
P^{-1}U(t)P=U(ct),\qquad
P^{-1}JU(t)J^{-1}P=L(-ct/\det J).
\]
取 $t=hs$ 和 $t=\det J\,hs$，在 $P$ 基下得到两条完整 $ghR$ 根。引理\ref{lem:deep}给 $P\Gamma(g^2h^2,R)P^{-1}$。
因 $P,gP^{-1}$ 整于 $R$，若 $\gamma=I+g^3h^2X$，$X\in M_2(R)$，则
\[
P^{-1}\gamma P-I
=g^2h^2\,(gP^{-1})XP
\]
各项都属于 $g^2h^2R$，并且行列式仍为 $1$。
故 $P^{-1}\gamma P\in\Gamma(g^2h^2,R)$，证明该共轭群包含 $\Gamma(g^3h^2,R)$。
这个换基只用于论证群的包含关系，不需要把 $P$ 本身实现为平均词。
合并时先约化到各层的公倍数。在每个素数幂分量上，主同余条件嵌套，两个像生成较浅的条件；CRT 及初等矩阵提升给出各分量的乘积。接上已知公倍数层的核，即得最大公因子层。
\end{proof}

\section{相同余数的本原列，为什么真的能精确互达}
在 $R$ 中，“列 $(a,b)$ 本原”表示可以找到 $u,v\in R$ 使 $ua+vb=1$。在整数环这就是 $\gcd(a,b)=1$。它的线性代数含义是：这列能够补成一个行列式为一的二阶矩阵。

给起点 $v$ 补一列 $x$，给目标 $w$ 补一列 $y$。若能让两个补成的矩阵也模 $q$ 相同，那么
\[(w,y)(v,x)^{-1}
\]
便既把 $v$ 送到 $w$，又模 $q$ 等于 $I$。下面证明的全部工作，就是把第二列稍作修改，以满足这两个要求。
\begin{example}[同余相同，实际数值可以相差很大]
取 $v=(1,0)^T$、$w=(1,q)^T$，两者模 $q$ 相同。
下剪切 $L(q)$ 已经把 $v$ 精确送到 $w$。
一般本原列可能没有这么直接的一个剪切，但下一个引理保证某个 $\Gamma(q,R)$ 元素总能完成。
\end{example}
\begin{lemma}[相同剩余类中的精确输送]\label{lem:fiber}
若 $v,w\in R^2$ 本原且 $v\equiv w\pmod q$，则存在 $\gamma\in\Gamma(q,R)$ 使 $\gamma v=w$。
\end{lemma}
\begin{proof}
取 $x$ 使 $\det(v,x)=1$，写 $\det(w,x)=1+qk$。
写 $w=(w_1,w_2)^T$，则 $\det(w,z)=w_1z_2-w_2z_1$。
因为 $w_1,w_2$ 生成整个 $R$，这个线性式能够取到任意环元素。
也就是说，本原性使 $z\mapsto\det(w,z)$ 满射，故取 $z$ 使其值为 $-k$。令 $y=x+qz$，则 $(w,y)\equiv(v,x)\pmod q$，两矩阵行列式皆为 $1$。取 $\gamma=(w,y)(v,x)^{-1}$。
\end{proof}
这是整套方法里最值得停下来理解的一步：有限同余计算负责把你送进正确的“房间”，这个引理负责在房间内送到那个\emph{确切的点}。最终是等号，不是取极限。

\section{取得主同余子群以后，逆更容易造了}
在有限矩阵群里，某矩阵的正整数次幂一定回到 $I$。如果 $Z^m\equiv I\pmod q$，那么真实 $Z^m$ 就落进已能实现的 $\Gamma(q,R)$，其逆也已可用。因此
\[Z^{-1}=Z^{m-1}(Z^m)^{-1}.
\]
这把“有限商中的周期”变成了一条真正的逆词。

行列式不是一时，要额外处理整数内容。Smith 标准形是通过整数可逆行列变换把矩阵化成对角形的定理；这里用它设计一个两边夹乘，使结果能除掉行列式而仍是整数矩阵。再由 CRT 让这个夹乘同时模 $q$ 正确。
\begin{lemma}[获得深同余群以后的正向逆]\label{lem:inverse}
若 $\Gamma(q,R)$ 可实现，则每个可实现 $Z\in\SL_2(R)$ 的逆可实现。若整数满秩返回 $A$ 满足 $\gcd(\det A,q)=1$，则 $A^{-1}$ 亦有射影实现。
\end{lemma}
\begin{proof}
第一项取 $m$ 使 $Z^m\equiv I\pmod q$，用 $Z^{-1}=Z^{m-1}(Z^m)^{-1}$。
第二项取偶数 $m$ 使 $P=A^m\equiv I\pmod q$，令 $\Delta=\det P>0$。选 $L,R'\in\SL_2(\Z)$ 使 $LPR'=\diag(a,b)$、$ab=\Delta$。置 $J_0=\left(\begin{smallmatrix}0&1\\-1&0\end{smallmatrix}\right)$、$g_*=R'J_0L$，则 $Pg_*P$ 各项被 $\Delta$ 整除。
利用 CRT 及 $\SL_2(\Z)$ 对有限模数的满射，取 $g\equiv g_*\pmod\Delta$、$g\equiv I\pmod q$。则 $H=\Delta^{-1}PgP\in\Gamma(q,\Z)$，而 $gPH^{-1}=\Delta P^{-1}$。接上 $A^{m-1}$ 得结论。所用约化满射由有限局部环上的初等消元、CRT 和整数参数提升得到。
\end{proof}
第一次读到这儿，可先记住它的用途和条件：\textbf{已经取得真实主同余子群后，行列式与控制模数互素的整数返回，可以获得正向射影逆。}具体夹乘等式留在上面的完整证明中，第二遍逐项相乘核对即可。

\chapter{最后那段有限计算，到底在调整什么}\label{sec:controllers}
\section{只知道比例正确，还差一步}
模五的 $(1,1)$ 与 $(2,2)$ 比例相同，但列并不相同。上一讲的精确输送要求\emph{整列}同余。因此，先把比例调对以后，还得说明允许怎样改变共同倍数。

这不是要求实际数值变大。例如某段平均可以输出一个很小的共同倍数；把这个倍数从分析代表中取出后，本原列却可能变成另一个有限尺度。我们始终跟踪实际完整配置和分析代表之间的对应。

\section{一张整数矩阵，可以产生几种本原输出}
设输入列 $(a,b)$ 的最大公因数为一。施加整数矩阵后，两坐标可能有公因子。把这个最大公因数叫作\textbf{内容}，记为 $\cont$。

最简单的例子是 $A=\diag(1,6)$。输入 $(s,t)$ 本原，输出 $(s,6t)$ 的公因子是 $\gcd(s,6)$，所以只可能为 $1,2,3,6$。

甚至可以固定输入模五都等于 $(1,1)$，而让四种内容分别出现：
\[
\begin{array}{c|c|c}
\text{输入}&A\text{ 的输出}&\text{本原输出}\\ \hline
(1,1)&(1,6)&(1,6)\\
(16,1)&(16,6)&(8,3)\\
(21,1)&(21,6)&(7,2)\\
(6,1)&(6,6)&(1,1)
\end{array}
\]
它们分别除去了 $1,2,3,6$。所以同一个模五输入类，在施加同一个 $A$ 并取本原代表以后，会得到几种不同列。这就是“内容分支”的具体意思。
\begin{example}[一个内容分支的具体样子]
令 $A=\diag(1,6)$，控制模数取 $q=5$。
对本原列 $(s,t)^T$，内容为 $\gcd(s,6)$，可能取 $1,2,3,6$。
固定任意本原模 $5$ 类，在模 $2,3$ 处分别选择 $s$ 是否为零，并取 $t$ 为单位，
就能安排这四种内容；与模 $5$ 的条件由中国剩余定理合并。
本原输出模 $5$ 因而分别是 $Ax,2^{-1}Ax,3^{-1}Ax,6^{-1}Ax$。
一般矩阵先经 Smith 标准形化为这个例子的形式。
\end{example}
\begin{lemma}[本原化的全部有限分支]\label{lem:content}
设 $A\in M_2(\Z)$ 满秩，条目最大公因子为 $1$，$d=|\det A|$、$\gcd(d,q)=1$。固定模 $q$ 本原列 $x$，则所有整数本原提升的本原输出模 $q$ 恰为
\begin{equation}\label{eq:content}
\{g^{-1}Ax:g\mid d,\ g>0\}.
\end{equation}
若 $\Gamma(q,\Z)$ 已实现，每个分支都能从当前列实现。
\end{lemma}
\begin{proof}
伴随恒等式说明 $\cont(Av)\mid d$。写 Smith 分解 $A=U\diag(1,d)V$；对本原 $Vv=(s,t)$，内容为 $\gcd(s,d)$。给定 $g\mid d$，在每个 $\ell^e\Vert d$ 指定 $s\equiv\ell^{v_\ell(g)}\pmod{\ell^e}$（指数为 $e$ 时取 $0$）、$t\equiv1\pmod{\ell^e}$。与模 $q$ 数据用 CRT 合并。所得本原剩余列可提升为本原整数列：固定非零第二坐标后，在第一坐标的等差类中避开第二坐标所含的其余有限素因子。其内容恰为 $g$。最后用引理\ref{lem:fiber}实现准备。
\end{proof}
该引理为什么要先有 $\Gamma(q,\Z)$？因为我们需要先把当前实际列送到具有指定内容的那个提升，再执行 $A$。没有这个先行准备，就不能把“某个提升存在”当成“当前输入能走过去”。

\section{有限集合上的正向操作为什么可以组成群}
假设一个置换在六步后回到原样，那么它的逆就是再做五步。这是有限群中最朴素的事实。在本题中，每条安全内容分支是有限标记集上的置换，所以无需额外要求它的有理矩阵具有统一逆词，仍能在这个有限集合里反向移动。

更妙的是，先走一次内容为 $g$ 的分支，再走若干次内容一分支，将矩阵本身绕回去，剩下的作用就是乘 $g^{-1}$。因此行列式的素因子可以给有限尺度变化。
\begin{corollary}[有限正向对应的群化]\label{cor:finite-group}
设若干返回 $A_i$ 满足引理\ref{lem:content}，并在模 $q$ 保持一个对单位标量封闭的合法列集。它们的正向内容对应轨道等于
\[
\langle A_i\bmod q,\ \ell I:\ell\mid\prod_i\det A_i\text{ 为素数}\rangle
\]
的群轨道。
\end{corollary}
\begin{proof}
每个分支 $g^{-1}A_i$ 是有限集上的置换。内容 $1$ 分支给 $A_i$，其逆为某个非负幂。若 $A_i$ 的阶为 $m$，一次内容 $g$ 分支后接 $m-1$ 次内容 $1$ 分支，给标量 $g^{-1}I$。这些恰生成所示群。
\end{proof}

\section{从模三走到模九：为什么有时会卡住}
普通平移 $x\mapsto x+1$ 在模 $\ell^b$ 下当然走遍所有类。一般分式变换模 $\ell$ 看起来像平移，却可能在更高精度过早循环。因此要额外检查每走 $\ell$ 步，步长是不是恰好多出一个 $\ell$ 因子。
\begin{definition}[局部整数]
$\Z_{(\ell)}$ 是分母不被 $\ell$ 整除的有理数环。
它的单位是赋值为零的元素，可同时约化到所有 $\Z/\ell^b\Z$。
当列的第二坐标是单位时，用其比值 $x=z/y$ 表示仿射方向。
矩阵 $\left(\begin{smallmatrix}a&b\\c&d\end{smallmatrix}\right)$ 对方向的作用为
$x\mapsto(ax+b)/(cx+d)$，只要分母为单位。
\end{definition}
若 $N^2=cI$，所有高次幂都能写成 $A I+B N$，因为偶次幂给标量，奇次幂给 $N$。于是很多矩阵计算只剩两项。例如
\[(I+N)^3=(1+3c)I+(3+c)N.
\]
当 $c=0\pmod3$，$3+c$ 通常恰含一个三；但 $c=6\pmod9$ 时它含至少两个三，就可能跳过一层。这解释了后面为什么会特意避开一个模九余类。
\begin{lemma}[局部平移的逐层提升]\label{lem:local-cycle}
设 $T$ 在 $\Z_{(\ell)}$ 上整可逆，模 $\ell$ 的射影作用为非平凡平移。写 $T=a(I+N)$、$N^2=cI$。若 $\ell\ge5$，或 $\ell=3$ 且 $c\not\equiv6\pmod9$，则其在每个 $\Z/\ell^b\Z$ 的仿射方向上为单循环。对 $\ell=2$，同一结论在 $a$ 为奇单位、$c$ 为偶数时成立。
\end{lemma}
\begin{proof}
此时 $c\in\ell\Z_{(\ell)}$。展开 $(I+N)^\ell=A I+B N$；$A$ 为单位，$v_\ell(B)=1$。在 $\ell=3$ 时 $B=3+c$，在 $\ell=2$ 时 $A=1+c,B=2$。反复取 $\ell$ 次幂后，规范化的 $N$ 系数赋值每次增加一。一般幂的系数赋值因而等于指数的 $\ell$ 进赋值。其在任一方向 $x$ 上的位移还乘一个模 $\ell$ 等于原平移步长的单位。因此固定 $x\bmod\ell^b$ 的指数恰为 $\ell^b$ 的倍数，轨道长度为 $\ell^b$。
\end{proof}

\section{我们到底要把参数送到哪里}
回到 $K_{p,r}$，令 $y=a+pz$，合法性为 $\gcd(y,n)=1$。想选 $j+1$ 份第一块、$j$ 份第二块，以及 $p-1-2j$ 份载体。将它们的和算出来：
\[
(j+1)a+j(rz-a)-(p-1-2j)pz=y-(p^2-nj)z.
\]
所以只要把参数方向送到
\begin{equation}\label{eq:terminal}
t_j=\binom{1}{p^2-nj},\qquad
\max(0,\lceil(p-1-r)/2\rceil)\le j\le(p-1)/2,
\end{equation}
这组数的和就是零。$j$ 的区间只是要求载体取用数在 $0,r$ 之间。

举例：$p=5,n=11,r=1$，区间只允许 $j=2$，终端列是 $(z,y)=(1,3)$。于是 $a=y-pz=-2$，第二块值为 $3$，载体为 $-5$。三份 $-2$ 加两份 $3$ 恰好为零。这是一个非零的配置，但已经能接第5讲的收尾。

\section{为什么只会生成少量小数字，也够用了}
互补混合中，份数最多是 $p$。因此容易得到的是不超过 $p$ 的安全数字，而不是模 $n$ 的每个单位直接给一个操作。下一个引理说，在 $2p<n<4p$ 这段范围，所有小安全数字已经能通过相乘和取逆生成模 $n$ 的全部单位。

证明不是找原根。奇数 $n$ 时，假设有一个最小缺失单位，通过乘二或四再折回较小正代表，会得到更小的缺失单位，矛盾。偶数时则按是否被四整除，利用平衡代表或数量超过群的一半来完成。
\begin{lemma}[小数字生成模n单位群]\label{lem:small-units}
设 $p\ge5$、$2p<n<4p$、$\gcd(p,n)=1$。$-1$ 和所有 $1\le b\le p$、$\gcd(b,n)=1$ 的类生成 $(\Z/n\Z)^\times$。若 $n$ 偶，还存在安全奇数 $b\le p$ 满足 $b\equiv3$ 或 $5\pmod8$。
\end{lemma}
\begin{proof}
记第一项的生成群为 $H$。$n$ 奇时，若 $H$ 不满，取最小正缺失单位 $p<a<n/2$。$a<n/3$ 时取 $4a-n$，$a>n/3$ 时取 $n-2a$，均为更小的正缺失单位，矛盾。
若 $n=2m$、$m$ 奇，选奇数 $b=(m\pm1)/2\le p$，则 $H$ 模 $m$ 包含 $2$ 及其逆。每个模 $m$ 单位的平衡代表绝对值小于 $p$，除尽 $2$ 因子后即是已知奇单位。
若 $4\mid n$，单位在 $(0,n/2)$ 中由 $a\mapsto n/2-a$ 配对，所以 $(0,n/4)$ 恰有 $\varphi(n)/4$ 个单位。它们及负数已属于 $H$，再加 $\pm p$，得到 $|H|>\varphi(n)/2$。

第二项先检查 $3,5,11,13$。$p=5,7$ 时若 $3,5$ 都不安全，则 $30\mid n<4p$，矛盾；$p=11$ 时同理由 $165\mid n$ 排除。若 $p\ge13$ 且四者均不安全，记 $n$ 的不同奇素因子数为 $d\ge4$，$\rho=\prod_{\ell\mid n,\ell\text{奇}}(1-1/\ell)$。等差列 $3+8t\le p$ 的项数 $L\ge(p-4)/8$，容斥给安全项数至少 $L\rho-2^d$。而
\[
L\rho>n\rho/32-1/2\ge60\cdot6^{d-4}-1/2>2^d.
\]
这里 $n\rho\ge1920\cdot6^{d-4}$，由四个指定素因子及其余奇素因子得到。
\end{proof}

\section{把所有工具接成一条有限路径}
下面定理列出三个输入条件，你可以把它看成后面构造的任务单：已经能在同余类内部精确移动；能在模 $n$ 下改变第一坐标；还要能处理有限共同尺度和二进信息。
\begin{theorem}[三块核心的共同终端定理]\label{thm:terminal}
在\eqref{eq:core}上，令 $h=\delta n$、$q=h^2$，其中 $\delta=1$（$n$ 奇）或 $\delta\in\{2,4\}$（$n$ 偶）。假设：
\begin{enumerate}
\item $L(h\Z)$ 和 $\Gamma(q,\Z)$ 可实现；
\item 若干条目本原的整数满秩返回，其行列式与 $q$ 互素，模 $q$ 保持合法域，且有限内容群包含 $V\equiv U(-p^{-2})\pmod n$；
\item $n$ 奇时该群包含标量 $2I$；$n$ 偶时\eqref{eq:terminal}的 $j$ 区间含 $\delta$ 个连续整数。
\end{enumerate}
则全部合法核心可归零。$p=3,n=7$ 时同样成立。
\end{theorem}
\begin{proof}

\par\noindent\textbf{第一步：获得模 $n$ 的全部有限标量。}
保留载体，两个 $p$ 块互补混合给 $(a,z)$ 坐标中的
\[
p^{-1}P_b,\qquad P_b=\begin{pmatrix}b&r(p-b)/2\\0&p\end{pmatrix},
\quad b\text{ 奇},\quad1\le b\le p.
\]
当 $\gcd(b,n)=1$ 时，$P_b$ 的行列式与 $q$ 互素，且 $y'\equiv by\pmod n$。$b<p$ 时条目本原，故推论\ref{cor:finite-group}提供 $p,b$ 的素因子标量；$P_1$ 已提供 $p$。结合假设中的 $2I$ 及引理\ref{lem:small-units}，有限标量模 $n$ 满射。

\par\noindent\textbf{第二步：先把列送对模 $n$。}
因为 $y$ 为单位，$z/y$ 有定义。
矩阵 $V\equiv U(-p^{-2})$ 把比值减去 $p^{-2}$，故可以一次指定所需平移次数。
对当前列，取 $k\equiv p^2zy^{-1}-1\pmod n$，用 $V^k$ 后再乘一个可用有限标量，得到 $(z,y)\equiv(1,p^2)\pmod n$。


\par\noindent\textbf{第三步：处理额外的二进信息。}
$n$ 奇时，$h=n$，本步骤无需操作，直接进入下一步。
$n$ 偶时写 $n=2^a m$、$m$ 奇，令 $B=2^{2v_2(h)}$。有限群中的 $P_bP_1^{-1}$ 及标量给
\begin{equation}\label{eq:split}
D_b=\begin{pmatrix}b^{-1}&0\\n(b^{-1}-1)/2&1\end{pmatrix}.
\end{equation}
引理\ref{lem:small-units}选出 $b\equiv3,5\pmod8$，因 $v_2(b^2-1)=3$，$-1,b$ 生成每个有限二进单位群。用\eqref{eq:split}及 $D_{-1}$ 使 $z=1\pmod B$，同时保持 $y\pmod n$。$V\bmod2$ 为上幺幂，逐次平方给 $V^B\equiv I\pmod B$；使用 $V$ 的 $B$ 倍次数恢复模 $m$ 的首坐标。于是 $z=1\pmod h$。

\par\noindent\textbf{第四步：从浅层修正到平方层。}
按 $y$ 选 $j\bmod\delta$，使当前列与 $t_j$ 模 $h$ 相同。置 $E=E_{21}$、$X=VEV^{-1}$，则模 $n$
\[
Et_j=(0,1)^T,\qquad Xt_j=(-2/p^2,2)^T.
\]

两列的行列式 $2/p^2$ 是模 $m$ 单位。
为什么这里仅需求解\emph{线性}同余？因为 $q=h^2$，对于任意矩阵 $A,B$，
\[
(I+hA)(I+hB)\equiv I+h(A+B)\pmod{h^2}.
\]
二次交叉项消失。两个独立切向列于是能修正任意两坐标的一阶误差。
在奇数部分，$2$ 可逆，所以这两个修正方向足够。用 $L(hk)$ 及其有限 $V$ 共轭，解一次二维线性同余，将两坐标修正到 $t_j\pmod{m^2}$；参数用 CRT 保持模 $B$。随后一次 $L(hk)$ 修正第二坐标模 $B$，并取 $k\equiv0\pmod m$。最终列与 $t_j$ 模 $q$ 相同。

\par\noindent\textbf{第五步：由同余相同到实际相同。}
每条有限内容分支由引理\ref{lem:content}提升，最后用引理\ref{lem:fiber}精确送到 $t_j$。平均其零和组，再用引理\ref{lem:zero-tail}收尾。$p=3,n=7$ 时，模 $7$ 的单位由 $-1,2$ 生成，其余步骤相同，收尾用引理\ref{lem:zero-base}。
\end{proof}

\section{再慢看一次：模n正确以后，误差怎样被消掉}
最容易被一句“线性修正”带过的是第四步。假设当前列 $w$ 与目标 $t$ 已模 $h$ 相同，写
\[w=t+he.
\]
施加 $I+hX$，模 $h^2$ 下得到
\[(I+hX)w\equiv t+h(e+Xt)\pmod{h^2}.
\]
于是修正整列，只需让 $Xt\equiv-e\pmod h$。如果我们手里的两个矩阵方向作用于 $t$ 后成为两条独立列，这就是一个二元线性方程组。

本题中的两条列为 $(0,1)^T$ 与 $(-2/p^2,2)^T$，行列式是 $2/p^2$。对奇素数部分可逆，所以够用；对二进部分不够，因此才先用乘法调整第一坐标，再用一次下剪切调整第二坐标。这就是定理分开处理奇偶部分的真正原因。

最后列已经与 $t_j$ 模 $h^2$ 完全相同，上一讲的精确纤维引理才出场。这时无限数值大小被处理掉，接下来平均零和组并收尾，全程只有有限次操作。

\chapter{先手算十一元五平均，再看一般奇数余量}\label{ch:odd}
\section{现在要做的，是把任务单上的资源造出来}
前面已经知道足够的剪切和同余子群能完成核心。现在不能继续把它们当作假设，必须给出真实位置选择。本讲取 $n=2p+r$，$r$ 为奇数且 $1\le r<p$，包括最紧的 $r=1$。

为方便交换两等大块，用对称坐标
\[u=z,qquad v=a-rz/2,qquad
A=ru/2+v,quad B=ru/2-v,quad C=-pu.
\]
交换 $A,B$ 恰为 $S=\diag(1,-1)$。这就是引入二分之一的目的：把一个原来稍复杂的角色交换变成直接给第二坐标变号。

\section{先看p等于五、r等于一}
从五份 $A$、五份 $B$、一份 $C$ 出发。
\begin{enumerate}
\item 先拿两份 $A$、三份 $B$ 平均，得到五份新值 $D$。
\item 保留一份 $D$ 作最终单点。剩下 $(A,B,C,D)$ 四类库存为 $(3,2,1,4)$。
\item 取输出组 $(2,1,0,2)$，即两份 $A$、一份 $B$、两份 $D$，做五平均。再把余下五个位置平均。
\end{enumerate}
做完又是两个五块加一个单点。这是一条三步返回。

另一条返回先拿四份 $A$、一份 $B$，保留一份新值后库存为 $(1,4,1,4)$，输出取 $(1,2,0,2)$，再平均补集。注意，两套表是\emph{可选的两条完整词}，不是要求同一轮同时取用两套库存。

在 $(u,v)$ 坐标里，两张清去共同分母的矩阵为
\[
J_{-1}=\begin{pmatrix}-5/2&1\\55/4&5/2\end{pmatrix},qquad
J_3=\begin{pmatrix}-5/2&-3\\55/4&5/2\end{pmatrix}.
\]
实际尺度为 $1/25$，两者平方为 $20I,-35I$。因此两条返回都取得了正向射影逆。两者的商为三角矩阵，倍率为 $-7/4$；与 $S$ 作交换子恰给 $U(1)$。倍率及其逆生成 $R=\Z[1/14]$，因而得到 $U(R)$。

这是一段完整的“计数变代数”：三个位置动作，得到一张矩阵；矩阵的平方给逆；两返回比较给倍率；再比较换块前后得到纯剪切。

\section{为什么偏要选择迹零矩阵}
二阶矩阵满足特征多项式恒等式
\[J^2-(\operatorname{tr}J)J+\det(J)I=0.
\]
因此迹为零时，$J^2=-\det(J)I$。只要行列式非零，重复同一返回就给射影逆。这比先搜索一个很长周期容易得多。

一般公式的第一项工作，就是把输出计数安排得使迹为零。第二项工作，是让两张反射只差一个简单方向。最后再证明所有取用数真的在库存内。你可以带着刚算过的 $(5,1)$ 例子，读下面每个参数。
\begin{proposition}[奇数余量的全部合法核心]\label{prop:odd}
若 $p\ge3$ 为奇素数、$1\le r<p$ 为奇数，则全部合法 $K_{p,r}$ 可归零。
\end{proposition}
\begin{proof}

\par\noindent\textbf{第一步：选择两组库存内的反射。}
采用 $u=z,v=a-rz/2$。写 $p-r=4k+2\epsilon$，$\epsilon\in\{0,1\}$。取 $\alpha_0=-r,\alpha_1=4-r$。第一步从 $A,B$ 各取 $(p+\alpha)/2,(p-\alpha)/2$ 份，产生 $p$ 份 $D$，保留 $r$ 份作载体；第二步按
\[
(x,y,c,v_D)=\left(\frac{p+2r-\alpha-2\epsilon}{4},\frac{p+\alpha-2\epsilon}{4},\epsilon,\frac{p-r}{2}\right)
\]
选择输出，第三步平均其 $p$ 元补集。两组输出分别为 $(r+k,k,\epsilon,2k+\epsilon)$ 和 $(r+k-1,k+1,\epsilon,2k+\epsilon)$；剩余库存分别为 $(r+2k+\epsilon,2k+\epsilon,r,4k+2\epsilon)$ 和 $(r+2k+\epsilon-2,2k+\epsilon+2,r,4k+2\epsilon)$。唯一紧的条件是 $k+\epsilon\ge1$，由 $p-r>0$ 保证。

这三次操作的顺序是：先产生一个新 $p$ 块，留出其中的 $r$ 份作载体，
再平均一个输出组，最后平均其补组。补组自动有 $p$ 个位置，
因此只需逐项检查所列输出组不超过剩余库存。
返回的实际尺度为 $p^{-2}$，矩阵为
\[
J_\alpha=\begin{pmatrix}-rp/2&-\alpha\\pn(r-2\epsilon)/4&rp/2\end{pmatrix},
\quad J_\alpha^2=pF_\alpha I,\quad F_\alpha=\frac{pr^2-n(r-2\epsilon)\alpha}{4}.
\]

\par\noindent\textbf{第二步：标量周期与参数环。}
平方标量模 $n$ 为 $p^4$，故非零且安全。令 $c_\epsilon=r$（$\epsilon=0$）或 $1$（$\epsilon=1$），$f_i=F_{\alpha_i}/c_\epsilon$。若 $\epsilon=0$，则 $f_0=r(3p+r)/4,f_1=f_0-n$；若 $\epsilon=1$，则 $f_0=r[pr+n(r-2)]/4,f_1=f_0-n(r-2)$。这给出
\[
\gcd(f_0,f_1)=\gcd(f_0f_1,n)=1,\qquad r\mid f_0,
\]
且二者奇偶相反。第二种情形的互素性由模 $n$、模 $r-2$ 的等式 $4f_0\equiv pr^2$ 得出。

相邻标量互素，故三角商的倍率 $f_1/f_0$ 约分后仍保留两者的全部素因子。
而 $f_0,f_1$ 奇偶相反、$r\mid f_0$，所以 $2$ 与 $r$ 都会成为参数环的单位。
在 $R=\Z[1/|f_0f_1|]$ 中，$2,r$ 为单位，故 $R$ 的单位无限。令 $S=\diag(1,-1)$，则
\begin{gather*}
T=J_{\alpha_1}J_{\alpha_0}^{-1}
=\begin{pmatrix}f_1/f_0&-2r/(c_\epsilon f_0)\\0&1\end{pmatrix},\\
[S,T]=U\left(\frac{4r}{c_\epsilon f_0}\right).
\end{gather*}
根参数为 $R$ 单位，引理\ref{lem:localize}给 $U(R)$。

\par\noindent\textbf{第三步：取得横向剪切。}
保留 $r$ 份旧 $B$，平均两组
\[
(A^{(p+r)/2},B^{(p-r)/2}),\qquad
(A^{(p-r)/2},B^{(p-r)/2},C^r),
\]
得到
\[
K=\begin{pmatrix}-r/2&1\\rn/4&r/2\end{pmatrix},\quad K^2=r(p+r)I/2.
\]
若 $H=KU(2/r)$，则 $[S,H]=L(n)$。置 $P=\left(\begin{smallmatrix}1&-2\\0&n\end{smallmatrix}\right)$，有
\[
P^{-1}U(t)P=U(nt),\quad
P^{-1}KU(8(p+r)t/r)K^{-1}P=L(nt).
\]

这里换基矩阵的行列式为 $n$，必须计入一层分母损失。
引理\ref{lem:deep}先给 $\Gamma(n^3,R)$。对 $t\in R$ 选整数 $k\equiv t\pmod{n^2R}$，分解 $L(nt)=L(nk)L(n(t-k))$，得到 $L(nR)$，再得 $\Gamma(n^2,R)$。

\par\noindent\textbf{第四步：接入共同终端。}
换到 $(z,y)$ 的矩阵 $C=\left(\begin{smallmatrix}1&0\\n/2&1\end{smallmatrix}\right)$。选 $4\mid k$ 且 $k\equiv-p^{-2}\pmod n$，则 $CU(k)C^{-1}$ 整数且满足定理\ref{thm:terminal}的方向条件。至少一个整数反射的行列式为偶数，而条目公因子为奇数（其左下项 $-\alpha_i$ 奇）；引理\ref{lem:content}提供有限标量 $2I$。定理\ref{thm:terminal}完成证明。
\end{proof}

\section{为什么先绕到n的三次方，再回到平方}
证明第三步有个看似绕远的地方：先得到 $\Gamma(n^3,R)$，才取得全部 $L(nR)$，最后得到 $\Gamma(n^2,R)$。

我们起初知道的只有整数步长 $L(nk)$，还不能直接把 $k$ 换成所有环元素 $t$。但每个 $t\in R$ 都能选整数 $k$ 使 $t-k\in n^2R$。所以
\[L(nt)=L(nk)L(n(t-k)).
\]
第一项已经能做，第二项属于已经取得的深子群。这样就补齐所有环参数，再应用两根定理回到平方层。每一次群的使用都有更早的真实操作作起点。

终端已经在上一讲算过：十一元五平均送到 $(-2^5,3^5,-5)$ 的方向后，三份 $-2$ 加两份 $3$ 平均成零，剩下交给含零收尾。大段矩阵最终服务的仍是这一个真实选择。

\chapter{只剩两个载体位置时，直接指定最后的余数}\label{ch:two}
\section{不要为所有输入预先选择同一条路}
现在 $n=2p+2$。核心写成
\[((u+v)^p,(u-v)^p,(-pu)^2).
\]
令 $m=p+1$，合法本原参数要求 $\gcd(u,v)=1$ 和 $\gcd(v,m)=1$，所以 $v$ 是奇数。

这一讲采用更有针对性的办法。先取得一些剪切和一个足够深的主同余子群；再根据当前输入选两次返回，直接把第二坐标设为 $1$ 或 $-1$ 模控制层。这样不必先证明所有有限尺度都可达。
\begin{proposition}[余量为2的核心]\label{prop:r2}
若 $p\ge5$，全部合法 $K_{p,2}$ 可归零。
\end{proposition}
\begin{proof}
引理\ref{lem:two-carrier}给出 $U(\Z[1/p])$ 和一个明确整数深核。其两次 CRT 数字返回使第二坐标等于 $\pm1$ 模该深核层，再由引理\ref{lem:fiber}精确送入相反值终端。
\end{proof}

\section{两次CRT分别负责什么}
已有上剪切就能把 $u$ 改成 $u+kv$。由于 $v$ 奇，可以把它写成 $2a-v$。某个数字返回给
\[(u,v)=(2a-v,v)quad\longmapsto\quad(-a,pv-ma)
\]
的方向。现在 $a$ 可以在一个等差数列中自由选，我们用它做两件事。

第一次，要让 $pv-ma$ 在控制模数 $q$ 的每个素因子处都非零。每个素数只提出一个条件，用 CRT 合并。不是让这个数变小，而是让它成为模 $q$ 的单位。

第二次，既然它已经可逆，就能选一个允许的数字 $b$，再选新的平移参数，让
\[bv_1-ma_1\equiv\epsilon\pmod q,qquad\epsilon=\pm1.
\]
与此同时还要让输出坐标互素；否则最后约分可能把刚安排的余数改掉。证明中的额外CRT条件正是在保护本原性。

\section{一组可以手算的CRT数据}
取 $p=5,m=6$，该构造可使用 $q=(p^2-1)^2=576$。从合法参数 $(u,v)=(1,5)$ 开始，先选 $a=8$，满足 $2a\equiv1\pmod5$。上平移把第一坐标变成 $2a-v=11$，再用第一个数字返回，得到本原分析列
\[(-8,25-48)=(-8,-23).
\]
$23$ 与 $576$ 互素，第一次任务完成。

$-23\equiv1\pmod6$。为避免不可用的数字 $b=1$，取 $b=5,\epsilon=-1$。第二个平移参数要满足
\[a_1\equiv-4\pmod{23},qquad a_1\equiv-19\pmod{96}.
\]
可取 $a_1=1997$，它还与五互素。第二次输出为
\[(-1997,5(-23)-6\cdot1997)=(-1997,-12097).
\]
两坐标最大公因数是一，且 $-12097\equiv-1\pmod{576}$。再用上剪切 $U(307)$，第一坐标成为 $-3715776=-6451\cdot576$。这就准确送进了 $(0,-1)$ 的同余类，精确输送引理接手。

数值在这个分析代表里变得很大，但实际平均词带着共同收缩因子。这个例子强调的是\textbf{有一个有限的解法}；它没有声称这条解法短，也没有声称每次本原参数都变小。

\section{把这些数字变成真实位置选择}
下面先给上平移和缩放逆的平均词，再给横向反射，最后证明两次CRT可对所有合法输入执行。刚才的例子对应证明的第三、第四阶段。
\begin{lemma}[两次CRT将尺度送到正负1]\label{lem:two-carrier}
在 $((u+v)^p,(u-v)^p,(-pu)^2)$ 上，合法本原参数满足 $\gcd(u,v)=\gcd(v,p+1)=1$。全部合法参数可送到 $(u,v)=(0,1)$ 的射影方向。
\end{lemma}
\begin{proof}

\par\noindent\textbf{第一阶段：上平移与缩放的正向逆。}
记 $m=p+1$。保留两份 $A=u+v$，平均
\[(A^{(p-3)/2},B^{(p+1)/2},C),\qquad
(A^{(p-1)/2},B^{(p-1)/2},C)
\]
实现 $-U(1)/p$；换块给 $U(-1)$。互补平衡混合给 $F=\diag(1,1/p)$。其逆由四次平均实现：先平均两组 $(A^{(p-1)/2},B^{(p-1)/2},C)$，产生 $2p$ 份 $D=-u/p$；保留两份 $D$，分别平均一份剩余 $A$ 或 $B$ 与 $p-1$ 份 $D$。输出为 $(u/p^2,v/p)$。因此 $U(\Z[1/p])$ 可实现。

\par\noindent\textbf{第二阶段：三步反射与深同余群。}
再平均 $(A^{(p+1)/2},B^{(p-3)/2},C)$。保留两份新值，按 $(x,x+1,b,(p-3)/2)$ 选下一组，其中 $b=1$（$p\equiv1\pmod4$）或 $0$（$p\equiv3\pmod4$），$x=(p-b-(p-1)/2)/2$，再平均补组。返回为 $J/p^2$，
\[
J=\begin{pmatrix}1&-2\\c&-1\end{pmatrix},\qquad
c=\begin{cases}-(p^2-1)/2,&p\equiv1\pmod4,\\(p+1)^2/2,&p\equiv3\pmod4.
\end{cases}
\]
两类计数分别是 $(k,k+1,1,2k-1)$ 和 $(k+1,k+2,0,2k)$，均在库存内。$J^2=(1-2c)I$。若 $S=\diag(1,-1)$、$H=JU(2)$，则 $[S,H]=L(-2c)$。借 $F$ 共轭得到完整 $L(2cR)$，$R=\Z[1/p]$。令 $q=(2c)^2$，引理\ref{lem:deep}给 $\Gamma(q,R)$；$m\mid2c$ 且 $\gcd(p,q)=1$。


\par\noindent\textbf{第三阶段：使第二坐标在整个控制层可逆。}
对 $2\le b\le p$，保留两份 $A$，平均 $(A^{b-2},B^{p-b},C^2)$ 和 $(A^{p-b},B^b)$，实现
\[
p^{-1}M_b,\qquad M_b=\begin{pmatrix}-1&-1\\-m&2b-m\end{pmatrix}.
\]

这个 $M_b$ 是可用的正向返回，但不预先假定它对所有本原输入安全或可逆。
下面每次使用它，都会先安排输入，使本原输出保持合法。
由于 $v$ 奇，先用整数上平移把 $u$ 改成 $2a-v$，即选 $2a\equiv u\pmod v$。
直接乘矩阵可见
\[
\frac1pM_p(2a-v,v)^T=\frac2p(-a,pv-ma)^T.
\]
其中 $2/p$ 是全部核心的共同尺度，所余整数列将被安排为本原。
执行 $M_p$ 后本原分析列为 $(-a,pv-ma)$。对 $\ell\mid q$，若 $\ell\mid m$，第二项自动非零；若 $\ell\mid v$，$a$ 自动非零；其余素数指定 $a=(pv-1)/m\pmod\ell$。
这些额外同余只在不整除 $v$ 的素数处新增，所以与 $2a\equiv u\pmod v$ 相容。
再令 $p\nmid a$；若 $p\mid v$，原来的同余已经保证这一点，否则另加 $a\equiv1\pmod p$。
由于 $\gcd(a,v)=1$，有
\[
\gcd(a,pv-ma)=\gcd(a,p)=1.
\]
这给出本原性，而对所有 $\ell\mid q$ 的选择给出第二坐标的单位性。
CRT 给出本原输出 $(u_1,v_1)$ 且 $\gcd(v_1,q)=1$。

\par\noindent\textbf{第四阶段：直接指定第二坐标的完整余数。}
取 $v_1^{-1}\pmod m$ 的代表 $b_0\in[1,p]$；若 $b_0=1$，改取 $b=p,\epsilon=-1$，否则取 $b=b_0,\epsilon=1$。于是 $b$ 奇、$2\le b\le p$、$bv_1\equiv\epsilon\pmod m$。选
\[
2a_1\equiv u_1\pmod{v_1},\qquad
a_1\equiv(bv_1-\epsilon)/m\pmod{q/m},\qquad \gcd(a_1,b)=1.
\]

除法 $(bv_1-\epsilon)/m$ 是整数除法，因为先前已选 $bv_1\equiv\epsilon\pmod m$。
此时第二个同余恰好等价于 $bv_1-ma_1\equiv\epsilon\pmod q$。
前两个模数互素；$b$ 的素因子若整除其中一个模数，相应条件已使 $a_1$ 非零，其余处由 CRT 补齐。平移后执行 $M_b$，得到本原列 $(-a_1,bv_1-ma_1)$，第二项为 $\epsilon\pmod q$。
因为 $\gcd(a_1,v_1)=\gcd(a_1,b)=1$，有
\[
\gcd(a_1,bv_1-ma_1)=\gcd(a_1,bv_1)=1.
\]
所以这里无须再约去任何内容，已经安排好的第二坐标余数确实保留。
再用上平移使第一项为零模 $q$，由引理\ref{lem:fiber}精确到 $(0,\epsilon)$。相反值收尾完成。
\end{proof}

\chapter{偶数余量为什么需要多照顾一个素数}\label{app:lower}
\section{先看清楚奇偶分支的来源}
写 $r=2t,m=p+t$，把两块对称地写成
\[A=tu+v,qquad B=tu-v,qquad C=-pu.
\]
两个大块的差是 $2v$，第一块与载体的差是 $mu+v$。所以本原配置的 $G$ 是
\[G=\gcd(2v,mu+v).
\]
若 $t$ 奇，则 $m$ 偶，$\gcd(v,m)=1$ 已保证 $v$ 奇，也保证所有危险差异。若 $t$ 偶，则 $m$ 奇，除了 $\gcd(v,m)=1$，还要 $u+v$ 奇。额外的模二条件就这样出现了。

因此这里的分类不是因为喜欢分奇偶，而是同一个差分公式在素数二处真的有不同表现。后面“二进深度”只是在进一步问：偶的那个坐标到底被多少个二整除？
\begin{proposition}[偶数余量的其余核心]\label{prop:even}
若 $p\ge5$、$2<r<p$ 为偶数，则全部合法 $K_{p,r}$ 可归零。
\end{proposition}
\begin{proof}
$4\le r\le p-3$ 使用引理\ref{lem:even-core}；$r=p-1$ 使用引理\ref{lem:upper-boundary}。两引理的根、局部化环和终端均在第\ref{app:lower}讲给出。
\end{proof}

\section{用一个严格下降的小步骤摆脱高二进深度}
在 $(a,z)$ 坐标里，合法性使两者奇偶相反。若偶数坐标被 $2^\nu$ 整除、$\nu\ge2$，希望一次完整返回后，两个输出共同含一个二，但不共同含四。约去这一个二，本原参数的深度就下降一。

这是很熟悉的思路：如果第一输出为“被四整除的一项”加“恰被二整除的一项”，那么它恰被二整除。我们只需要选择返回矩阵的某个系数也恰含一个二。下面的配比就是为这个小目标设计的。
\begin{lemma}[降低二进深度]\label{lem:binary-depth}
设 $r<p$ 为正偶数。在合法 $K_{p,r}(a,z)$ 上，若 $az\ne0$，可使本原参数满足 $v_2(az)=1$。
\end{lemma}
\begin{proof}
合法性使 $a,z$ 奇偶相反。若深度 $\nu=v_2(az)\ge2$ 且 $a$ 奇，先用
\[
(a,z)\mapsto p^{-1}(dz,-a),\qquad d=[pr-n(r-1)]/2
\]
将偶坐标放在 $a$；该两组平均是下文\eqref{eq:even-reflection}取 $s=r-1$ 的返回，$d$ 奇且与 $n$ 互素。于是 $4\mid a,z$ 奇。
保留 $r$ 份旧 $A$，平均组成 $(p-j-s,j,s)$ 的一组及其补组，给
\[
M_{j,s}=\begin{pmatrix}p-2j-s&rj-ps\\-1&0\end{pmatrix}.
\]
若 $r\equiv2\pmod4$，取 $(j,s)=(p,0)$。若 $4\mid r$，令 $Q_0=\lfloor n/3\rfloor$，在 $Q_0,Q_0-1,Q_0-2$ 中按顺序取与 $n$ 奇部互素的 $b$：$3\nmid n$ 时取 $Q_0$，否则取不被 $3$ 整除的 $Q_0-1$ 或 $Q_0-2$。有 $r-1\le b\le p-1$；令 $(j,s)=(b-1,2)$。后一范围由 $p-r>0$，以及 $3\mid n$ 时 $p-r\ge3$ 得出。
两种选择都使 $D=rj-ps$ 恰含一个 $2$，且与 $n$ 奇部互素；后一种用 $D\equiv-2pb\pmod\ell$。
所选 $D$ 的二进赋值恰为 $1$，而当前 $4\mid a$、$z$ 为奇数。
因此第一坐标是“被 $4$ 整除的数”加“恰被 $2$ 整除的数”，其赋值必为 $1$。
整数输出为 $((p-2j-s)a+Dz,-a)$，两项的二进赋值为 $1,\nu$，故本原化后深度减一。见证坐标的输出为 $njz-(2j+s)(a+pz)$，奇素数处保持单位，模 $2$ 亦保持合法。迭代完成。
\end{proof}

\section{两个相邻反射先造出一个足够深的群}
标准反射由两个不交平均组组成，再留下 $r$ 份旧值作载体。参数 $s$ 表示第一组取几份旧载体；不同的可行 $s$ 给不同反射。

取 $s=r-1$ 和 $s=r-3$，得到的平方标量相差 $n$，而两者互素。比较它们，就取得一个三角倍率；再与换块比较，取得剪切。这和奇数余量的证明做的是同一件事，只是系数不同。

开始得到的根步长含一个额外的 $r$ 因子。设 $t=2^a t_0$，其中 $t_0$ 奇，接下来需要把 $t_0$ 变成环的单位。难点在于它可能还整除当前控制模数，不能直接用“行列式安全”引理。证明因此设计一个夹乘 $C_cJC_c/c$，使它与 $J$ 行列式相同；除以 $J$ 后就是行列式一矩阵，可以取得逆。

\section{读完整证明时，分清三条路}
当 $t$ 奇时，局部化掉奇部以后已经足以接共同终端，路线很短。当 $t$ 偶且有三个标准反射时，要先调整二进方向，再调整奇素数方向，最后匹配共同单位。两个边界只有两个标准反射，则加一个迹为 $\pm2$ 的返回，直接把根步长改善到 $2n$，又能接共同终端。

第二条路最长。读它时可以把“方向”理解为比值 $z/a$，把“尺度”理解为同一方向的哪个本原列。先把比值调对，再把共同倍数调对，两步缺一不可。
\begin{lemma}[偶数内部的完整构造]\label{lem:even-core}
对 $p\ge7$、偶数 $4\le r\le p-3$，全部合法 $K_{p,r}$ 可归零。
\end{lemma}
\begin{proof}

\par\noindent\textbf{第一阶段：两个标准反射给出根。}
写 $r=2t,m=p+t$，用 $a=tu+v,z=u$。合法性为 $\gcd(v,m)=1$；$t$ 偶时还要求 $u+v$ 奇。保留 $r$ 份旧 $A$，对可行奇数 $\max(1,4t-p)\le s\le2t$，平均
\begin{equation}\label{eq:even-reflection}
(A^i,B^i,C^s),\quad(A^{p-r-i},B^{p-i},C^{r-s}),\qquad i=(p-s)/2.
\end{equation}
所得矩阵为 $J_s/p$，
\[
J_s=\begin{pmatrix}-t&-1\\m(t-s)&t\end{pmatrix},\qquad
J_s^2=-d_sI,\quad d_s=pt-ms.
\]
取 $s_0=r-1,s_1=r-3$，对应 $d,e=d+n$ 互素且为奇数，与 $n$ 互素。在 $(a,z)$ 坐标中，同一反射为 $\left(\begin{smallmatrix}0&d_s\\-1&0\end{smallmatrix}\right)$，相邻商为 $\diag(1,e/d)$，与换块 $\Sigma=\left(\begin{smallmatrix}-1&r\\0&1\end{smallmatrix}\right)$ 的交换子为 $U(rn/e)$。因此在 $R_0=\Z[1/|de|]$ 上有双向 $rnR_0$ 根。$|de|>1$，引理\ref{lem:deep}适用。


\par\noindent\textbf{第二阶段：用夹乘反演奇数倍率。}
互补混合给 $F_c=\diag(1,c/p)$，$c$ 为 $|c|\le p$ 的奇数。引理\ref{lem:inverse}先激活 $F_1^{-1}$，得到 $C_c=\diag(1,c)$。若奇数 $c\mid t$，则
\[
H=c^{-1}C_cJ_{s_0}C_c
=\begin{pmatrix}-t/c&-1\\m(t-s_0)&tc\end{pmatrix},\quad
HJ_{s_0}^{-1}\in\SL_2(R_0).
\]

因为 $\det H=\det J_{s_0}=d$ 在 $R_0$ 中已经为单位，
矩阵 $HJ_{s_0}^{-1}$ 的行列式为 $1$，可以使用已有深同余群来实现它的逆。
这一步避免了直接把可能整除旧控制层的 $c$ 当作安全单位。
引理\ref{lem:inverse}使 $H^{-1}$ 可实现，而 $J_{s_0}C_cH^{-1}=cC_c^{-1}$。写 $t=2^a t_0$，由引理\ref{lem:localize}可取 $R=R_0[1/(pt_0)]$ 并得到
\begin{equation}\label{eq:even-deep}
\Gamma((2^{a+2}m)^2,R).
\end{equation}
以后任何与 $n$ 互素的小奇数字均可由引理\ref{lem:inverse}激活；同一根共轭在扩环后保留此层。


\par\noindent\textbf{第三阶段：$t$ 为奇数的情形。}
若 $t$ 奇，已有 $U(2nR)$ 及 $\Gamma(4n^2,R)$。在 $(z,y)$ 基下，标准反射规范化模 $n$ 为 $\left(\begin{smallmatrix}1&-1/p\\0&-1\end{smallmatrix}\right)$，乘换块后取适当整数幂得到定理\ref{thm:terminal}的 $V$。规范化所需的标量 $p$ 由 $P_1$ 的内容分支提供，故该定理适用。


\par\noindent\textbf{第四阶段：有三个反射的偶 $t$。}
设 $t$ 偶且 $8\le r\le p-5$。第三反射 $s_2=r-5$ 亦可用，将 $d+n,d+2n$ 的素因子一并加入 $R$。取奇数 $c_*=(m+\sigma)/2$、$\sigma=\pm1$，则 $c_*\le p$、$\gcd(c_*,n)=1$，而单位 $\epsilon_2=\sigma/c_*$ 模 $m$ 等于 $2$。引理\ref{lem:binary-depth}先令 $v_2(az)=1$；必要时再用标准反射，使 $a$ 奇、$u=z$ 恰含一个 $2$，故 $v$ 奇。取 $v^{-1}\pmod m$ 的平衡代表 $2^h c$，$c$ 奇、$|c|<p$；用 $U(Qk)$ 使第一坐标与 $c$ 互素后施加 $C_c$，便使 $v\bmod m$ 属于实际单位像。这里 $Q=(2^{a+2}m)^2$。

\par\noindent\emph{先控制二进方向。}
比值 $(d+n)/d,(d+2n)/(d+n)$ 模 $8$ 分别为 $3,7$；取其中的 $\kappa\equiv3$。它作用于 $x=z/a$ 为乘 $\kappa$，换块作用为 $x/(-1+rx)\equiv-x\pmod{16}$。
把 $x$ 写成 $2w$、$w$ 奇，则模 $16$ 的 $x$ 对应模 $8$ 的四个奇单位。
乘 $\kappa\equiv3\pmod8$ 和变号生成这四个类，因此覆盖赋值一的全部模 $16$ 方向；$v_2(\kappa^2-1)=3$ 又给全部高层提升。可将二进方向送到
\[
u_T=2,\quad a_T=p-r(p-3)/2,\quad v_T=a_T-2t.
\]

\par\noindent\emph{再控制奇素数方向，同时固定二进方向。}
奇素数处用 $G_s=-\diag(1,-1)J_s$。其 $G_s/t=I+N$ 满足 $N^2=m(t-s)I/t^2$；相邻 $s$ 可避开引理\ref{lem:local-cycle}的模 $9$ 例外。$G_s$ 的二进矩阵阶为 $2$ 幂，故 CRT 选择幂可在保持二进方向时命中所有奇素数方向。

\par\noindent\emph{最后匹配单位尺度。}
此时方向相同，所以当前列与目标列模控制层只差一个单位。
需要证明这个有限单位来自真正的 $R$ 单位。
尺度模 $m$ 的变化由 $R$ 单位给出，其中 $t$ 的类可用 $\epsilon_2^a t_0$ 表示，且 $v_T\equiv t^2\pmod m$。剩余单位核由 $(d+n)/d$ 和 $(d+2n)/d$ 生成：后者在 $2$ 处从深度 $2$ 开始，在各奇处从 $v_\ell(m)$ 开始；前者补模 $4$ 的另一个类。
这里使用的升位事实如下。若 $\ell$ 为奇素数，$u=1+\ell^b c$、$\ell\nmid c$，
则二项式展开给
\[
v_\ell(u^\ell-1)=b+1.
\]
对 $\ell=2$，同样结论在 $b\ge2$ 时成立。
所以这样的单位的连续素数幂逐层产生整个主单位核。
各局部核的阶为对应素数的幂，用中国剩余定理选指数就能同时实现所有分量。
逐次取素数幂并用指数 CRT 完成所有精度。引理\ref{lem:fiber}精确送到该终端。取 $(p+1)/2$ 份 $A$、$(p-3)/2$ 份 $B$ 和一份 $C$ 即得零和组。


\par\noindent\textbf{第五阶段：只有两个标准反射的内部边界。}
剩下 $r=4$，或 $r=p-3$ 且 $p\equiv3\pmod4$。保留旧 $A$ 的返回 $M=\left(\begin{smallmatrix}a_0&D\\-1&0\end{smallmatrix}\right)$ 可按下表选择，其中第一组为 $(p-j-s,j,s)$、第二组为补集：
\[
\begin{array}{c|rrrr}
&s&j&a_0&D\\\hline
r=4,p\ge11&1&(p-3)/2&2&p-6\\
r=p-3,p\equiv1\ (3)&p-4&1&2&-p^2+5p-3\\
r=p-3,p\equiv2\ (3)&p-6&4&-2&-p^2+10p-12
\end{array}
\]
$p=7,r=4$ 用第二行。代入可验容量非负及 $\gcd(D,rn)=1$。由引理\ref{lem:inverse}取得 $M^{-1}$，相邻商局部化 $D$；与换块的交换子给
\[U\bigl(2a_0+r(1-D/d_s)\bigr).
\]

这个新根参数的作用是减少旧根理想中过多的 $2$ 因子。
该参数属于 $nR$，二进赋值恰为 $2$。与原根 $2^{a+1}nR$ 合并得到 $U(2nR)$，标准反射给另一方向，故有 $\Gamma(4n^2,R)$。标准反射仍提供前述有限 $V$，直接应用定理\ref{thm:terminal}。这覆盖全部偶数内部参数。
\end{proof}

\section{乘三、变号，为什么能处理所有二进单位}
证明里有 $\kappa\equiv3\pmod8$。模八的奇数只有 $1,3,5,7$，乘三和变号就能从一到达全部四类。更高位怎么办？

最简单地看 $3^2=9=1+8$。平方一次，$9^2-1=80$ 恰被 $16$ 整除；再平方，赋值继续增加一。一般地，若 $u=1+2^bc$、$b\ge2,c$ 奇，则
\[u^2-1=2^{b+1}c+2^{2b}c^2
\]
恰含 $b+1$ 个二。这样逐位选择幂，就能修正每个更高二进位。所谓“生成主单位核”，在这里就是这件具体的逐位调整。

\section{r接近p时，重新选择实际分组}
当 $r=p-1$，要保留的载体太大，刚才某些标准分组不再有足够位置。于是先平均一组制造新值，再把其中 $p-1$ 份留作载体，之后用第二组及补集恢复两个 $p$ 块。

它仍是相邻反射、倍率、剪切的同一机制，只是必须重新核对取用数。最小的 $p=5,n=14$ 再单列三个具体矩阵；这些矩阵的行列式因子和夹乘等式都列在下面。
\begin{lemma}[余量为p−1的重叠反射]\label{lem:upper-boundary}
若 $p\ge5$、$r=p-1$，全部合法 $K_{p,r}$ 可归零。
\end{lemma}
\begin{proof}
置 $n=3p-1,d=p^2-n$，标准两步反射为 $A=\left(\begin{smallmatrix}0&-d\\-1&0\end{smallmatrix}\right)$。先平均 $(A^{p-\sigma},C^\sigma)$，保留 $r$ 份新值，再平均 $(A^{h-\sigma},B^h,C^{p+\sigma-2h})$ 及其补集，给
\[
J_{\sigma,h}=\begin{pmatrix}-p\sigma&p[nh-p(p+\sigma)]\\-(p-\sigma)&p\sigma\end{pmatrix},
\quad J_{\sigma,h}^2=pf_{\sigma,h}I,\quad f_{\sigma,h}=p^3-n(p-\sigma)h.
\]
$p\ge7$ 时取 $(\sigma,h)=(2,3),(2,4),(3,4),(3,5)$；容量等价于 $0\le h-\sigma\le\sigma$、$0\le p+\sigma-2h\le p-1-\sigma$，均成立。记相邻标量 $f_\sigma,g_\sigma$，它们互素且与 $n$ 互素。相邻商与换块的交换子给 $U(n(n\sigma-pr)/g_\sigma)$。两整数系数的最大公因子为 $\gcd(2n-pr,3n-pr)=2$，故在 $R_0=\Z[1/|f_2g_2f_3g_3|]$ 上有 $U(2nR_0)$。由 $AU(dx)A^{-1}=L(x)$ 得双向根和 $\Gamma(4n^2,R_0)$。
$A/p$ 在 $(z,y)$ 坐标中模 $n$ 为 $\left(\begin{smallmatrix}1&-1/p\\0&-1\end{smallmatrix}\right)$。与换块组合及有限标量 $p^{-1}$ 提供定理\ref{thm:terminal}所需方向，因此完成 $p\ge7$。

\par\noindent\textbf{最小边界 $p=5,n=14$。}
这里直接列出三个已有真实分组的矩阵，先得到含额外因子 $3$ 的根。
随后一次夹乘让 $3$ 成为可逆倍率，根理想便由 $84R$ 化为 $28R$。
对 $p=5,n=14$，取
\[
A=\begin{pmatrix}0&-11\\-1&0\end{pmatrix},\quad
S=\begin{pmatrix}-1&4\\0&1\end{pmatrix},\quad
J=\begin{pmatrix}-5&-10\\-4&5\end{pmatrix}.
\]

这里 $A$ 的两组来源是 $(A,B,C^3)$、$(B^4,C)$，保留四份旧 $A$，实际尺度为 $1/5$。
$S$ 是交换两个五块的名称。
$J$ 取重叠构造的 $(\sigma,h)=(1,2)$：先平均 $(A^4,C)$，
保留四份新值，再平均 $(A,B^2,C^2)$ 及其补组，实际尺度为 $1/25$。
所以这三个矩阵都已由明确原位置选择实现。
$A^2=11I,J^2=65I$，且 $W=ASASJ=\left(\begin{smallmatrix}-121&330\\0&65\end{smallmatrix}\right)$，$[S,W]=U(84/65)$。于是 $R_0=\Z[1/(5\cdot11\cdot13)]$ 上有双向 $84R_0$ 根及 $\Gamma(7056,R_0)$。由 $F_1$ 激活 $C_3=\left(\begin{smallmatrix}3&-4\\0&1\end{smallmatrix}\right)$；
\[
H=C_3JC_3/3=\begin{pmatrix}1&-18\\-4&7\end{pmatrix},\quad
\det H=\det J=-65.
\]
$HJ^{-1}\in\SL_2(R_0)$，故 $H^{-1}$ 可实现，继而 $JC_3H^{-1}=3C_3^{-1}$。局部化 $3$ 后 $84R=28R$，给 $\Gamma(784,R)$。同一标准反射 $A$ 提供有限方向，仍接定理\ref{thm:terminal}。
\end{proof}
这一讲可以看出，一般代数工具已经很统一，长度主要花在“现有几份数够不够实现那个工具”。位置容量不能由群论自动补出来。

\chapter{把一张取数表，逐行翻译成矩阵}\label{app:upper}
\section{上半带的大载体从哪里来}
设 $n=3p+s$，$r=p+s$。核心是 $(A^p,B^p,C^r)$。现在载体有 $r>p$ 份，如果想全部换成一个新均值，做一次平均只会产生 $p$ 份，显然不够。因此先做两组\emph{组成相同、位置不交}的平均，共产生 $2p$ 份新值 $D$。

每组从 $A,B,C$ 取 $i,j,k$ 份，$i+j+k=p$。新值为
\[D=\frac{iA+jB+kC}{p}
=\frac{(i-j)a+(rj-pk)z}{p}.
\]
保留 $r$ 份 $D$ 作新载体，其余恰有 $2p$ 个位置，分成两组输出即可。这件事的全部自由度，就是来源三项计数，以及第一输出组的四项计数。

\section{先把库存列出来，再做任何代数}
两组来源用去 $2i,2j,2k$ 份旧值。保留 $r$ 份 $D$ 后，剩余四类库存是
\[(p-2i, p-2j, r-2k, 2p-r).
\]
第一输出取 $(x,y,u,v)$。它必须逐项不超过库存，且四项和为 $p$。第二输出自动是剩下的 $p$ 项。此处 $y$ 是一项计数，暂时不是先前的见证坐标 $a+pz$；读局部公式时要按本段定义使用。

新载体满足 $-pz'=D$，所以矩阵第二行直接得到。第一输出的均值是 $a'$，所以把 $(xA+yB+uC+vD)/p$ 展开，就得到第一行。
\begin{lemma}[由库存表计算返回矩阵]\label{lem:count-interface}
在 $K_{p,r}$ 上，选择 $q_0\in\{1,2\}$ 组不交且组成相同的来源 $(i,j,k)$，$i+j+k=p$，保留 $r$ 份新值作载体。剩余库存为
\begin{equation}\label{eq:stock}
(p-q_0i,p-q_0j,r-q_0k,q_0p-r).
\end{equation}
若输出 $(x,y,u,v)$ 在库存内、和为 $p$，再平均其补集，返回矩阵为 $J/p^2$，其中
\begin{equation}\label{eq:count-matrix}
J=\begin{pmatrix}p(x-y)+\alpha v&pry-p^2u+\beta v\\-\alpha&-\beta\end{pmatrix},
\quad \alpha=i-j,\quad\beta=rj-pk.
\end{equation}
若 $p(x-y)+\alpha v=\beta$，则 $J^2=pfI$，$pf\equiv p^4\pmod n$。若一组这样的表与其相邻输出 $(x-1,y-1,u+2,v)$ 同时可用，记相邻标量为 $pf,pg$，则
\begin{equation}\label{eq:root-count}
g-f=\alpha n,\qquad
[\Sigma,J_f^{-1}J_g]=U(nH/g),\quad H=n(2j+\alpha)-2p^2,
\end{equation}
其中 $\Sigma=\left(\begin{smallmatrix}-1&r\\0&1\end{smallmatrix}\right)$。
\end{lemma}
\begin{proof}

新载体值满足 $-pz'=D$，所以
\[
z'=-\frac{\alpha a+\beta z}{p^2}.
\]
第一输出组的均值为
\begin{align*}
a'&=\frac{xA+yB+uC+vD}{p}\\
&=\frac{p(x-y)+\alpha v}{p^2}a
 +\frac{pry-p^2u+\beta v}{p^2}z.
\end{align*}
这分别给出矩阵的第二行与第一行；将输出组的和除以 $p$ 给出\eqref{eq:count-matrix}。迹零即得标量平方；代入 $r\equiv-2p\pmod n$ 得其模 $n$ 值。相邻输出使右上项减少 $pn$，故 $g-f=\alpha n$。相乘得到
\[
J_f^{-1}J_g=\begin{pmatrix}1&-\beta n/f\\0&g/f\end{pmatrix},
\]
再与 $\Sigma$ 作交换子即得\eqref{eq:root-count}。
\end{proof}

\section{一整套数字：十一平均、三十四个位置}
取 $p=11,s=1,r=12,n=34$。来源 $(i,j,k)=(4,3,4)$。做两次相同组成的平均后，留下十二份 $D$，库存为
\[(3,5,4,10).
\]
可选择两个不同输出方案 $(3,4,1,3)$ 和 $(2,3,3,3)$。每个方案各自形成一条四次平均的完整返回。

$\alpha=i-j=1$，$\beta=rj-pk=-8$。矩阵为
\[
J_f=\begin{pmatrix}-8&383\\-1&8\end{pmatrix},qquad
J_g=\begin{pmatrix}-8&9\\-1&8\end{pmatrix}.
\]
平方分别为 $-319I$ 和 $55I$，故 $f=-29,g=5$，且 $g-f=34=\alpha n$。它们的商给倍率 $-5/29$，环为 $\Z[1/145]$。再算
\[H=n(2j+\alpha)-2p^2=34\cdot7-242=-4,
\]
得到上剪切 $U(-136/5)$。横向系数 $\alpha=1$，根步长是 $4nR$，便能进入共同终端允许的偶数层。

这张表里的数字不是“某几个输入能平均”的例子。它给出的矩阵等式对任意参数 $a,z$ 恒成立，所以这一个计数表可以处理整个无界的二维配置族。

\section{两张很像的表，为什么比一张强很多}
从第一输出改成 $(x-1,y-1,u+2,v)$，总项数没变，$x-y$ 也没变，迹零条件仍成立；只有右上项少了 $pr+2p^2=pn$。

于是比较两反射会产生三角矩阵；再用换块的对称性消去三角矩阵的倍率，剩下一个剪切。若有几组表，利用剪切参数相加和裴蜀等式，就把它们的步长合成最大公因数。

下面的推论把后面的工作归结为两个清楚的目标：横向系数的最大公因数为一，根系数的最大公因数足够小。满足这两项，群论和终端部分已经不用重新证明。
\begin{corollary}[相邻输出表的充分条件]\label{cor:resource}
设 $p\ge11$、$p<r<2p$。若有限组相邻表满足引理\ref{lem:count-interface}，$0<\alpha<p$，所有 $\alpha$ 的最大公因子为 $1$，所有 $H$ 的最大公因子为 $\delta$，其中 $\delta=1$（$n$ 奇）或 $\delta\in\{2,4\}$（$n$ 偶），则全部合法核心可归零。
\end{corollary}
\begin{proof}

\par\noindent\textbf{根与控制层。}
令 $R$ 为所有约分倍率 $g/f$ 及其逆生成的整数环。它是有限局部化环，反演素数均不除 $n$。某一倍率必不为 $\pm1$：$g=f$ 与 $\alpha\ne0$ 矛盾，而 $g=-f$ 与 $g-f=\alpha n$、$\gcd(f,n)=1$、$n>2$ 矛盾。因此 $R$ 有无限单位。
由\eqref{eq:root-count}先取整数倍得到 $U(nH)$，再作倍率共轭及裴蜀合并，取得 $U(\delta nR)$。引理\ref{lem:transverse}及 $\gcd(\alpha)=1$ 给 $\Gamma(\delta^2n^2,R)$。换到 $(z,y)$，上根变为下根。

\par\noindent\textbf{有限方向与尺度。}
规范反射 $J/p^2$ 模 $n$ 为 $\left(\begin{smallmatrix}1&-\alpha/p^2\\0&-1\end{smallmatrix}\right)$。其条目公因子整除 $\alpha<p$；所需有限标量由安全小数字提供。换块、反射和 $\alpha$ 的裴蜀组合给 $V\equiv U(-p^{-2})$。$n$ 奇时，选奇 $\alpha$；相邻 $f,g$ 奇偶相反，而条目公因子为奇数，所以引理\ref{lem:content}给 $2I$。因 $r>p$，终端的 $j$ 区间包含至少六个整数，覆盖 $\delta$ 个连续值。定理\ref{thm:terminal}适用。
\end{proof}

\chapter{上半带的公式，先看它们在节省哪种位置}\label{ch:upper}
\section{这一讲还要证明什么}
上一讲说明，合格的计数表会自动接出剪切和同余输送。本讲只负责证明：对每个允许的 $p,s$，真的有这样的表。我们仍记 $r=p+s$，$1\le s<p$。

不必一开始记住所有公式。先按新值 $D$ 还剩多少份来理解三种选择。新载体保留 $p+s$ 份，因此可用于输出的 $D$ 只有 $p-s$ 份。
\begin{enumerate}
\item $s$ 很大、接近 $p$ 时，$D$ 剩得少，直接设计不使用 $D$ 的输出。
\item $s$ 在中段时，让旧三类取用更均衡，用一个整数 $e$ 在两种输出之间调整。
\item $s$ 很小时，用另外两级相邻参数保证根系数的最大公因数足够小。
\end{enumerate}
三类设计的算术收尾相同，不同的是剩余位置怎么安排。
\begin{proposition}[上半带的充分性]\label{prop:upper}
若 $p\ge11$、$3p<n<4p$，全部合法输入可归零。
\end{proposition}
\begin{proof}
$p\ge83$ 时，由引理\ref{lem:upper-supply}得到定理\ref{thm:terminal}的资源。$11\le p<83$ 且 $n=3p+s$、$s\le9$ 时用引理\ref{lem:small-supply}；其余小参数由引理\ref{lem:upper-small}处理。三者都接命题\ref{prop:entry}的入口及引理\ref{lem:zero-tail}的收尾。
\end{proof}

\section{一个可行整数不够，需要相邻两个}
典型输出写成 $o_e=(\alpha-e,p-3\alpha-e,2\alpha-s+2e,s)$。当 $e$ 增加一，恰好是一份 $A$、一份 $B$ 换成两份 $C$。这就是上一讲产生剪切的相邻输出。

所有库存限制会变成 $L\le e\le U$。我们必须证明 $L+1\le U$，因为一个可行点只给一个反射，两点才给反射之间的比较。

例如 $p=83,s=10$，取 $\alpha=14$。来源 $(28,14,41)$，剩余库存为 $(27,55,11,73)$。可选 $e=-9,-8$，得到输出
\[(23,50,0,10),qquad(22,49,2,10).
\]
两者都非负、和为83、没有超库存。这里 $e$ 为负完全可以：$e$ 是参数，不是直接取用份数；真正不能为负的是输出的四项。

第二个相邻 $\alpha=15$ 同样给两张表。两个 $\alpha$ 互素；根系数 $H=3n\alpha-2p^2$ 分别为 $-2900,-2123$，最大公因数为一。这就满足上一讲的资源条件。
\begin{lemma}[大素数的三组计数公式]\label{lem:upper-supply}
若 $p\ge83$、$p<r<2p$，则存在满足推论\ref{cor:resource}且 $\delta=\gcd(n,2)$ 的相邻表。
\end{lemma}
\begin{proof}
全部取 $q_0=2$。按下列三种范围选计数。

\par\noindent\textbf{第一类：靠近 $4p$。} 若 $r>3p/2$，置 $d=2p-r$，取 $\alpha=\lceil d/2\rceil+2,\lceil d/2\rceil+3$，来源 $(\alpha,0,p-\alpha)$，输出
\[(l,p-\alpha+l,\alpha-2l,0).
\]
其精确区间为 $L\le l\le U$，其中
\[L=\max(0,\lceil(d-\alpha)/2\rceil),\qquad
U=\min(\lfloor\alpha/2\rfloor,p-2\alpha).
\]
取 $l=L+1,L$。有 $L+1\le\lfloor\alpha/2\rfloor$；由 $2\alpha\le d+7$、$L+1\le d/4+1$，另一约束由 $p\ge5d/4+8$ 保证。$d\le(p-1)/2,p\ge83$ 足够。

\par\noindent\textbf{第二类：上半带中部。} 写 $r=p+s$、$10\le s\le(p-1)/2$。取 $\alpha=\lceil p/6\rceil,\lceil p/6\rceil+1$，来源 $(2\alpha,\alpha,p-3\alpha)$，输出
\[(\alpha-e,p-3\alpha-e,2\alpha-s+2e,s).
\]
可行区间为 $L\le e\le U$，
\[L=\max(5\alpha-p,\lceil s/2\rceil-\alpha),\qquad
U=\min(\alpha,2\alpha+s-\lceil p/2\rceil).
\]
取 $e=L,L+1$。$L+1\le U$ 的四个比较是
\begin{align*}
4\alpha&\le p-1,&3\alpha&\le\lfloor p/2\rfloor+s-1,\\
\lceil s/2\rceil+1&\le2\alpha,&\lceil p/2\rceil+1&\le3\alpha+\lfloor s/2\rfloor.
\end{align*}

下面说明四个比较的余量。因为 $p/6\le\alpha\le(p+11)/6$，第一式在 $p\ge25$ 时成立；
第二式左侧至多比 $\lfloor p/2\rfloor$ 大 $6$，而 $s-1\ge9$；
第三式使用 $s\le(p-1)/2$，在 $p\ge15$ 时足够；
第四式由 $3\alpha\ge\lceil p/2\rceil$ 及 $\lfloor s/2\rfloor\ge5$ 成立。
此外 $4\alpha\le p$ 保证第一类来源不超量，$6\alpha\ge p$ 保证载体来源容量，
$s\le p/2$ 保证新值 $D$ 的库存。因此所列整数区间确实包含相邻两个输出。

\par\noindent\textbf{第三类：紧邻 $3p$ 的九个偏移。} 若 $1\le s\le9$，令 $a_0=\lfloor s/3\rfloor+1$，取 $\alpha=a_0,a_0+1$、$h=\lfloor p/(3\alpha)\rfloor-1,\lfloor p/(3\alpha)\rfloor$。令 $t=p-3\alpha h$，来源和输出为
\[
(\alpha(h+1),\alpha h,p-\alpha(2h+1)),\quad
(\alpha h-e,p-\alpha(2h+1)-e,(\alpha-s)h+\alpha+2e,sh).
\]
输出区间为
\[
L=\max(0,2\alpha-t,\lceil((s-\alpha)h-\alpha)/2\rceil),\quad
U=\min(\alpha h,\alpha h+t-\alpha,\lfloor(sh+\alpha+s-t)/2\rfloor).
\]
可取 $e=L,L+1$，事实上 $L+2\le U$。逐项移项，充分条件是
\[
\begin{gathered}
h\ge4,\quad(3\alpha-s)h\ge\max(s,\alpha+4),\\
sh\ge\max(5\alpha+3-s,3\alpha+4-s),\\
\alpha h\ge\max(4\alpha+4-s,3\alpha+2).
\end{gathered}
\]
对 $\alpha=1,2,3,4,5$，所需 $h$ 下界依次为 $7,12,8,4,4$；$p\ge83$ 时实际最小值为 $26,12,8,5,4$。这同时保证来源及 $sh\le p-s$ 的库存。


\par\noindent\textbf{三类计数的最大公因子。}
到这里已经有真实反射，但还需把它们合成足够浅的根与同余群。
这一步只计算前一讲的 $H=n(2j+\alpha)-2p^2$。
三种公式均满足迹零条件。第一种的 $H=n\alpha-2p^2$，第二种为 $3n\alpha-2p^2$，相邻 $\alpha$ 给 $\gcd(H)=\gcd(n,2)$。第三种相邻 $h$ 的差为 $2\alpha n$，相邻 $\alpha$ 互素；再用 $H\equiv-2p^2\pmod n$ 及奇偶性得到同一结论。所有横向系数含两个相邻整数。
\end{proof}

\section{83是一个统一够用的界，不是数学相变点}
三套计数都可以用 $p\ge83$ 一次保证，但其中两套在小素数上早已能用。偏移 $s\ge10$ 的情形，只需回到四个容量不等式逐项检查，小参数就能补齐。
\begin{lemma}[偏移至少10时的小素数]\label{lem:upper-small}
若 $p\ge11$、$10\le s<p$，引理\ref{lem:upper-supply}的前两种计数已提供推论\ref{cor:resource}的全部资源。
\end{lemma}
\begin{proof}
中部 $s\le(p-1)/2$ 时，四个容量比较对 $p\ge29$ 成立：第一项只需 $p\ge25$；第二项由 $s\ge10$ 保证；第三项只需 $p\ge15$；第四项由 $3\alpha\ge\lceil p/2\rceil$ 保证。唯一较小素数是 $p=23,s=10,11$，此时 $\alpha=4,5$，直接代入四式均成立。
靠近 $4p$ 时，$d=p-s\le(p-1)/2$。充分条件 $p\ge5d/4+8$ 对 $p\ge23$ 成立。若 $p=11,13,17$，还可用 $d\le p-10$，分别得到右侧上界 $9.25,11.75,16.75$。若 $p=19$，$d\le8$ 同样成立；仅余 $d=9$，其 $\alpha=7,8$ 均有 $L=1,U=3$。故两张相邻输出仍在库存内。根和横向系数的最大公因子计算与引理\ref{lem:upper-supply}相同。
\end{proof}
这个引理的作用是缩小真正要查表的区域。不能因为某个统一证明选了83，就把所有小于83的参数都当作没有结构的特殊例子。

\section{直接造出n倍剪切：用一次守恒关系}
记来源为 $(i,j,k)$，新均值满足 $iA+jB+kC=pD$，全局还满足 $pA+pB+rC=0$。如果一个输出组的四项计数允许增加
\[\Delta_t=(ti-p, tj-p, tk-r, n-tp),
\]
那么项数不变，而输出总和增加
\[
t(iA+jB+kC)-(pA+pB+rC)+(n-tp)D=nD.
\]
所以新第一块均值增加 $nD/p=-nz'$，载体 $z'$ 不变。两条返回比较，恰好是 $(a',z')\mapsto(a'-nz',z')$，即 $U(-n)$。

这个根参数 $n$ 直接来自总和守恒。唯一还需检查的是：移动前后的两张输出表都不能超库存。五返回公式就是把这项检查与两个横向参数一起安排。
\begin{lemma}[九个低偏移的有限补充]\label{lem:small-supply}
若 $11\le p<83$、$1\le s\le9$，则全部合法 $K_{p,p+s}$ 可归零。
\end{lemma}
\begin{proof}
先给五返回公式。令 $a=\lceil s/2\rceil$，对 $\alpha=a,a+1$ 写 $p=3\alpha d+\rho$，$0\le\rho<3\alpha$。来源为 $(\alpha(d+1),\alpha d,\alpha d+\rho-\alpha)$，输出为
\[
o_e=(\alpha d-e,\alpha d+\rho-\alpha-e,(\alpha-s)d+\alpha+2e,sd).
\]
对于第一组 $\alpha=a$，还使用 $o_e+\Delta$，其中 $\Delta=(3\alpha-\rho,-\rho,2\rho-3\alpha-s,s)$。可同时容纳 $o_e,o_{e+1},o_e+\Delta$ 的区间是
\begin{align*}
L={}&\max\left(0,5\alpha-2\rho,
\left\lceil\frac{(s-\alpha)d-\alpha}{2}\right\rceil,
\left\lceil\frac{(s-\alpha)d+2\alpha+s-2\rho}{2}\right\rceil\right),\\
U={}&\min\left(\alpha d-\alpha-\mathbf1_{\rho=0},
\left\lfloor\frac{sd+\alpha+s-\rho-2}{2}\right\rfloor,
\left\lfloor\frac{sd+4\alpha+2s-3\rho}{2}\right\rfloor\right).
\end{align*}
逐项比较给充分条件 $d\ge7,sd\ge9\alpha+2-s,(3\alpha-s)d\ge4\alpha+s+2$。例如最紧的三式是 $5\alpha\le\alpha d-\alpha-1$、$10\alpha\le sd+\alpha+s-2$、$(s-\alpha)d+2\alpha+s\le2(\alpha d-\alpha-1)$；其余比较的余量至少分别为 $\alpha d+2\alpha+s-\rho-3$ 和 $\alpha d+2\alpha+s-\rho-1$。第二组只需相邻输出，条件降为 $d\ge4,sd\ge3\alpha+2-s$。由此得到九个充分截点
\[
T_s=(42,24,42,42,63,63,84,84,105).
\]
当 $p\ge T_s$ 时，以上计数直接可用。守恒关系 $pA+pB+rC=0$ 使沿 $\Delta$ 的输出和增加 $nD$，故两返回之商为 $U(-n)$。相邻倍率扩展为 $U(nR)$，两横向系数 $a,a+1$ 由引理\ref{lem:transverse}给 $\Gamma(n^2,R)$，接定理\ref{thm:terminal}。


\par\noindent\textbf{有限剩余。}
余下的114组参数及233对相邻表完整列在附录\ref{app:certificates}。它们的非负库存、迹零等式、$\alpha$ 的最大公因子和 $H$ 的最大公因子由命题\ref{prop:root-certificate}核对。除 $(p,s)=(11,3)$ 外均满足推论\ref{cor:resource}。

\par\noindent\textbf{唯一单列参数 $(p,s)=(11,3)$。}
此时 $n=36$，一般最大公因子接口改由两组明确矩阵完成。
在此例，来源 $(4,2,5)$、输出 $(2,5,1,3)$ 与 $(3,0,2,6)$ 给
\[
J=\begin{pmatrix}-27&568\\-2&27\end{pmatrix},\quad
M=\begin{pmatrix}45&-404\\-2&27\end{pmatrix},\quad MJ^{-1}=U(-36).
\]
库存为 $(3,7,4,8)$，两输出均可用；$J^2=-407I$。相邻输出 $(1,4,3,3)$ 的反射 $G$ 满足
\[J^{-1}G=\begin{pmatrix}1&-972/37\\0&-35/37\end{pmatrix}.
\]
故 $R_0=\Z[1/(5\cdot7\cdot37)]$ 上有 $U(36R_0)$。另一个表 $(4,1,6;2,7,1,1)$ 给 $K=\left(\begin{smallmatrix}-52&905\\-3&52\end{smallmatrix}\right)$、$K^2=-11I$。横向系数 $2,3$ 合并给 $\Gamma(36^2,R_0)$，再取其子群 $\Gamma(4\cdot36^2,R_0)$ 接定理\ref{thm:terminal}。
\end{proof}

\section{剩余114组表，承担的是哪一个有限命题}
到此，无穷多个大 $p$ 已由公式覆盖，小 $p$ 的不少参数也由五返回公式覆盖。还保留114组物理根资源。它们各列一对或几对具体计数，检查非负、总数、迹零以及两个最大公因数。

查表以后，对任意有理参数的终止仍由同余输送证明。表里不枚举数值高度，不说“随机一千次都成功所以全部成功”。附录保留所有计数，你可以用上一讲三十四元的方式逐行读它。

\chapter{位置很多时，为什么还要补一小段}\label{sec:wide}
\section{看上去接近4p，实际还差一个条件}
整数下降证明的是
\[n\ge4p-2+d_p(n).
\]
它不是无条件的 $n\ge4p$。例如 $p=761,n=3045=4p+1$，而 $n=3\cdot5\cdot7\cdot29$ 有四个危险素因子。右边要求 $n\ge3046$，差一个位置。

所以即使临界带 $2p<n<4p$ 已全部完成，也不能忘了这个窄缝。我们需要证明宽维数中被旧不等式漏掉的部分也成立。

\section{一个失败的不等式，反而告诉我们有更多资源}
旧结论已经覆盖所有 $n\ge4p+\lceil\log_2p\rceil$。因此漏掉的只能写成 $n=4p+s$，其中 $s$ 很小。原不等式不成立，又意味着
\[\omega(n)\ge s+3.
\]
也就是说，位置数含很多\emph{不同}素因子。不同素数的乘积长得很快，于是 $n$、进而 $p$ 必须相当大。这给出 $p\ge16s+21$ 等强估计。

这一步改变了我们看剩余参数的方式。不是处理任意 $p,s$，而是处理“$s$ 小，$p$ 被素因子条件逼得很大”的参数。它们有足够的实际副本做一组五步反射；同时欧拉函数足够大，使安全数字有足够密度。

\section{这一讲的七步路线}
先用同样的见证集中和三行折叠，得到加权核心
\[(a^{2p},(-2a+rz)^p,(-pz)^r),\qquad r=p+s.
\]
这次第一块是 $2p$ 份，不能把它当作两个可随意交换的 $p$ 块去套前一讲的对称公式。我们重新写两套真实返回：一套五步反射和一套保留旧载体的三步反射。比较它们，取得双向剪切和明确同余层。

接下来，来源计数允许一整段安全整数 $q$。其对应的行列式成为环中单位。只要这段及其负数包含超过一半的模 $n$ 单位，每个单位就是两者之比：集合 $S$ 与 $uS$ 都超过半数，必相交，交点给 $u=a/b$。

最后一个返回提供局部平移。模五处的步长可能消失，所以先用另一返回处理五进分量，再用指数CRT处理其余素数而不动刚做好的部分。完整列匹配后，精确送到 $(\epsilon^{2p},(-2\epsilon)^p,0^{p+s})$，在其中的 $4p$ 元零和子系统收尾。

下面保留七步的全部计数。第一次阅读可以先对照这条路线确认每组公式承担哪项任务，再第二遍检查上下界。
\begin{proposition}[所有n至少4p的配置]\label{prop:wide}
对 $p\ge5$ 的每个素数，所有 $n\ge4p$ 满足定理\ref{thm:prime}的判据。
\end{proposition}
\begin{proof}

\par\noindent\textbf{第一步：剩余范围自动带来足够的容量。}
由引理\ref{lem:integer-range}，只须考虑
\begin{equation}\label{eq:tail-range}
n=4p+s,\quad1\le s<\lceil\log_2p\rceil,\quad d=\omega(n)\ge s+3.
\end{equation}
$n$ 奇时至少含四个不同奇素数，$n\ge1155,\varphi(n)\ge480$；$n$ 偶时 $s$ 为正偶数，至少含五个素数，$n\ge2310,\varphi(n)\ge480$。每增加一个不同素因子，最小乘积至少乘 $13$，最小欧拉函数至少乘 $12$。归纳给
\begin{equation}\label{eq:tail-estimates}
p\ge16s+21,\qquad \varphi(n)>8\cdot2^d+7s+36.
\end{equation}

具体地，先以这两个基点归纳得到
\[
n\ge65d-111,\qquad \varphi(n)>8\cdot2^d+7d+15.
\]
后续最小乘积每增加一项至少乘 $13$，大于第一式右侧的线性增量；
最小欧拉函数至少乘 $12$，也大于第二式右侧的增量。
代入 $d\ge s+3$，得到 $n\ge65s+84$，于是
\[
p=(n-s)/4\ge16s+21.
\]
第二个估计同样由 $7d+15\ge7s+36$ 得到。容斥从而给
\[(p-1)\varphi(n)/n\ge[\varphi(n)-s-4]/4>2^d-1.
\]
故引理\ref{lem:short-choice}的配比选择在这些宽维数仍有效。


\par\noindent\textbf{第二步：建立加权三块核心。}
保护单点 $w$ 及见证，建立两个块，用引理\ref{lem:concentrate}集中 $a-w$ 的支持，再避开 $w$ 建立第三、第四个块。按命题\ref{prop:entry}的见证修复，使前三块不全同余。写 $p=3m+\epsilon$。先混合 $b,c$，再混合 $a,b$，使 $m(a+b+c)+\epsilon a-pd_0$ 在全部危险处非零；各条件至多禁止一个配比。三行折叠后得到 $(A^{2p},B^p,d_0^p,x_1,\ldots,x_s)$ 且 $A-d_0$ 全支持。平均 $p-s$ 份 $A$ 与全部单点，留下 $p+s$ 份 $A$ 作载体；将其余三个 $p$ 块再折叠一次，进入
\begin{equation}\label{eq:wide-core}
(a^{2p},(-2a+rz)^p,(-pz)^r),\qquad r=p+s.
\end{equation}
本原参数的合法性仍为 $\gcd(a+pz,n)=1$。


\par\noindent\textbf{第三步：五次平均组成反射。}
取 $\epsilon=\pm1$、$p\equiv\epsilon\pmod4$，$t_0=(p-\epsilon)/4,t_1=t_0+s$。对 $t=t_0,t_1$，令 $b=p-s-3t$、$\beta=nt-p^2+ps$。先平均两组 $(A^{p-b-t},B^t,C^b)$，产生 $2p$ 份 $D=(sa+\beta z)/p$，保留 $r$ 份。再平均两组
\[
(A^{2h+t-b},B^h,C^{p-2t+b-3h},D^t)
\]
及全部余下 $p$ 份。五次操作给 $J_{t,h}/p^2$，其中
\begin{equation}\label{eq:tail-J}
J_{t,h}=\begin{pmatrix}\beta&prh-p^2(p-2t+b-3h)+t\beta\\-s&-\beta\end{pmatrix},
\quad J_{t,h}^2=pf_{t,h}I,
\end{equation}
\[
f_{t,h}=p^3+n[4t^2+(s-2p)t-sh].
\]

为看清五步构造的库存，第一阶段后旧 $A,B,C$ 分别剩
$2(b+t),p-2t,r-2b$ 份。留出最终载体后，新值 $D$ 还剩 $p-s$ 份。
第二阶段要做两组组成相同的平均，所以条件分别是
\[
0\le2h\le p-2t,\qquad
0\le2(p-2t+b-3h)\le r-2b,
\]
\[
0\le2h+t-b\le b+t,\qquad 2t\le p-s.
\]
这四项分别限制旧 $B$、旧 $C$、旧 $A$、新 $D$ 的用量。
将它们合并，容量允许
\[
\max(0,\lceil(5p-5s-16t)/6\rceil)\le h\le b.
\]
写 $t=t_0+\delta$、$\delta=0,s$，未截断端点差为 $(p+\epsilon-2s-4\delta)/12>2$，且 $b\ge4$；故有相邻 $h,h+1$。代入还得 $b-t\le0$、$b\le(p-2t)/2$，这些不等式正好保证两阶段所有库存约束。


\par\noindent\textbf{第四步：第二组反射与共同根理想。}
保留旧 $A$ 的三步反射为
\[
A_j=\begin{pmatrix}0&-d_j\\-1&0\end{pmatrix},\quad d_j=p^2-nj,\qquad
\lceil(p-s)/6\rceil\le j\le\lfloor(p-s)/4\rfloor.
\]
其两重复组各取 $(A^{2j},B^j,C^{p-3j})$，余组取补集，保留 $r$ 份 $A$。因 $p-s\ge36$，可取三个相邻 $j$，记标量 $d_0,d_1,d_2$。对固定 $t$，相邻 $h$ 的标量为 $f,g_t=f-sn$，且互素：$f\equiv p(p-4t)^2\pmod s$，而 $p-4t\equiv\epsilon\pmod s$。
相邻反射商与 $A_j^{-1}A_{j+1}$ 的交换子给
\[
U\left(\frac{\beta_t n^2}{d_1g_t}\right),\qquad
\gcd(\beta_{t_0},\beta_{t_1})=\gcd(n,5)=g.
\]
后式用差为 $ns$、$\gcd(\beta_{t_0},s)=1$ 以及 $\beta_t\equiv p(s-p)\pmod n$。相邻倍率及引理\ref{lem:localize}给 $R_0$，其反演素数来自这些互素标量且不除 $n$。于是有双向 $gn^2R_0$ 根和
\begin{equation}\label{eq:tail-level}
\Gamma(M,R_0),\qquad M=g^2n^4.
\end{equation}
$R_0$ 含绝对值不同于 $1$ 的相邻标量，所以满足定理\ref{thm:arithmetic-external}的单位条件。保载体返回以 $p-3j=\pm1$ 激活 $p$，第三个 $A$ 的倍率激活 $d_2$。


\par\noindent\textbf{第五步：可用数字覆盖足够多的单位。}
此时已经有深同余群，可以安全地激活与控制层互素的返回的逆。
我们从位置计数读出一个连续数字区间，证明其安全数字足以生成模 $n$ 单位群。
一般保留 $A$ 的三步返回为
\[
M_{q,j}=\begin{pmatrix}p-q&nj-pq\\-1&0\end{pmatrix},\qquad q=3j+k,
\]
其中 $0\le j\le\lfloor p/2\rfloor$、$0\le k\le\lfloor r/2\rfloor$、$\lceil r/2\rceil\le j+k\le p$。这些计数覆盖全部整数区间
\[
I=[L,2p-1],\qquad L=r/2+2\ (r\text{偶}),\quad L=(r+9)/2\ (r\text{奇}).
\]
固定 $j$ 所给区间为 $[2j+\lceil r/2\rceil,\min(3j+\lfloor r/2\rfloor,2j+p)]$，相邻区间相接即证覆盖。安全 $q$ 的行列式与 $M$ 互素，引理\ref{lem:inverse}和\ref{lem:localize}允许将其加入最终环 $R$，并重建\eqref{eq:tail-level}。由\eqref{eq:tail-estimates}及 $|I|\ge(3n-7s-36)/8$，$I$ 中单位数超过 $\varphi(n)/4$。
由于 $I\subset(0,n/2)$，区间中的单位及其负数互不重复。
令它们组成 $S$，则 $|S|>\varphi(n)/2$。
给定任意单位 $u$，两个集合 $S,uS$ 在大小为 $\varphi(n)$ 的单位群中必相交；
若 $a=ub$，$a,b\in S$，则 $u=a/b$。
因此每个模 $n$ 单位是其中两者之比，$R^\times$ 模 $n$ 满射。


\par\noindent\textbf{第六步：逐素数调整方向。}
在 $(z,y)=(z,a+pz)$ 中，$V=J_{t,h}A_j/p^3$ 满足
\[
V\equiv U(-5/p)\pmod n,\quad
\operatorname{tr}V=2-\frac n{p^3}(pt+t^2+ph+sj).
\]
若 $n\equiv2\pmod4$，选偶 $h$，则迹为 $2\pmod4$；$4\mid n$ 时自动成立。若 $v_3(n)=1$，规范化平方参数满足
\[
\frac c3\equiv\frac{n/3}{p^3}[(p-s)(j+t-h)-5t^2]\pmod3.
\]
相邻 $j$ 避开引理\ref{lem:local-cycle}的唯一坏类。若 $5\mid n$，先在 $I$ 选安全 $q\equiv-p\pmod5$；固定该类的容斥主项为 $|I|\rho/4$，误差小于 $2^{d-1}$，而\eqref{eq:tail-estimates}使主项更大。其 $M_{q,j}$ 模 $5$ 给非平凡平移，先完成五进方向；$V\bmod5=I$，随后使其指数为完整五进阶的倍数，再用 CRT 完成其余局部方向。得到 $z/y=0\pmod M$。

\par\noindent\textbf{第七步：匹配尺度并送到含零子系统。}
单位 $d_1/d_0=1-n/d_0$ 和 $d_2/d_0=1-2n/d_0$ 生成模 $M$ 到模 $n$ 的单位核。奇素数处由二项式展开逐层升位；$n\equiv2\pmod4$ 时第二个单位生成 $1\pmod4$ 的部分，第一个补另一个类。于是 $R^\times$ 模 $M$ 满射。引理\ref{lem:fiber}将列精确送到 $(0,\epsilon)$。最终配置为 $(\epsilon^{2p},(-2\epsilon)^p,0^{p+s})$；选全部非零位置和 $p$ 个零，使用引理\ref{lem:integer-range}已独立证明的 $4p$ 维结论。
\end{proof}

\section{为什么这个收尾没有绕圈}
最终调用的是 $4p$ 元判据。它早在整数下降中独立完成，因为 $d_p(4p)=1$，故 $4p\ge4p-2+1$。我们不是用“所有大于等于 $4p$ 已解决”来证明“所有大于等于 $4p$ 已解决”，而是用其中那个先已证明的固定维数 $4p$ 消去所有非零位置。

这种先后关系值得在整份证明里一直保持：先证明一个含零或等块的小工具，后面的矩阵构造只需把状态送进它能处理的范围。

\chapter{回到三平均：几种终止方法放在一起看}\label{ch:three}
\section{最小的奇素数，不能简单代入所有大p公式}
三平均的最终阈值是七。七元已被奇数余量构造覆盖，九元可排成 $3\times3$，先平均三行，再平均三列。因此真正还需处理的是八元、十元和所有十一元以上。

它们展示三种不同的终止方法：十一元以上保持整数平方和下降；八元让一个本原参数的绝对值下降；十元让一个由局部类型决定的整数二次量下降。有限步证明并不一定要所有分支用同一个下降量。

\section{十一元以上：三个余数类使分类更直接}
模三只有三个类。有四个不同三重值，就必有两个同余；有三个三重值，就看它们是否各占一类。很多一般 $p$ 的长选择，在这里变成 $(a,a,x)$ 或 $(a,b,x)$。

最细的地方是危险素数可能为二。如果选择 $(a,x,x)$，只靠“还留着一份 $a$”不能从 $2x\equiv2a\pmod2$ 推出 $x\equiv a$。证明会额外使用碰撞已经给出的 $x=3b-2a$，并留下 $a,b$ 的副本。这个\emph{已知的整数仿射关系}修复了模二失效。

全部分类的剩余容量来自 $n\ge d_3(n)+10$。所以一旦需要保护最多 $d_3(n)$ 个位置，仍有足够轻值可选。
\begin{lemma}[三平均的全部n至少11]\label{lem:triple-large}
三平均在全部 $n\ge11$ 上满足 $G$ 为三幂的充要条件。
\end{lemma}

\begin{proof}
\par\noindent\textbf{准备：统一的保护容量。}
记 $d=d_3(n)$，$f(a)$ 表示实际数值 $a$ 的重数。
有
\[
n\ge d+10.
\]
当 $d\le1$ 时显然；$d=2$ 时 $n\ge2\cdot7=14$；
当 $d\ge3$ 时，非三素因子的最小乘积给 $n\ge2\cdot5^{d-1}\ge d+10$。
每个危险主类的补集有两或三个位置，预留至多 $d$ 个位置便能同时保护。
在单一模三余数类的轻值池中，每个实际值至多两份：
池大小达到 $d+3$ 时能安全地选一个非恒等整数三元组，达到 $d+4$ 时还能避开一个指定均值。
后一结论使用四个候选数的删一比较。

从 $3X$ 开始。若已有两个三重值，无须初始化。
只有一个三重值 $a$ 且 $f(a)\ge5$ 时，平均 $(a,a,b)$，$b\ne a$，
留下三份 $a$ 并产生新的三重值。
若 $f(a)=3,4$，其外至少 $n-4\ge d+6$ 个位置，可用保护池选出均值不同于 $a$ 的三元组。
若没有三重值，先用保护池造出一个 $a$；它至多有五份，
其外至少 $d+5$ 个未改变的初值仍被 $3$ 整除，第二步再造一个不同三重值。

下面给每个非终端状态一个保持两个三重值的安全整数续步。
先说明一种将反复用到的安全性：若未选位置留下 $a,b$，而 $x=3b-2a$，
则假设输出全同余于 $c$ 会先得到 $a=b=c$，再得到 $x=c$。
因此选择 $(a,x,x)$ 时，即使危险素数为 $2$，也无须除以 $2$ 来证明安全。
同理，$x=3b-2a,y=3a-2b$ 时，$(x,y,z)$ 和 $(x,b,z)$ 均可用未选锚点及组和证明安全。

\par\noindent\textbf{第一类：至少三个三重值。}
若至少有四个，模三只有三个类，故有两个不同重值 $a,b$ 同余。
平均 $(a,a,b)$，至少另外两个三重值未动。
若恰三个，且 $a,b$ 同余，记第三个为 $c$。
两个候选 $(a,a,b)$、$(b,b,a)$ 的均值不同，所以至少一个不同于 $c$，保留 $c$ 与新均值即可。

若三个重值分处三个模三类，且另有轻值 $x$，取同余于 $x$ 的重值 $a$，平均 $(a,a,x)$，
另外两个重值保留。
若只有三个实际值，平均 $(a,b,c)$；若至少两个值有四份，它们仍有三份。
只剩重数为 $3,3,n-6$ 的情形，设最后一个值为 $c$。
如果新均值碰撞到 $c$，便有 $a+b=2c$；零和进一步给 $nc=0$。
于是得到明确终端
\[
(a^3,(-a)^3,0^{n-6}),
\]
三次 $(a,-a,0)$ 即可完成。

\par\noindent\textbf{第二类：两个三重值模三同余。}
设它们为 $a,b$。若 $f(a)\ge4$，用 $(b,b,a)$；若 $f(b)\ge4$，对称处理。
至少一个旧值仍有三份，新均值与它不同。
因此只剩 $f(a)=f(b)=3$，其外的轻值池 $S$ 满足 $|S|=n-6\ge d+4$。

若 $S$ 中有与锚点同余的 $x$，两个候选 $(a,a,x)$、$(b,b,x)$ 都安全。
若都碰撞，必有 $2a+x=3b$、$2b+x=3a$，相减得到 $5(a-b)=0$，矛盾。
若 $S$ 只占另外一个余数类，直接用保护池。
若它占另外两个类 $S_1,S_2$，考虑跨类对 $(c,e)$。
候选 $(a,c,e)$ 和 $(b,c,e)$ 不能选中的危险，仅可能是某个二元危险集恰等于 $\{c,e\}$。
每个危险素数至多阻塞一个跨类位置对，而
\[
|S_1||S_2|\ge |S|-1=n-7>d.
\]
故有未被阻塞的对。两个安全候选若同时碰撞，会给
$a+c+e=3b$、$b+c+e=3a$，推出 $4(a-b)=0$，仍矛盾。

\par\noindent\textbf{第三类：两个三重值模三不同。}
若轻值 $x\equiv a\pmod3$，先试 $(a,a,x)$。
失败只可能是 $f(a)\le4$ 且 $x=3b-2a$。
若这个 $x$ 出现两份，改用 $(a,x,x)$，得到 $2b-a$，与 $a,b$ 均不同。
其安全性由前述仿射见证保证。
因此还未成功时，$a$ 类至多有一个例外单点 $x=3b-2a$，
$b$ 类至多有一个例外单点 $y=3a-2b$，其余位置都在第三类，记其池为 $Z$。

若 $x,y$ 都存在，$f(a)+f(b)+2\le10<n$，所以 $Z$ 非空。
任取 $z\in Z$，平均 $(x,y,z)$，两个锚点完整保留，且 $x-y=5(b-a)\ne0$。

若只有 $x$，先说明 $Z$ 非空。否则安全操作 $(a,a,x)$ 后只有两个实际值 $a,b$，
而它们模三不同。合法性使 $a-b$ 与 $n$ 的每个非三素因子互素，
模三不同又补上素数 $3$，故 $\gcd(a-b,n)=1$。
零和迫使 $n\mid f(a)$，与 $0<f(a)<n$ 矛盾。
所以可选 $(x,b,z)$。若 $f(b)\ge4$，两个旧重值均保留。
若 $f(b)=3$，则 $|Z|\ge n-8\ge3$，而每个轻值至多两份，
可以选 $z\ne5a-4b$，使新均值不同于保留的 $a$。
只有 $y$ 的情形对称。

若 $x,y$ 都不存在，$Z$ 仍不能为空，理由同上。
当 $f(a)=f(b)=3$ 时，用 $|Z|\ge d+4$ 的保护池；
当二者均至少四份时，任取 $(a,b,z)$，两个旧重值仍保留。
最后设 $f(a)=3,f(b)\ge4$，选 $z\ne2b-a$ 即可。
若无此选择，$Z$ 全部等于 $2b-a$，其位置数 $r$ 只能为 $1$ 或 $2$。
零和变成
\[
(3-r)(a-b)+nb=0.
\]
所有值都在 $b+(a-b)\Z$ 中，所以合法性与模三不同仍给 $\gcd(a-b,n)=1$。
上式却迫使 $n\mid3-r$，不可能。

\par\noindent\textbf{有限终止。}
三类已经穷尽全部储备状态。每个非终端操作都安全、非恒等，并保持整数与两个不同三重值。
平方和严格下降且为非负整数，故有限步后必到前述相反值终端。
对 $3X$ 得到的同一平均词也使 $X$ 归零。
\end{proof}

\section{八元：从加减倍数回到欧几里得算法}
八元先进入 $(u^4,v^3,-4u-3v)$。取 $(x,y)=(u,u+v)$ 后，合法性是 $y$ 奇。目标是把正整数 $|y|$ 降到一。

两条真实返回提供 $x\mapsto x\pm2y$。这让我们能在一个模 $2y$ 的剩余类里选一个离 $2y$ 最近的 $x_0$。第三条返回把第二坐标改成 $2(x_0-2y)$；两个输出都有公因子二，约去以后，新的 $|y|$ 比原来小。这和欧几里得算法“减去合适倍数后取小余数”的逻辑一样。

例如分析列 $(x,y)=(1,3)$。加一次 $2y$，选 $x_0=7$，此时 $|7-6|=1<3$。施加第三返回得到 $(30,2)$，取本原代表是 $(15,1)$。一轮便把 $|y|$ 从三降到一。选择平移的准则很具体：让 $x_0$ 尽量接近 $2y$，以便下一步的第二坐标变小。
\begin{lemma}[三平均八元的欧几里得下降]\label{lem:eight}
三平均八元的合法配置均可归零。
\end{lemma}
\begin{proof}
若无五重值，存在安全非恒等整数三元组。奇偶分布为 $4+4$ 或 $2+6$；前者任意三元组补集都保留奇偶见证。若没有非恒等整三元组，则至多两个模三类，且每个至少三项的类实际常值；无五重值迫使 $(a^4,b^4)$，由零和得 $b=-a$，违反合法性。在 $2+6$ 情形，若多数池没有整三元组，无五重值迫使它的两模三类各至少两项；其两项和覆盖全部模三类，任一少数位置可与其中两项配成安全整三元组。
整数能量下降至五重值 $u$ 后，写例外为 $a,b,c$，选 $c-u$ 奇。平均 $(u,a,b)$，得到合法 $B_8(u,v)=(u^4,v^3,-4u-3v)$。

\par\noindent\textbf{在核心上作下降。}
令 $x=u,y=u+v$，本原参数满足 $y$ 奇。下列实际返回给三个作用：
\[
R:(x,y)\mapsto(x+y,-y),\quad
S:(x,y)\mapsto(y-x,y),\quad
F:(x,y)\mapsto(3x+3y,2x-4y).
\]
为列出操作，置 $w=-4u-3v$、$a_0=(u+2v)/3$、$b_0=(2u+v)/3$、$q_0=-u-2v/3$、$h_0=-u/3-4v/9$、$k_0=u/3-2v/9$。三个词依次为
\begin{align*}
R:&\ (v,v,u),(v,u,u),(w,a_0,b_0),(a_0,b_0,u);\\
S:&\ (w,v,u),(q_0,q_0,u),(q_0,u,u),
(h_0,v,k_0),(h_0,v,k_0),(v/9,v/9,k_0);\\
F:&\ (v,u,u),(w,b_0,b_0),(v,u,u).
\end{align*}

这里每个括号表示在\emph{当时}选取的三个值。
新值 $a_0,b_0,\ldots$ 是前面步骤产生的值，反复写同一值时要取不同位置。
这些词的输出重新具有重数 $4,3,1$，并在整体非零尺度下给出所列参数矩阵。
因此它们可以在任意后继核心上重放。
逐次代入可核对库存及最终重数 $4,3,1$。$SR,RS$ 给平移 $x\mapsto x\pm2y$。若 $|y|>1$，先用 $S$ 使 $x$ 奇，再平移至 $|x_0-2y|<|y|$；等号仅在 $y\mid x$ 时发生，与本原性矛盾。
余数的平衡选择与通常的欧几里得算法相同：允许 $x$ 加减 $2y$ 后，
把 $x_0$ 选在以 $2y$ 为中心、宽度为 $2|y|$ 的区间内。
边界等号意味着 $y\mid x$，本原性会迫使 $|y|=1$，与当前阶段矛盾。
施加 $F$ 后两坐标均偶，约去 $2$ 后第二坐标奇，故本原输出满足 $|y'|\le|x_0-2y|<|y|$。最终 $|y|=1$，平移和 $S$ 送到 $(0,\pm1)$。配置 $(0^4,1^3,-3)$ 先平均 $(-3,0,0)$，再三次平均 $(-1,0,1)$。
\end{proof}

\section{十元：有限表为什么可以覆盖无限有理输入}
十元的危险素数只有二和五。先用安全整数下降取得六重值，再经一个四步桥进入 $(u^6,v^3,-6u-3v)$。

取方向坐标 $(X,Y)=(u,3u+2v)$ 后再取本原代表。不同奇偶类型对应不同的尺度 $h=1$ 或 $2$，势能写成
\[h^2(5X^2+Y^2).
\]
不要把 $h$ 忽略：同一形式的有理方向，在清分母与取本原代表时可能需要不同的整数尺度，比较下降必须用同一真实规则。
\begin{lemma}[三平均十元的二次势能下降]\label{lem:ten}
三平均十元的合法配置均可归零。
\end{lemma}
\begin{proof}

\par\noindent\textbf{第一步：取得六重值。}
先证明无六重值时有安全整三元组。危险奇偶集只能有两项，记为 $P$；危险模五集有两或三项，记为 $F$。若二者均不存在，任意整三元组安全；不存在整三元组会迫使 $(a^5,b^5)$，违反合法性。仅有一种危险集时，在多数池内用模三两类的两项和；若无法配出安全三元组，多数池中至少六项为同一实际值。
两危险集同时存在时，令 $K=(P\cup F)^c$。其内部整三元组安全。若 $|K|\ge7$，无六重值使两个模三类均至少两项，与任一外部位置配对即得安全选择。$|K|=6$ 时只剩 $r^5s$，全部外部余数必须为 $s$；从 $P\setminus F,F\setminus P$ 各取一项，连同 $K$ 中的 $s$ 即可。$|K|=5$ 时 $P,F$ 不交，只剩 $r^4s$ 或 $r^5$；前者同理。后者若 $P,F$ 有跨类余数，配一项 $K$ 得三类各一；否则外部五项同余，取一项 $P$、两项 $F$。每种选择补集均保留两种见证。


\par\noindent\textbf{第二步：从六重值到二维核心。}
能量下降至 $(u^6,a,b,c,d)$。令 $e_a=a-u$ 等，和为 $-10u$。若某个例外差是模 $10$ 单位，平均其余三个例外，一步进入合法 $B_{10}(u,v)=(u^6,v^3,-6u-3v)$。否则四个差模 $10$ 为两份 $5$ 及一对相反非零偶数，故有单位和的跨类配对，设为 $e_a+e_b$。令
\[q=(-5u-a-b)/3,\quad r=(-5u+2a+2b)/9,\quad s=(-11u-a-b)/27.
\]
词 $(u,c,d),(q,a,b),(u,q,r),(u,q,r)$ 给 $u^3r s^6=B_{10}(s,u)$，且 $s-u=-(40u+e_a+e_b)/27$ 是模 $10$ 单位。


\par\noindent\textbf{第三步：用有限局部类型选择下降返回。}
在本原 $(X,Y)$ 方向 $X=u,Y=3u+2v$ 上，合法类型为：两者奇且 $X+Y\equiv0\pmod4$，取 $h=1$；或奇偶相反，取 $h=2$。另要求 $5\nmid Y$。整数势能为 $h^2(5X^2+Y^2)$。命题\ref{prop:ten-certificate}给七个实际返回及十二个完整二次覆盖：除 $X=0$ 外，总有返回严格降低该势能。$X=0$ 时已有零坐标，固定一个零，对其余九个位置采用两层三进网络。由整数下降完成。
\end{proof}

\section{具体看一个二次下降条件}
设输入属于 $h=1$ 的一类，采用
\[M_4=\begin{pmatrix}1&1\\5&-3\end{pmatrix},
\]
该类的输出内容为四，输出尺度仍为一。乘掉内容分母后，输入减输出的量为
\[
16(5X^2+Y^2)-[5(X+Y)^2+(5X-3Y)^2]
=50X^2+20XY+2Y^2.
\]
它等于 $2(5X+Y)^2$。所以除了直线 $Y=-5X$，这张矩阵都严格下降。在那条直线上不能就停止，还要检查同一局部类型的其他矩阵；完整表做的正是多个下降区域的联合覆盖。

原问题的有理数可以任意大，但决定内容和模四类型只需要有限余数。固定一种类型后，是否下降变成实变量 $t=Y/X$ 上的二次不等式。二次函数只在至多两个根处变号，因此可精确检查全部实区间，进而覆盖其中所有有理方向。
\begin{proposition}[十元二次覆盖证书]\label{prop:ten-certificate}
引理\ref{lem:ten}的合法本原方向上，以下七个矩阵中至少一个严格降低 $h^2(5X^2+Y^2)$，唯一相容例外为 $X=0$：
\[
\begin{gathered}
M_1=\begin{pmatrix}1&1\\-5&-1\end{pmatrix},\quad
M_2=\begin{pmatrix}1&0\\-5&2\end{pmatrix},\quad
M_3=\begin{pmatrix}1&-2\\0&3\end{pmatrix},\quad
M_4=\begin{pmatrix}1&1\\5&-3\end{pmatrix},\\
M_5=\begin{pmatrix}4&1\\0&-3\end{pmatrix},\quad
M_6=\begin{pmatrix}1&-2\\5&2\end{pmatrix},\quad
M_7=\begin{pmatrix}1&0\\-10&-1\end{pmatrix}.
\end{gathered}
\]
\end{proposition}
\begin{proof}
置 $U=u,V=v,W=-6u-3v$，$Q=(-5u-2v)/3,R=(u+2v)/3,S=(2u+v)/3,T=-2u-v/3$。返回依次平均的当前三元组为
\begin{center}\small
\begin{tabular}{cl}\toprule
矩阵&平均词\\\midrule
$M_1$&$(W,V,U),(V,U,U),(V,U,U)$\\
$M_2$&$(W,V,U),(Q,U,U),(V,V,U),(Q,R,U),(Q,R,U)$\\
$M_3$&$(W,U,U),(V,U,U),(V,S,U),(V,S,U)$\\
$M_4$&$(W,U,U),(V,U,U),(V,U,U)$\\
$M_5$&$(W,V,U),(V,V,U),(R,U,U),(R,U,U)$\\
$M_6$&$(V,U,U),(W,S,U),(V,S,U),(V,S,U)$\\
$M_7$&$(W,V,V),(V,U,U),(T,S,U),(T,S,U),(S,U,U)$\\\bottomrule
\end{tabular}\end{center}
代入两个独立参数可核对全部线性形式和库存。行列式为 $4,2,3,-8,-12,12,-1$，所以本原化内容只含 $2,3$。模 $96$ 记录各内容及输出模四类型，得到5120个合法局部类和12种签名。
若 $M_i=(a,b;c,d)$、内容为 $g_i$、输出尺度为 $h_i$，
输出本原列是 $M_i(X,Y)^T/g_i$，所以输出势能为
\[
\frac{h_i^2}{g_i^2}
\left(5(aX+bY)^2+(cX+dY)^2\right).
\]
把它与输入势能比较，乘正数 $g_i^2$，再在 $X\ne0$ 时除以 $X^2$，
下降条件是
\[
g_i^2h^2(5+t^2)-h_i^2[5(a+bt)^2+(c+dt)^2]>0,\qquad t=Y/X.
\]
表\ref{tab:ten-data}列出所有系数。每行各多项式的实根把实直线分成有限区间，逐区间及逐根检查符号；二次项检查无穷方向。剩余的有限点只可能为 $0$，被 $5\nmid Y$ 排除；两行留下无穷方向，其中仅一行与 $(0,1)$ 的局部签名相容，正是终端 $X=0$。由此得到完整覆盖。
\end{proof}
\begingroup\small
\begin{longtable}{r r l p{83mm} p{24mm}}
\caption{十二个二次覆盖：$i:(a,b,c)$ 表示 $a+bt+ct^2>0$}\label{tab:ten-data}\\
\toprule 类&个数&代表&$i:(a,b,c)$&未严格覆盖\\\midrule\endfirsthead
\toprule 类&个数&代表&$i:(a,b,c)$&未严格覆盖\\\midrule\endhead
1&768&$(0,1)$&1:(-10,-20,-2); 2:(-40,80,0); 3:(0,80,-112); 6:(50,0,-8); 7:(-400,-80,0)&$\infty$\\
2&768&$(1,0)$&1:(-10,-20,-2); 2:(-10,20,0); 3:(0,80,-112); 5:(0,-160,8); 7:(-400,-80,0)&0\\
3&256&$(1,2)$&1:(-10,-20,-2); 2:(-10,20,0); 3:(160,80,-80); 5:(640,-40,130); 7:(-400,-80,0)&$\varnothing$\\
4&384&$(1,3)$&1:(-40,-80,-8); 3:(0,20,-28); 4:(50,20,2); 7:(-100,-20,0)&$\varnothing$\\
5&768&$(1,6)$&1:(-10,-20,-2); 2:(-10,20,0); 3:(0,80,-112); 5:(0,-40,2); 7:(-400,-80,0)&0\\
6&384&$(1,7)$&1:(-40,-80,-8); 3:(0,20,-28); 4:(200,80,8); 7:(-100,-20,0)&$\varnothing$\\
7&256&$(1,8)$&1:(-10,-20,-2); 2:(-10,20,0); 3:(160,80,-80); 5:(2560,-160,520); 7:(-400,-80,0)&$\varnothing$\\
8&128&$(1,11)$&1:(-40,-80,-8); 3:(40,20,-20); 4:(50,20,2); 7:(-100,-20,0)&$\varnothing$\\
9&128&$(1,23)$&1:(-40,-80,-8); 3:(40,20,-20); 4:(200,80,8); 7:(-100,-20,0)&$\varnothing$\\
10&256&$(2,1)$&1:(-10,-20,-2); 2:(-40,80,0); 3:(160,80,-80); 6:(2760,0,480); 7:(-400,-80,0)&$\varnothing$\\
11&768&$(2,3)$&1:(-10,-20,-2); 2:(-40,80,0); 3:(0,80,-112); 6:(200,0,-32); 7:(-400,-80,0)&$\infty$\\
12&256&$(4,5)$&1:(-10,-20,-2); 2:(-40,80,0); 3:(160,80,-80); 6:(690,0,120); 7:(-400,-80,0)&$\varnothing$\\
\bottomrule\end{longtable}
\endgroup
表里的 $\infty$ 表示 $X=0$，这正是含零终点。有限未覆盖点 $0$ 表示 $Y=0$，违反 $5\nmid Y$。核验应同时检查“没有下降”与“这个方向是否真的属于该局部类型”；不能仅凭一个实图上看起来被盖住了就结束。

\chapter{小参数中的计算机证明，实际证明了什么}\label{app:small}
\section{先列出缺口，避免把小例子当作全部证明}
五平均只需额外处理十六、十七、十八元；七平均需要二十二至二十六元。其他位置数已经由一般范围覆盖。

其中五平均的三个点与七平均的二十三至二十六元共用一条路线：假设所有安全候选都发生碰撞，把碰撞写成线性方程，再证明这些方程与零和、同余见证不相容。七平均二十二元另用矩阵群和一个带尺度标签的有限图。

因此这里有两种有限证书：一类穷尽线性方程\emph{模式}，另一类穷尽一个已经证明足够的有限\emph{状态空间}。两者都不对实际输入数值规定高度上限。

\section{一条碰撞等式，怎样变成线性方程}
设两个锚点为 $a,b$，令 $\Delta=a-b$。把每个轻值写成
\[x_i=b+\Delta t_i.
\]
这样 $b$ 对应参数零，$a$ 对应参数一。选择一个轻值子集 $S$，再加 $i$ 份 $a$、若干份 $b$ 补足 $p$ 项，均值为
\[b+\frac{\Delta}{p}\left(i+\sum_{j\in S}t_j\right).
\]
若它恰好碰撞到 $a$，就必须有 $\sum_{j\in S}t_j=p-i$。这是关于实际有理参数的严格线性方程，不只是模 $p$ 同余。

实际 $t_i$ 有无穷多个，但它们的余数 $z_i=t_i\pmod p$ 只有有限种。每个余数模式决定哪些候选和必须写哪些方程。于是可以对有限种模式逐一做精确消元，而消元得到的是每个系统的\emph{全部有理解}。

\section{保护集中的每个位置，为什么有一个私人素数}
先为危险支持各留一个位置，然后删掉所有多余的保护位置，得到一个“再删任何一项都不行”的集合 $T$。对其中某个位置 $v$，总有一个危险素数的例外支持只由它保护；否则 $v$ 还能删。这就是私人素数。

在这个素数处，其他保护位置与锚点 $b$ 同余，只有 $v$ 不同。因此，如果未保护部分相对 $b$ 的总和为零，整个配置的相对总和就只剩一个非零项，违背全局零和。这个矛盾只需要知道\emph{总和}，不必强迫每个轻值都与 $b$ 同余。
\begin{lemma}[五平均和七平均的碰撞分类]\label{lem:collision-small}
五平均的 $n=16,17,18$ 及七平均的 $n=23,24,25,26$ 均满足判据。
\end{lemma}
\begin{proof}
引理\ref{lem:initialize}适用。引理\ref{lem:background}、\ref{lem:heavy}、\ref{lem:two-heavy}只留下不同余的两个重值 $a,b$，$f(a)=p+\alpha,f(b)=p+\beta$、$\alpha\ge\beta\ge0$、$\alpha+\beta\le p-3$。
五平均十六元另有三值 $(a^5,b^6,c^5)$，其中 $a<b<c$，三者模五不同。
选整数调节 $(a^i,c^{4-i},b)$，保留五份 $b$。
若均值碰撞到 $b$，$i=1$ 给 $a+3c=4b$，代入零和得 $13b-5c=0$，
从而 $b\equiv c\pmod2$，再由前式得到 $a\equiv c\pmod2$，违反合法性。
$i=3$ 对称；$i=2$ 给 $a+c=2b$，零和则给 $16b=0$，故 $b=0,a=-c$。
因此这项容量边界也已经处理。

\par\noindent\textbf{保护集与私人素数。}
从需要保护的轻值危险集选一个包含极小横截集 $T$，$h=|T|\le d_p(n)$。
这里“包含极小”是指删去其中任意一个位置就不再同时击中全部危险支持。
因此对每个被保护位置，总有某个危险支持只在这个位置与保护集相交；
否则可以把该位置删掉。这就是所谓私人素数的来源。
每个 $v\in T$ 有私人素数 $\ell$，其危险支持与 $T$ 恰交于 $v$；此处 $a=b\pmod\ell$，其余保护位置与 $b$ 同余。
故未保护部分相对 $b$ 的总和不可能为 $0\pmod\ell$。
若它为零，总和
\[
\sum_i(x_i-b)=\sum_{i\notin T}(x_i-b)+\sum_{i\in T}(x_i-b)
\]
模 $\ell$ 就恰剩私人位置 $v$ 的一个非零差。
但左侧因原配置零和且 $\ell\mid n$ 应为零，矛盾。
这个论证只需要未保护部分的\emph{总和}为零，不要求其中每个值都同余于 $b$。

\par\noindent\textbf{把轻值写成相对锚点的参数。}
令 $\Delta=a-b$、$t_i=(x_i-b)/\Delta$、$z_i=t_i\pmod p$。锚点同余修复后，余数 $0,1$ 各至多一个，实际参数为 $p,1-p$；各余类重数小于 $p$。对轻值子集 $S$，$k=|S|\le p$，若 $i\equiv-\sum_{S}z\pmod p$ 满足 $i\le\min(\alpha,p-k)$，
所选组由轻值子集 $S$、$i$ 份 $a$ 和 $p-k-i$ 份 $b$ 组成，
它的和为 $pb+\Delta(i+\sum_S t)$。
均值碰撞到 $a=b+\Delta$，恰好要求 $i+\sum_S t=p$。
因此碰撞要求
\begin{equation}\label{eq:collision}
\sum_{S}t=p-i;
\end{equation}
若还满足 $p-k-i\le\beta$，两个旧重值均保留，因而该模式不可能失败。从另一锚点，$j\equiv\sum_S z-k\pmod p$、$j\le\min(\beta,p-k)$ 时，碰撞要求 $\sum_S t=k+j-p$。

\par\noindent\textbf{有限线性分类与总和矛盾。}
每个轻值的模 $p$ 余数只有 $p$ 种，各余类的重数也有明确上界。
因此方程模式只有有限个，尽管实际轻值 $t_i$ 可以是任意有理数。
命题\ref{prop:collision-certificate}穷尽这些线性系统。五平均令 $c=n-10-h\in\{5,6,7,8\}$。$c=7,8$ 无解；$c=6$ 的两个解均为整数且 $D=5+\alpha+\sum t_i=15$；$c=5$ 时所有 $D$ 的分母为奇数。对 $n=17,18$，有保护时私人素数整除 $\Delta$，而 $D=15$ 给总和矛盾；无保护时没有相容模式。对 $n=16,h=0$，$16b+15\Delta=0$ 迫使 $\Delta$ 偶，全部整数参数使配置同奇偶；$h=1$ 时 $\Delta$ 偶而 $D$ 分母奇，同样违反私人素数 $2$ 的总和条件。
七平均可在每个余数类内对\eqref{eq:collision}取对称平均，压缩为变量 $y_r$、重数 $f_r$；总和保持。某类在一个方程中部分选取时，交换位置进一步强迫该类实际常值。令 $n-h=21+e$，则 $1\le e\le5$。证书给十一种相容模式，
每个实际参数均为整数，且 $0<D=7+\alpha+\sum f_ry_r<21+e$。$h>0$ 时私人素数给总和矛盾；$h=0$ 时所有值属于 $b+\Delta\Z$，合法性给 $\gcd(\Delta,n)=1$，零和迫使 $n\mid D$，又矛盾。因此总有安全续步，整数能量下降完成。
\end{proof}

\section{如何读一行碰撞表}
五平均表的一行是 $c=6,\alpha=1,\beta=0$，余数列 $(3,4,4,4,4)$，唯一实际解为 $(-7,4,4,4,4)$。这些有理未知数一旦全部满足碰撞方程，就被迫取这五个值。加上锚点贡献，
\[D=5+1+(-7+4+4+4+4)=15.
\]
当位置数为十七或十八、存在保护位置时，私人素数整除 $\Delta$，未保护相对总和 $\Delta D$ 便在该素数处为零，正好触发上面的矛盾。

这不是验证了五个具体输入，而是证明一种可能的失败模式只能有这个解，然后把它排除。
\begin{proposition}[碰撞分类证书]\label{prop:collision-certificate}
在引理\ref{lem:collision-small}的方程下，五平均的903个无序模式恰有表\ref{tab:five-data}所列22个相容解；七平均的57026个压缩模式恰有表\ref{tab:seven-data}所列11个相容解。每个相容系统满秩。
\end{proposition}
\begin{proof}
五平均遍历 $(\alpha,\beta)=(0,0),(1,0),(1,1),(2,0)$、$c=5,6,7,8$，以及长度 $c-\alpha-\beta$ 的非降余数列。删除含五个相同余数、或余数 $0,1$ 各出现超过一次的列。对每个非空至多五项子集，写入\eqref{eq:collision}及另一锚点方程；同时保留两个旧重值时写入矛盾方程 $0=1$。按最左非零元逐列消元，每次整数行消元后约去最大公因子，最后以有理数回代。所得全部相容解如下表。
七平均遍历 $1\le e\le5$、$0\le\beta\le\alpha$、$\alpha+\beta\le4$，长度为 $7+e-\alpha-\beta$。每类重数小于七，余数 $0,1$ 各至多一次。按各类选择数 $0\le k_r\le f_r$、$1\le\sum k_r\le7$ 建立压缩方程，用同一消元。每个重复类都出现在至少一条部分选取方程中，故其实际值等于表列解。程序逐行回代并比较整个解集，得到所述计数及完整表。
\end{proof}
\begingroup\small
\begin{longtable}{rrrlll}
\caption{五平均的全部相容系统}\label{tab:five-data}\\
\toprule $c$&$\alpha$&$\beta$&余数列&唯一解&$D$\\\midrule\endfirsthead
\toprule $c$&$\alpha$&$\beta$&余数列&唯一解&$D$\\\midrule\endhead
5&0&0&$(0,2,2,2,2)$&$(5,-8,-8,-8,-8)$&$-22$\\
5&0&0&$(0,3,3,3,3)$&$(5,-2,-2,-2,-2)$&$2$\\
5&0&0&$(0,4,4,4,4)$&$(5,\frac{-7}{2},\frac{-7}{2},\frac{-7}{2},\frac{-7}{2})$&$-4$\\
5&0&0&$(1,2,2,2,2)$&$(-4,\frac{9}{2},\frac{9}{2},\frac{9}{2},\frac{9}{2})$&$19$\\
5&0&0&$(1,3,3,3,3)$&$(-4,3,3,3,3)$&$13$\\
5&0&0&$(1,4,4,4,4)$&$(-4,9,9,9,9)$&$37$\\
5&0&0&$(2,2,2,2,3)$&$(-3,-3,-3,-3,8)$&$1$\\
5&0&0&$(2,2,2,2,4)$&$(7,7,7,7,-16)$&$17$\\
5&0&0&$(2,3,3,3,3)$&$(12,-7,-7,-7,-7)$&$-11$\\
5&0&0&$(2,4,4,4,4)$&$(17,-6,-6,-6,-6)$&$-2$\\
5&0&0&$(3,3,3,3,4)$&$(8,8,8,8,-11)$&$26$\\
5&0&0&$(3,4,4,4,4)$&$(-7,4,4,4,4)$&$14$\\
5&1&0&$(2,2,2,2)$&$(2,2,2,2)$&$14$\\
5&1&0&$(3,3,3,3)$&$(\frac{4}{3},\frac{4}{3},\frac{4}{3},\frac{4}{3})$&$\frac{34}{3}$\\
5&1&0&$(3,4,4,4)$&$(-7,4,4,4)$&$11$\\
5&1&0&$(4,4,4,4)$&$(4,4,4,4)$&$22$\\
6&1&0&$(3,4,4,4,4)$&$(-7,4,4,4,4)$&$15$\\
5&1&1&$(0,3,3)$&$(5,-7,-7)$&$-3$\\
5&1&1&$(1,3,3)$&$(-4,8,8)$&$18$\\
5&2&0&$(2,2,2)$&$(2,2,2)$&$13$\\
5&2&0&$(2,2,3)$&$(2,2,3)$&$14$\\
6&2&0&$(2,2,2,2)$&$(2,2,2,2)$&$15$\\
\bottomrule\end{longtable}\endgroup

\begingroup\small
\begin{longtable}{rrrllll}
\caption{七平均的全部相容系统}\label{tab:seven-data}\\
\toprule $e$&$\alpha$&$\beta$&余数&重数&唯一解&$D$\\\midrule\endfirsthead
\toprule $e$&$\alpha$&$\beta$&余数&重数&唯一解&$D$\\\midrule\endhead
1&2&0&$(2)$&$(6)$&$(2)$&$21$\\
1&3&0&$(2)$&$(5)$&$(2)$&$20$\\
1&3&0&$(2,3)$&$(4,1)$&$(2,3)$&$21$\\
1&4&0&$(2)$&$(4)$&$(2)$&$19$\\
1&4&0&$(2,3)$&$(3,1)$&$(2,3)$&$20$\\
1&4&0&$(2,4)$&$(3,1)$&$(2,4)$&$21$\\
1&4&0&$(2,3)$&$(2,2)$&$(2,3)$&$21$\\
2&3&0&$(2)$&$(6)$&$(2)$&$22$\\
2&4&0&$(2)$&$(5)$&$(2)$&$21$\\
2&4&0&$(2,3)$&$(4,1)$&$(2,3)$&$22$\\
3&4&0&$(2)$&$(6)$&$(2)$&$23$\\
\bottomrule\end{longtable}\endgroup

\section{二十二元的3872个状态是什么}
控制模数为 $M=1936$。列记为 $(z,y)$，其中 $y$ 是模 $M$ 单位。先记录比值 $x=z/y\pmod M$，有1936种。

但实际可用的环 $R=\Z[1/35]$ 的单位模 $M$ 只覆盖全部单位的一半，记这半为 $U$，另一半为 $3U$。所以再加一位 $\epsilon$：$y\in U$ 时为零，$y\in3U$ 时为一。总共 $1936\times2=3872$ 个状态。

这是为何单纯遍历1936个方向不够。两个方向相同的列，可能一个处于允许单位倍数类，另一个不在；必须看真实宏能否把它们之间的尺度差也调整掉。
\begin{lemma}[七平均二十二元的方向与尺度]\label{lem:seven22}
七平均二十二元的合法配置均可归零。
\end{lemma}
\begin{proof}

\par\noindent\textbf{第一步：三个真实返回与深层。}
入口使用命题\ref{prop:entry}证明中转换前的 $(a^{14},(-2a+z)^7,-7z)$。三条三步返回为
\[
A=\begin{pmatrix}4&1\\-1&0\end{pmatrix},\quad
F=\begin{pmatrix}1&2\\0&7\end{pmatrix},\quad
J=\begin{pmatrix}1&2\\2&-1\end{pmatrix}.
\]
$A$ 保留一份旧 $A$、平均两组 $(A^6,B)$ 再平均 $(A,B^5,C)$；$F$ 保留 $C$、平均两组 $(A^5,B^2)$ 再平均 $(A^4,B^3)$；$J$ 保留一份旧 $B$、平均两组 $(A^5,B^2)$ 再平均 $(A^4,B^2,C)$。实际尺度均为 $1/7$。有 $(FA^2)^2=-7I,J^2=5I$，故各逆可实现。
令 $H=AFJA^{-1}=\left(\begin{smallmatrix}-7&-62\\0&5\end{smallmatrix}\right)$，则 $[F,H]=U(396/35)$。倍率生成 $R=\Z[1/35]$；用 $C=\left(\begin{smallmatrix}1&4\\0&-1\end{smallmatrix}\right)$ 换基，$CU(t)C=U(-t)$、$CAU(t)A^{-1}C=L(t)$，得到 $\Gamma(396^2,R)$。

\par\noindent\textbf{第二步：去掉多余的三进控制层。}
记 $M=1936$、$X=F^2/7$。命题\ref{prop:seven-certificate}给 $A^{440}=X^{440}=I\pmod M$，且两者模 $9$ 生成全部 $\SL_2(\Z/9\Z)$。其中三个方向 $I+3E_{12},I+3E_{21},I+3\diag(1,-1)$ 的立方逐层生成模 $81$ 的高层核。因此对任意 $g\in\Gamma(M,R)$，可选可实现 $w\equiv I\pmod M$、$w\equiv g\pmod{81}$；由 $396^2=81M$ 得 $w^{-1}g$ 已在深核中，故 $\Gamma(M,R)$ 可实现。

\par\noindent\textbf{第三步：有限方向与尺度同时命中终端。}
同一证书给 $R^\times$ 模 $M$ 的像 $U$ 有 $440$ 个元素，全部单位为 $U\sqcup3U$。合法列 $(z,y)=(z,a+7z)$ 的有限标记是 $(z/y,\epsilon)$，$\epsilon$ 记录 $y$ 所属单位陪集。证书证明全部 $3872$ 个标记都能到 $(0,0)$。引理\ref{lem:fiber}再精确送到 $(0,s)$，其中 $s\in R^\times$。配置 $(s^{14},(-2s)^7,0)$ 中，四份 $s$、两份 $-2s$ 和一个零组成零和七元组，接引理\ref{lem:zero-tail}。
\end{proof}

\section{有限图的边，怎样提升成无限输入的操作}
矩阵 $(a,b;c,d)$ 作用于列后，第二坐标乘上 $cx+d$。于是 $x$ 按分式线性变换更新，而 $\epsilon$ 要加上 $cx+d$ 属于哪一类的标记。

证书从终端 $(0,0)$ 出发，把每条允许边都展开，直到没有新状态；结果恰好覆盖3872个。生成元已经有真实逆，所以这条父边树可以反向走回终端。至于实际整数有多大，最后由 $\Gamma(M,R)$ 的精确纤维引理处理。
\begin{proposition}[二十二元有限轨道证书]\label{prop:seven-certificate}
引理\ref{lem:seven22}中的 $A^{440},(F^2/7)^{440}$ 模 $1936$ 均为恒等，模 $9$ 生成 $648$ 个行列式一矩阵。环 $\Z[1/35]$ 的单位像模 $1936$ 为 $440$ 元子群 $U$，单位群为 $U\sqcup3U$。三个返回及其逆在 $1936\times2$ 个方向--尺度标记上作用传递，最大父边深度为八。
\end{proposition}
\begin{proof}
快速整数模幂检查前一等式。从恒等矩阵出发，逐个乘两个模九生成元，保留尚未出现的矩阵，闭包大小为648；逐项检查行列式为一。因 $|\SL_2(\Z/9\Z)|=3^3(3^2-1)3=648$，这就是全群。
同理从 $1$ 乘 $-1,5,7$ 得 $U$，并检查 $U,3U$ 不交且并为全部880个单位。令 $\chi$ 为其二元商标记。在 $(z,y)$ 坐标中使用
\[
A=\begin{pmatrix}7&-1\\22&-3\end{pmatrix},\quad
F=\begin{pmatrix}7&0\\44&1\end{pmatrix},\quad
J=\begin{pmatrix}-15&2\\-110&15\end{pmatrix}.
\]
矩阵 $(a,b;c,d)$ 对标记的作用是
\[
(x,\epsilon)\longmapsto((ax+b)(cx+d)^{-1},\epsilon+\chi(cx+d)).
\]
从 $(0,0)$ 对 $A,A^{-1},F,F^{-1},J$ 逐层闭包，得到全部3872个状态，最深层为八。对每个状态的全部边重新检查分母可逆及目标在集合内，即给完整有限证书。
\end{proof}

“最大父边深度为八”只说有限标记图最多走八条这样的边，不能直接说原题只需八次平均。每条宏本身有若干原子操作，精确纤维输送也还需要平均词。有限图深度与原子词长度是两个不同计数。

\chapter{把整条证明重新连起来}\label{sec:completion}
\section{现在可以不用公式，先把逻辑讲一遍}
起点是任意有理配置。先减去均值、清分母、约去公因子，得到本原整数零和代表。

如果 $G$ 含非 $p$ 素因子，那个素数下的共同非零余数永远不变，立即判不可解。若 $G$ 只有 $p$ 因子，我们按位置数选择证明路线。

位置足够多，保持两个重复块，选择安全整数平均，使整数平方和下降；旧容量不等式遗漏的窄段，则利用其大量素因子带来的强计数估计另行补齐。

临界带位置紧张，先安全地压缩成三个块。通过真实返回得到剪切、倍率和逆，进而取得足够大的主同余子群。有限层面将参数与某个零和终端列对齐，再由主同余子群把它们精确对齐。最后真的平均一个零和组，接含零收尾。

小素数 $3,5,7$ 中一般容量不适用的部分，由前两讲的独立下降或有限证书补上。每一块都提供指定范围内的全部输入结论，下面把范围合并。
\begin{proof}[定理\ref{thm:prime}的证明]
必要性由引理\ref{lem:obstruction}给出，下界由命题\ref{prop:lower}给出。对于充分性，$p\ge11$ 时：$n\ge4p$ 用命题\ref{prop:wide}；$n=3p$ 用命题\ref{prop:threep}；其余临界带先用命题\ref{prop:entry}压缩，再按下表完成。
\begin{center}
\begin{tabular}{ll}
\toprule 参数范围 & 核内结论\\\midrule
$n=2p+r,\ 1\le r<p,\ r$ 奇 & 命题\ref{prop:odd}\\
$r=2$ & 命题\ref{prop:r2}\\
$4\le r<p,\ r$ 偶 & 命题\ref{prop:even}\\
$3p<n<4p$ & 命题\ref{prop:upper}\\\bottomrule
\end{tabular}
\end{center}
$p=5$ 时，$11\le n\le15$ 已由上述下半带及等块结论处理；$n=16,17,18$ 用引理\ref{lem:collision-small}；$n=19$ 满足引理\ref{lem:integer-range}；$n\ge20$ 用命题\ref{prop:wide}。
$p=7$ 时，下半带处理 $15\le n\le21$；$n=22$ 用引理\ref{lem:seven22}；$23\le n\le26$ 用引理\ref{lem:collision-small}；$n=27$ 满足引理\ref{lem:integer-range}；更大维数用命题\ref{prop:wide}。
$p=3$ 时，$n=7$ 用命题\ref{prop:entry}及\ref{prop:odd}，$n=8,10$ 用引理\ref{lem:eight}、\ref{lem:ten}，$n\ge11$ 用引理\ref{lem:triple-large}。$n=9$ 将位置排成 $3\times3$，先平均各行再平均各列，六次操作即达到全局均值。
这些范围覆盖全部奇素数及全部 $n\ge2p+1$。推论\ref{cor:threshold}随定理和第\ref{sec:obstruction}讲的二平均论证成立。
\end{proof}

\section{如果有人问：你们到底用了多深的数学}
必要条件用最大公因数和模运算；整数路线用重复值、零和选择、抽屉原理和平方和；压缩用线性组合和CRT；核心矩阵的制造用真实整数计数、标量周期和交换子。

最深的外部输入出现在从“已经实现上下两整族剪切”到“已经实现某个主同余子群”的那一步，使用 Morris 和 Serre 的定理。取得它以后，把同余相同的两列精确连接，又是裴蜀和二阶矩阵的直接计算。

所以整个证明既不是纯粹平均值不等式，也不是套一个数论黑箱就结束。具体的位置构造提供可用矩阵，数论把这些矩阵的能力扩展到任意有理参数，最后仍要回到实际位置归零。

\section{读完以后，最值得自己回答的几个问题}
\begin{enumerate}
\item 为什么 $p$ 和 $n$ 的素因子承担不同角色？答：$p$ 决定实际分母，非 $p$ 的 $n$ 素因子承担共同余数障碍。
\item 为什么平方和下降不总能证明有限终止？答：有理正数可以无限下降；整数格才给离散性。
\item 为什么可逆矩阵不直接把列送到零？答：它只先送到含零和组的非零核心，最后由奇异平均归零。
\item 为什么同余子群出现以后仍需要有限计算？答：它能在一个剩余类里精确输送，但初始剩余类还未必是终端类。
\item 为什么证明还分很多计数情形？答：不同块重数提供的真实副本数量不同；同一抽象矩阵不一定在每种库存下都能实现。
\end{enumerate}
这些问题说明了判据简短而充分性很长的原因。公式 $G=p^e$ 只描述障碍；要证明它也是充分条件，还必须造出不会丢掉余数见证、能在有限步内到达终点的路径。

\appendix
\chapter{剩余根计数表：怎样逐行复查}\label{app:certificates}
附录保留原讲义的114组参数、233对相邻表。第一列 $p$ 是每次平均元数，第二列 $s$ 给 $n=3p+s$。七元组前面三项是来源 $(i,j,k)$，后面四项是输出 $(x,y,u,v)$。

读一行时，先算四类库存 $(p-2i,p-2j,p+s-2k,p-s)$，确认输出逐项不超量、和为 $p$。然后算 $\alpha=i-j$、$\beta=(p+s)j-pk$，检查迹零条件 $p(x-y)+\alpha v=\beta$。最后比较同一参数的几对表，计算 $\alpha$ 和根系数 $H$ 的最大公因数。

第13讲的 $p=11,s=1$ 已经示范了一整行。这里每行承担同一种有限算术检查，不另引入新的终止原理。
\begin{proposition}[根计数证书]\label{prop:root-certificate}
对 $11\le p<83$ 的十八个素数及 $1\le s\le9$，共有162组参数。其中48组满足引理\ref{lem:small-supply}的五返回截点；余下114组在表\ref{tab:root-data}中给出233对相邻输出。各表满足\eqref{eq:stock}及\eqref{eq:count-matrix}的迹零条件。除 $(11,3)$ 外，横向系数的最大公因子为 $1$，根系数的最大公因子为推论\ref{cor:resource}允许的 $\delta$。
\end{proposition}
\begin{proof}
每行七元组为 $(i,j,k,x,y,u,v)$。对两表分别验证
\[
i+j+k=x+y+u+v=p,\quad
0\le(x,y,u,v)\le(p-2i,p-2j,p+s-2k,p-s),
\]
以及 $p(x-y)+(i-j)v=(p+s)j-pk$。第二输出与第一输出相差 $(-1,-1,2,0)$。随后按\eqref{eq:root-count}计算 $H$ 并取最大公因子；这些均为表列整数的直接运算。对每个参数检查它恰属于“达到截点”或“表内”之一，证明覆盖无遗漏。例外 $(11,3)$ 已由引理\ref{lem:small-supply}中的显式矩阵处理。
\end{proof}
\begingroup\small
\begin{longtable}{rrll}
\caption{114组参数的相邻反射计数}\label{tab:root-data}\\
\toprule $p$&$s$&第一输出&相邻输出\\\midrule\endfirsthead
\toprule $p$&$s$&第一输出&相邻输出\\\midrule\endhead
11&1&$(4,3,4,3,4,1,3)$&$(4,3,4,2,3,3,3)$\\
11&2&$(4,3,4,2,3,0,6)$&$(4,3,4,1,2,2,6)$\\
11&2&$(4,2,5,3,6,0,2)$&$(4,2,5,2,5,2,2)$\\
11&3&$(4,2,5,2,5,1,3)$&$(4,2,5,1,4,3,3)$\\
11&4&$(4,3,4,3,3,4,1)$&$(4,3,4,2,2,6,1)$\\
11&4&$(4,2,5,2,5,0,4)$&$(4,2,5,1,4,2,4)$\\
11&5&$(4,3,4,3,3,1,4)$&$(4,3,4,2,2,3,4)$\\
11&5&$(4,2,5,1,4,1,5)$&$(4,2,5,0,3,3,5)$\\
11&6&$(3,2,6,3,6,1,1)$&$(3,2,6,2,5,3,1)$\\
11&6&$(3,1,7,1,7,0,3)$&$(3,1,7,0,6,2,3)$\\
11&7&$(3,2,6,2,5,1,3)$&$(3,2,6,1,4,3,3)$\\
11&7&$(5,2,4,1,2,7,1)$&$(5,2,4,0,1,9,1)$\\
11&8&$(4,3,4,3,2,4,2)$&$(4,3,4,2,1,6,2)$\\
11&8&$(3,0,8,1,9,1,0)$&$(3,0,8,0,8,3,0)$\\
11&9&$(2,0,9,1,10,0,0)$&$(2,0,9,0,9,2,0)$\\
11&9&$(3,0,8,1,9,1,0)$&$(3,0,8,0,8,3,0)$\\
13&1&$(5,4,4,3,3,3,4)$&$(5,4,4,2,2,5,4)$\\
13&1&$(5,3,5,1,4,0,8)$&$(5,3,5,0,3,2,8)$\\
13&2&$(4,3,6,2,5,0,6)$&$(4,3,6,1,4,2,6)$\\
13&2&$(5,4,4,2,2,1,8)$&$(5,4,4,1,1,3,8)$\\
13&3&$(5,2,6,3,7,1,2)$&$(5,2,6,2,6,3,2)$\\
13&3&$(5,1,7,1,8,0,4)$&$(5,1,7,0,7,2,4)$\\
13&4&$(5,4,4,3,2,5,3)$&$(5,4,4,2,1,7,3)$\\
13&4&$(4,2,7,2,7,0,4)$&$(4,2,7,1,6,2,4)$\\
13&5&$(4,3,6,4,6,1,2)$&$(4,3,6,3,5,3,2)$\\
13&5&$(5,4,4,3,2,1,7)$&$(5,4,4,2,1,3,7)$\\
13&6&$(4,3,6,3,5,0,5)$&$(4,3,6,2,4,2,5)$\\
13&6&$(4,2,7,1,6,0,6)$&$(4,2,7,0,5,2,6)$\\
13&7&$(3,2,8,3,8,1,1)$&$(3,2,8,2,7,3,1)$\\
13&7&$(5,4,4,3,1,7,2)$&$(5,4,4,2,0,9,2)$\\
13&8&$(3,2,8,2,7,1,3)$&$(3,2,8,1,6,3,3)$\\
13&8&$(5,2,6,3,6,3,1)$&$(5,2,6,2,5,5,1)$\\
13&9&$(4,3,6,5,6,1,1)$&$(4,3,6,4,5,3,1)$\\
17&1&$(6,5,6,5,6,1,5)$&$(6,5,6,4,5,3,5)$\\
17&1&$(7,6,4,3,1,7,6)$&$(7,6,4,2,0,9,6)$\\
17&2&$(5,4,8,2,6,1,8)$&$(5,4,8,1,5,3,8)$\\
17&2&$(6,5,6,3,4,0,10)$&$(6,5,6,2,3,2,10)$\\
17&3&$(6,4,7,4,7,0,6)$&$(6,4,7,3,6,2,6)$\\
17&3&$(6,3,8,4,9,1,3)$&$(6,3,8,3,8,3,3)$\\
17&4&$(6,5,6,5,5,4,3)$&$(6,5,6,4,4,6,3)$\\
17&4&$(5,3,9,2,8,1,6)$&$(5,3,9,1,7,3,6)$\\
17&4&$(6,4,7,3,6,0,8)$&$(6,4,7,2,5,2,8)$\\
17&5&$(5,4,8,5,8,1,3)$&$(5,4,8,4,7,3,3)$\\
17&5&$(6,5,6,4,4,1,8)$&$(6,5,6,3,3,3,8)$\\
17&6&$(4,3,10,5,11,0,1)$&$(4,3,10,4,10,2,1)$\\
17&6&$(5,4,8,3,6,1,7)$&$(5,4,8,2,5,3,7)$\\
17&7&$(4,3,10,3,9,1,4)$&$(4,3,10,2,8,3,4)$\\
17&7&$(6,5,6,5,4,7,1)$&$(6,5,6,4,3,9,1)$\\
17&8&$(4,3,10,2,8,0,7)$&$(4,3,10,1,7,2,7)$\\
17&8&$(6,5,6,5,4,2,6)$&$(6,5,6,4,3,4,6)$\\
17&9&$(5,4,8,6,8,1,2)$&$(5,4,8,5,7,3,2)$\\
17&9&$(5,3,9,3,8,1,5)$&$(5,3,9,2,7,3,5)$\\
19&1&$(6,5,8,5,8,1,5)$&$(6,5,8,4,7,3,5)$\\
19&1&$(7,6,6,5,5,3,6)$&$(7,6,6,4,4,5,6)$\\
19&2&$(6,5,8,3,6,0,10)$&$(6,5,8,2,5,2,10)$\\
19&2&$(7,6,6,3,3,1,12)$&$(7,6,6,2,2,3,12)$\\
19&3&$(6,4,9,4,9,0,6)$&$(6,4,9,3,8,2,6)$\\
19&3&$(8,6,5,3,2,5,9)$&$(8,6,5,2,1,7,9)$\\
19&3&$(7,4,8,5,9,1,4)$&$(7,4,8,4,8,3,4)$\\
19&4&$(6,5,8,7,9,2,1)$&$(6,5,8,6,8,4,1)$\\
19&4&$(7,6,6,5,4,5,5)$&$(7,6,6,4,3,7,5)$\\
19&5&$(5,4,10,6,11,1,1)$&$(5,4,10,5,10,3,1)$\\
19&5&$(6,5,8,5,7,1,6)$&$(6,5,8,4,6,3,6)$\\
19&6&$(5,4,10,4,9,1,5)$&$(5,4,10,3,8,3,5)$\\
19&6&$(6,5,8,3,5,0,11)$&$(6,5,8,2,4,2,11)$\\
19&7&$(5,4,10,2,7,1,9)$&$(5,4,10,1,6,3,9)$\\
19&7&$(7,6,6,5,3,7,4)$&$(7,6,6,4,2,9,4)$\\
19&8&$(4,3,12,3,11,0,5)$&$(4,3,12,2,10,2,5)$\\
19&8&$(6,5,8,7,8,2,2)$&$(6,5,8,6,7,4,2)$\\
19&9&$(4,3,12,1,9,1,8)$&$(4,3,12,0,8,3,8)$\\
19&9&$(6,5,8,5,6,1,7)$&$(6,5,8,4,5,3,7)$\\
23&1&$(7,6,10,6,10,1,6)$&$(7,6,10,5,9,3,6)$\\
23&1&$(8,7,8,7,8,1,7)$&$(8,7,8,6,7,3,7)$\\
23&2&$(7,6,10,3,7,1,12)$&$(7,6,10,2,6,3,12)$\\
23&2&$(8,7,8,4,5,0,14)$&$(8,7,8,3,4,2,14)$\\
23&3&$(8,6,9,5,8,1,9)$&$(8,6,9,4,7,3,9)$\\
23&3&$(8,5,10,6,11,1,5)$&$(8,5,10,5,10,3,5)$\\
23&4&$(7,6,10,8,11,3,1)$&$(7,6,10,7,10,5,1)$\\
23&4&$(8,7,8,7,7,4,5)$&$(8,7,8,6,6,6,5)$\\
23&5&$(6,5,12,7,13,1,2)$&$(6,5,12,6,12,3,2)$\\
23&5&$(7,6,10,6,9,1,7)$&$(7,6,10,5,8,3,7)$\\
23&6&$(6,5,12,5,11,0,7)$&$(6,5,12,4,10,2,7)$\\
23&6&$(7,6,10,3,6,1,13)$&$(7,6,10,2,5,3,13)$\\
23&7&$(6,5,12,2,8,1,12)$&$(6,5,12,1,7,3,12)$\\
23&7&$(8,7,8,7,6,7,3)$&$(8,7,8,6,5,9,3)$\\
23&8&$(5,4,14,2,11,1,9)$&$(5,4,14,1,10,3,9)$\\
23&8&$(7,6,10,9,11,1,2)$&$(7,6,10,8,10,3,2)$\\
23&9&$(7,6,10,6,8,1,8)$&$(7,6,10,5,7,3,8)$\\
23&9&$(5,3,15,5,16,0,2)$&$(5,3,15,4,15,2,2)$\\
29&1&$(9,8,12,8,12,1,8)$&$(9,8,12,7,11,3,8)$\\
29&1&$(10,9,10,9,10,1,9)$&$(10,9,10,8,9,3,9)$\\
29&3&$(10,8,11,7,10,0,12)$&$(10,8,11,6,9,2,12)$\\
29&3&$(12,10,7,5,2,7,15)$&$(12,10,7,4,1,9,15)$\\
29&3&$(9,6,14,7,15,1,6)$&$(9,6,14,6,14,3,6)$\\
29&4&$(9,8,12,10,13,3,3)$&$(9,8,12,9,12,5,3)$\\
29&4&$(10,9,10,9,9,4,7)$&$(10,9,10,8,8,6,7)$\\
29&5&$(8,7,14,8,14,1,6)$&$(8,7,14,7,13,3,6)$\\
29&5&$(9,8,12,7,10,1,11)$&$(9,8,12,6,9,3,11)$\\
29&6&$(7,6,16,6,15,1,7)$&$(7,6,16,5,14,3,7)$\\
29&6&$(8,7,14,5,11,0,13)$&$(8,7,14,4,10,2,13)$\\
29&7&$(7,6,16,3,12,1,13)$&$(7,6,16,2,11,3,13)$\\
29&7&$(8,7,14,1,7,1,20)$&$(8,7,14,0,6,3,20)$\\
29&8&$(9,8,12,10,12,1,6)$&$(9,8,12,9,11,3,6)$\\
29&8&$(10,9,10,8,7,0,14)$&$(10,9,10,7,6,2,14)$\\
29&9&$(8,7,14,9,14,1,5)$&$(8,7,14,8,13,3,5)$\\
29&9&$(9,8,12,6,8,1,14)$&$(9,8,12,5,7,3,14)$\\
31&1&$(9,8,14,8,14,1,8)$&$(9,8,14,7,13,3,8)$\\
31&1&$(10,9,12,9,12,1,9)$&$(10,9,12,8,11,3,9)$\\
31&3&$(10,8,13,7,12,0,12)$&$(10,8,13,6,11,2,12)$\\
31&3&$(12,10,9,7,6,3,15)$&$(12,10,9,6,5,5,15)$\\
31&3&$(10,7,14,8,15,1,7)$&$(10,7,14,7,14,3,7)$\\
31&4&$(9,8,14,10,15,5,1)$&$(9,8,14,9,14,7,1)$\\
31&4&$(10,9,12,11,13,2,5)$&$(10,9,12,10,12,4,5)$\\
31&5&$(8,7,16,9,17,1,4)$&$(8,7,16,8,16,3,4)$\\
31&5&$(9,8,14,8,13,1,9)$&$(9,8,14,7,12,3,9)$\\
31&6&$(8,7,16,6,14,0,11)$&$(8,7,16,5,13,2,11)$\\
31&6&$(9,8,14,4,9,1,17)$&$(9,8,14,3,8,3,17)$\\
31&7&$(8,7,16,2,10,1,18)$&$(8,7,16,1,9,3,18)$\\
31&7&$(10,9,12,11,12,7,1)$&$(10,9,12,10,11,9,1)$\\
31&8&$(7,6,18,1,12,1,17)$&$(7,6,18,0,11,3,17)$\\
31&8&$(9,8,14,11,15,3,2)$&$(9,8,14,10,14,5,2)$\\
31&9&$(8,7,16,10,17,3,1)$&$(8,7,16,9,16,5,1)$\\
31&9&$(9,8,14,8,12,1,10)$&$(9,8,14,7,11,3,10)$\\
37&1&$(11,10,16,10,16,1,10)$&$(11,10,16,9,15,3,10)$\\
37&1&$(12,11,14,11,14,1,11)$&$(12,11,14,10,13,3,11)$\\
37&3&$(10,8,19,7,18,0,12)$&$(10,8,19,6,17,2,12)$\\
37&3&$(12,10,15,8,13,1,15)$&$(12,10,15,7,12,3,15)$\\
37&3&$(11,8,18,9,19,1,8)$&$(11,8,18,8,18,3,8)$\\
37&4&$(11,10,16,12,17,5,3)$&$(11,10,16,11,16,7,3)$\\
37&4&$(12,11,14,13,15,2,7)$&$(12,11,14,12,14,4,7)$\\
37&5&$(10,9,18,10,18,1,8)$&$(10,9,18,9,17,3,8)$\\
37&5&$(11,10,16,9,14,1,13)$&$(11,10,16,8,13,3,13)$\\
37&6&$(9,8,20,7,18,1,11)$&$(9,8,20,6,17,3,11)$\\
37&6&$(10,9,18,6,14,0,17)$&$(10,9,18,5,13,2,17)$\\
37&7&$(9,8,20,3,14,1,19)$&$(9,8,20,2,13,3,19)$\\
37&7&$(10,9,18,1,9,1,26)$&$(10,9,18,0,8,3,26)$\\
37&8&$(11,10,16,13,17,1,6)$&$(11,10,16,12,16,3,6)$\\
37&8&$(12,11,14,11,12,0,14)$&$(12,11,14,10,11,2,14)$\\
37&9&$(10,9,18,11,18,1,7)$&$(10,9,18,10,17,3,7)$\\
37&9&$(11,10,16,8,12,1,16)$&$(11,10,16,7,11,3,16)$\\
41&1&$(12,11,18,11,18,1,11)$&$(12,11,18,10,17,3,11)$\\
41&1&$(13,12,16,12,16,1,12)$&$(13,12,16,11,15,3,12)$\\
41&3&$(12,10,19,8,17,1,15)$&$(12,10,19,7,16,3,15)$\\
41&3&$(14,12,15,10,13,0,18)$&$(14,12,15,9,12,2,18)$\\
41&3&$(12,9,20,10,21,1,9)$&$(12,9,20,9,20,3,9)$\\
41&4&$(12,11,18,13,19,6,3)$&$(12,11,18,12,18,8,3)$\\
41&4&$(13,12,16,14,17,3,7)$&$(13,12,16,13,16,5,7)$\\
41&5&$(11,10,20,11,20,1,9)$&$(11,10,20,10,19,3,9)$\\
41&5&$(12,11,18,10,16,1,14)$&$(12,11,18,9,15,3,14)$\\
41&6&$(10,9,22,8,20,0,13)$&$(10,9,22,7,19,2,13)$\\
41&6&$(11,10,20,6,15,1,19)$&$(11,10,20,5,14,3,19)$\\
41&7&$(10,9,22,3,15,1,22)$&$(10,9,22,2,14,3,22)$\\
41&7&$(11,10,20,1,10,1,29)$&$(11,10,20,0,9,3,29)$\\
41&8&$(12,11,18,14,19,2,6)$&$(12,11,18,13,18,4,6)$\\
41&8&$(13,12,16,12,14,1,14)$&$(13,12,16,11,13,3,14)$\\
41&9&$(11,10,20,12,20,1,8)$&$(11,10,20,11,19,3,8)$\\
41&9&$(12,11,18,9,14,1,17)$&$(12,11,18,8,13,3,17)$\\
43&5&$(11,10,22,12,23,1,7)$&$(11,10,22,11,22,3,7)$\\
43&5&$(12,11,20,11,19,1,12)$&$(12,11,20,10,18,3,12)$\\
43&6&$(11,10,22,7,18,1,17)$&$(11,10,22,6,17,3,17)$\\
43&6&$(12,11,20,6,14,0,23)$&$(12,11,20,5,13,2,23)$\\
43&7&$(11,10,22,2,13,1,27)$&$(11,10,22,1,12,3,27)$\\
43&7&$(14,13,16,15,16,7,5)$&$(14,13,16,14,15,9,5)$\\
43&8&$(12,11,20,14,21,6,2)$&$(12,11,20,13,20,8,2)$\\
43&8&$(13,12,18,14,18,1,10)$&$(13,12,18,13,17,3,10)$\\
43&9&$(11,10,22,13,23,3,4)$&$(11,10,22,12,22,5,4)$\\
43&9&$(12,11,20,11,18,1,13)$&$(12,11,20,10,17,3,13)$\\
47&5&$(12,11,24,13,25,1,8)$&$(12,11,24,12,24,3,8)$\\
47&5&$(13,12,22,12,21,1,13)$&$(13,12,22,11,20,3,13)$\\
47&6&$(12,11,24,8,20,0,19)$&$(12,11,24,7,19,2,19)$\\
47&6&$(13,12,22,6,15,1,25)$&$(13,12,22,5,14,3,25)$\\
47&7&$(12,11,24,2,14,1,30)$&$(12,11,24,1,13,3,30)$\\
47&7&$(15,14,18,17,19,7,4)$&$(15,14,18,16,18,9,4)$\\
47&7&$(16,15,16,15,14,7,11)$&$(16,15,16,14,13,9,11)$\\
47&8&$(13,12,22,15,23,7,2)$&$(13,12,22,14,22,9,2)$\\
47&8&$(14,13,20,16,21,0,10)$&$(14,13,20,15,20,2,10)$\\
47&9&$(12,11,24,14,25,3,5)$&$(12,11,24,13,24,5,5)$\\
47&9&$(13,12,22,12,20,1,14)$&$(13,12,22,11,19,3,14)$\\
53&5&$(14,13,26,14,26,1,12)$&$(14,13,26,13,25,3,12)$\\
53&5&$(15,14,24,13,22,1,17)$&$(15,14,24,12,21,3,17)$\\
53&6&$(13,12,28,9,24,1,19)$&$(13,12,28,8,23,3,19)$\\
53&6&$(14,13,26,8,20,0,25)$&$(14,13,26,7,19,2,25)$\\
53&7&$(13,12,28,3,18,1,31)$&$(13,12,28,2,17,3,31)$\\
53&7&$(14,13,26,1,13,1,38)$&$(14,13,26,0,12,3,38)$\\
53&8&$(15,14,24,17,25,5,6)$&$(15,14,24,16,24,7,6)$\\
53&8&$(16,15,22,17,22,0,14)$&$(16,15,22,16,21,2,14)$\\
53&9&$(14,13,26,15,26,1,11)$&$(14,13,26,14,25,3,11)$\\
53&9&$(15,14,24,12,20,1,20)$&$(15,14,24,11,19,3,20)$\\
59&5&$(15,14,30,16,31,1,11)$&$(15,14,30,15,30,3,11)$\\
59&5&$(16,15,28,15,27,1,16)$&$(16,15,28,14,26,3,16)$\\
59&6&$(15,14,30,9,24,1,25)$&$(15,14,30,8,23,3,25)$\\
59&6&$(16,15,28,8,20,0,31)$&$(16,15,28,7,19,2,31)$\\
59&7&$(15,14,30,2,17,1,39)$&$(15,14,30,1,16,3,39)$\\
59&7&$(18,17,24,20,25,13,1)$&$(18,17,24,19,24,15,1)$\\
59&7&$(19,18,22,21,23,7,8)$&$(19,18,22,20,22,9,8)$\\
59&8&$(17,16,26,19,27,3,10)$&$(17,16,26,18,26,5,10)$\\
59&8&$(18,17,24,18,23,0,18)$&$(18,17,24,17,22,2,18)$\\
59&9&$(15,14,30,17,31,3,8)$&$(15,14,30,16,30,5,8)$\\
59&9&$(16,15,28,15,26,1,17)$&$(16,15,28,14,25,3,17)$\\
61&5&$(16,15,30,16,30,1,14)$&$(16,15,30,15,29,3,14)$\\
61&5&$(17,16,28,15,26,1,19)$&$(17,16,28,14,25,3,19)$\\
61&6&$(15,14,32,10,27,1,23)$&$(15,14,32,9,26,3,23)$\\
61&6&$(16,15,30,9,23,0,29)$&$(16,15,30,8,22,2,29)$\\
61&7&$(15,14,32,3,20,1,37)$&$(15,14,32,2,19,3,37)$\\
61&7&$(16,15,30,1,15,1,44)$&$(16,15,30,0,14,3,44)$\\
61&8&$(17,16,28,19,29,7,6)$&$(17,16,28,18,28,9,6)$\\
61&8&$(18,17,26,20,27,0,14)$&$(18,17,26,19,26,2,14)$\\
61&9&$(16,15,30,17,30,1,13)$&$(16,15,30,16,29,3,13)$\\
61&9&$(17,16,28,14,24,1,22)$&$(17,16,28,13,23,3,22)$\\
67&7&$(17,16,34,2,19,1,45)$&$(17,16,34,1,18,3,45)$\\
67&7&$(21,20,26,23,27,11,6)$&$(21,20,26,22,26,13,6)$\\
67&8&$(19,18,30,21,31,5,10)$&$(19,18,30,20,30,7,10)$\\
67&8&$(20,19,28,21,28,0,18)$&$(20,19,28,20,27,2,18)$\\
67&9&$(17,16,34,19,35,3,10)$&$(17,16,34,18,34,5,10)$\\
67&9&$(18,17,32,17,30,1,19)$&$(18,17,32,16,29,3,19)$\\
71&7&$(18,17,36,2,20,1,48)$&$(18,17,36,1,19,3,48)$\\
71&7&$(22,21,28,24,29,13,5)$&$(22,21,28,23,28,15,5)$\\
71&8&$(20,19,32,22,33,6,10)$&$(20,19,32,21,32,8,10)$\\
71&8&$(21,20,30,22,30,1,18)$&$(21,20,30,21,29,3,18)$\\
71&9&$(18,17,36,20,37,3,11)$&$(18,17,36,19,36,5,11)$\\
71&9&$(19,18,34,18,32,1,20)$&$(19,18,34,17,31,3,20)$\\
73&7&$(18,17,38,3,23,1,46)$&$(18,17,38,2,22,3,46)$\\
73&7&$(19,18,36,1,18,1,53)$&$(19,18,36,0,17,3,53)$\\
73&8&$(20,19,34,22,35,10,6)$&$(20,19,34,21,34,12,6)$\\
73&8&$(21,20,32,23,33,3,14)$&$(21,20,32,22,32,5,14)$\\
73&9&$(19,18,36,20,36,1,16)$&$(19,18,36,19,35,3,16)$\\
73&9&$(20,19,34,17,30,1,25)$&$(20,19,34,16,29,3,25)$\\
79&7&$(20,19,40,2,22,1,54)$&$(20,19,40,1,21,3,54)$\\
79&7&$(24,23,32,26,33,17,3)$&$(24,23,32,25,32,19,3)$\\
79&8&$(22,21,36,24,37,8,10)$&$(22,21,36,23,36,10,10)$\\
79&8&$(23,22,34,25,35,1,18)$&$(23,22,34,24,34,3,18)$\\
79&9&$(20,19,40,22,41,3,13)$&$(20,19,40,21,40,5,13)$\\
79&9&$(21,20,38,20,36,1,22)$&$(21,20,38,19,35,3,22)$\\
\bottomrule\end{longtable}
\endgroup

\backmatter
\begin{thebibliography}{9}
\addcontentsline{toc}{chapter}{参考文献}
\bibitem{mixing} M. Coviello Gonzalez and M. Chrobak,
\emph{Towards a Theory of Mixing Graphs: A Characterization of Perfect Mixability},
arXiv:1806.08875v4. 所用为引言的完全混合刻画定理(a)及条件MC。
\bibitem{morris} D. W. Morris,
\emph{Bounded generation of $\mathrm{SL}(n,A)$ (after D. Carter, G. Keller and E. Paige)},
New York Journal of Mathematics \textbf{13} (2007), 383--421.
所用为 Theorem 6.1(2), pp.416--417。
\bibitem{serre} J.-P. Serre,
\emph{Le probl\`eme des groupes de congruence pour $\mathrm{SL}_2$},
Annals of Mathematics (2) \textbf{92} (1970), 489--527.
\bibitem{prasad} G. Prasad and A. S. Rapinchuk,
\emph{Developments on the congruence subgroup problem after the work of Bass, Milnor and Serre},
in \emph{Collected Papers of John Milnor, Vol. V: Algebra}, AMS, 2010, 307--325.
所用强同余子群性质见 pp.308--309 的 Proposition 1 与 Section 2.3。
\end{thebibliography}


\end{document}
